% La Première Guerre mondiale et la crise économique de 1929 — Histoire 1re Tech — Cours signé OnBuch+
% Programme officiel MINESEC. Compiler avec : tectonic N03-premiere-guerre-mondiale-crise-economique-1929.tex
\newcommand{\DOCMATIERE}{Histoire}
\newcommand{\DOCNIVEAU}{1re Tech}
\newcommand{\DOCMODULE}{Histoire – classe de Première technique}
\newcommand{\DOCLECON}{N03}
\newcommand{\DOCTITRE}{La Première Guerre mondiale et la crise économique de 1929}
% OnBuch+ — préambule « cours signés » ENRICHI (compilé avec tectonic / XeLaTeX)
% Système visuel « Pop » d'OnBuch+ : fond crème (jamais blanc), bordures encre
% épaisses, ombres « dures » décalées, titres Archivo Black en capitales.
% Version MATHÉMATIQUES : graphes (pgfplots/tikz), boîtes pédagogiques
% (définition, propriété, méthode, attention, le savais-tu, exercice résolu,
% exercice, TP, bilan), page de garde et signature OnBuch+ ancrée en bas.
%
% Macros attendues dans main.tex AVANT \input{preamble} :
%   \DOCMATIERE  \DOCNIVEAU  \DOCMODULE  \DOCLECON (numéro)  \DOCTITRE
\documentclass[11pt]{article}

\usepackage{fontspec}
\usepackage[french]{babel}
\usepackage[a4paper,margin=2cm,top=3.4cm,bottom=2.8cm,headheight=1.4cm,footskip=1.3cm]{geometry}
\usepackage{amsmath,amssymb,amsfonts}
\usepackage{mathtools} % \xmapsto, \coloneqq... souvent utilisés par les modèles
\usepackage{cancel} % simplifications barrées (bilans de forces)
% \abs{z} (module, valeur absolue) et \norm{u} (norme), utilisés par les modèles
\providecommand{\abs}{}\let\abs\relax\DeclarePairedDelimiter{\abs}{\lvert}{\rvert}
\providecommand{\norm}{}\let\norm\relax\DeclarePairedDelimiter{\norm}{\lVert}{\rVert}
\usepackage[table]{xcolor} % [table] : fournit aussi \rowcolors (alternance de lignes)
\usepackage{pagecolor}
\usepackage{tikz}
\usetikzlibrary{arrows.meta,positioning,calc,shapes.geometric,shapes.misc,shapes.arrows,decorations.pathmorphing,decorations.pathreplacing,decorations.markings,patterns,shadows,mindmap,trees,fit,backgrounds,matrix,intersections,angles,quotes}
\usepackage{pgfplots}
\usepackage[version=4]{mhchem} % équations nucléaires \ce{^{235}_{92}U}
\usepackage[european,siunitx]{circuitikz} % schémas électriques
\pgfplotsset{compat=1.18}
\usepgfplotslibrary{fillbetween,groupplots} % aires entre courbes (calcul intégral)
\usepackage{enumitem}
\usepackage{fancyhdr}
% [most] charge toutes les bibliothèques tcolorbox (listings, minted, raster,
% documentation...) et gonfle inutilement la mémoire TeX sur un document long
% (~300 boîtes) : on ne charge que ce qui est réellement utilisé.
\usepackage[skins,breakable]{tcolorbox}
\usepackage{graphicx}
\usepackage{colortbl}
\usepackage{booktabs}
\usepackage{tabularx}
\usepackage{multirow}
\usepackage{diagbox} % case d'en-tête diagonale des tableaux à double entrée
% Colonnes extensibles centrée / gauche / droite (C, L, R), souvent utilisées par les modèles
\newcolumntype{C}{>{\centering\arraybackslash}X}
\newcolumntype{L}{>{\raggedright\arraybackslash}X}
\newcolumntype{R}{>{\raggedleft\arraybackslash}X}
\usepackage{array}
\usepackage{multicol}
\usepackage{siunitx}
% parse-numbers=false : accepte aussi \SI{2,5 \times 10^{-3}}{...}
\sisetup{locale=FR,output-decimal-marker={,},group-minimum-digits=4,parse-numbers=false}
\DeclareSIUnit\ohm{\ensuremath{\symup{\Omega}}} % U+2126 absent de la police
\DeclareSIUnit\division{div} % réglages d'oscilloscope (ms/div, V/div)
\DeclareSIUnit\tr{tr} % tours (vitesse de rotation, ex. \SI{1500}{\tr\per\minute})
\DeclareSIUnit\molar{M} % concentration molaire (ex. \SI{3}{\molar}, \SI{1}{\micro\molar})
\DeclareSIUnit\an{an} % année (le modèle confond parfois avec \year, un compteur TeX) — ex. \SI{5}{\kilo\meter\per\an}
\DeclareSIUnit\FCFA{FCFA} % franc CFA (ex. \SI{500}{\FCFA\per\kilo\gram})
\DeclareSIUnit\degreeBrix{\textdegree Brix} % teneur en sucre d'un sirop/jus (réfractométrie)
\DeclareSIUnit\month{mois} % le modèle utilise parfois \month, un compteur TeX, comme unité de durée
\DeclareSIUnit\microliter{\micro\liter} % le modèle écrit parfois \microliter en un seul token
\DeclareSIUnit\million{millions} % dénombrement (ex. \SI{15}{\million\per\mL} spermatozoïdes)
\usepackage{qrcode}
% \degree : commande fréquemment utilisée par les modèles pour un angle ou une
% température en dehors de \SI{}{} (non fournie par les paquets chargés).
\providecommand{\degree}{^{\circ}}
% \qed (fin de démonstration), utilisé par les modèles sans amsthm
\providecommand{\qed}{\hfill$\square$}
% \barre : double barre des valeurs interdites dans un tableau de variations (tkz-tab)
\providecommand{\barre}{\ensuremath{\|}}
% environnement proof (amsthm non chargé) utilisé par les modèles
\ifdefined\proof\else\newenvironment{proof}[1][Démonstration]{\par\noindent\textit{#1.}\ }{\qed\par}\fi
\usepackage{lastpage}
\usepackage{hyperref}
\hypersetup{hidelinks}

% ============================================================
% Palette « Pop » OnBuch+ — mêmes hex que la classe P (plus_ui.dart)
% ============================================================
\definecolor{poporange}{HTML}{F59321}
\definecolor{popdark}{HTML}{E07A0C}
\definecolor{popcream}{HTML}{FFF6E8}
\definecolor{popink}{HTML}{1C1714}
\definecolor{poppurple}{HTML}{7A5AE0}
\definecolor{popgreen}{HTML}{2E9E6A}
\definecolor{poppink}{HTML}{E0457B}
\definecolor{popblue}{HTML}{2F7DE1}
\definecolor{popgold}{HTML}{C98A12}
\definecolor{popmuted}{HTML}{8A7F76}
\definecolor{popline}{HTML}{EADFD2}
\definecolor{popgray}{HTML}{8A7F76}
\colorlet{popgrey}{popgray}
% Alias tolérants (le modèle nomme parfois les couleurs différemment)
\colorlet{popwhite}{white}
\colorlet{popblack}{popink}
\colorlet{popred}{poppink}
\colorlet{popyellow}{popgold}
% Teintes claires (fonds de tableaux/encadrés) — le modèle nomme parfois
% les couleurs avec un suffixe L (ex. popblueL) sans les avoir définies.
\colorlet{poporangeL}{poporange!20}
\colorlet{popdarkL}{popdark!20}
\colorlet{poppurpleL}{poppurple!20}
\colorlet{popgreenL}{popgreen!20}
\colorlet{poppinkL}{poppink!20}
\colorlet{popblueL}{popblue!20}
\colorlet{popgoldL}{popgold!20}
\colorlet{popredL}{popred!20}
\colorlet{popgrayL}{popgray!20}
% Teintes claires pour fonds de tableaux / zones
\colorlet{popblueL}{popblue!10!white}
\colorlet{poporangeL}{poporange!14!white}
\colorlet{popgreenL}{popgreen!12!white}
\colorlet{poppurpleL}{poppurple!12!white}
\colorlet{poppinkL}{poppink!12!white}
\colorlet{popgoldL}{popgold!14!white}

\pagecolor{popcream}

% ============================================================
% Typographie
% ============================================================
\defaultfontfeatures{Ligatures=TeX}
\setmainfont{PlusJakartaSans-Medium.ttf}[Path=../../../fonts/,BoldFont=PlusJakartaSans-Bold.ttf]
% Maths en sans-serif (Fira Math) avec chiffres et lettres droites de Plus
% Jakarta, pour que les formules se fondent dans le texte.
\usepackage[math-style=ISO,bold-style=ISO]{unicode-math}
\setmathfont{FiraMath-Regular.otf}[Path=../../../fonts/]
\setmathfont{PlusJakartaSans-Medium.ttf}[Path=../../../fonts/,range={up/num,up/{Latin,latin},\mathup}]
\usepackage{icomma} % virgule décimale sans espace en mode math
% Caractères mathématiques Unicode absents des polices de texte (Plus Jakarta,
% Archivo Black) : rendus par leur équivalent mathématique, en texte comme en
% formule — sinon ils s'affichent en carré vide (ex. « Arithmétique dans ℤ »).
\usepackage{newunicodechar}
\newunicodechar{ℝ}{\ensuremath{\mathbb{R}}}
\newunicodechar{ℤ}{\ensuremath{\mathbb{Z}}}
\newunicodechar{ℂ}{\ensuremath{\mathbb{C}}}
\newunicodechar{ℕ}{\ensuremath{\mathbb{N}}}
\newunicodechar{ℚ}{\ensuremath{\mathbb{Q}}}
\newunicodechar{𝔻}{\ensuremath{\mathbb{D}}}
\newunicodechar{α}{\ensuremath{\alpha}}
\newunicodechar{β}{\ensuremath{\beta}}
\newunicodechar{γ}{\ensuremath{\gamma}}
\newunicodechar{δ}{\ensuremath{\delta}}
\newunicodechar{ε}{\ensuremath{\varepsilon}}
\newunicodechar{θ}{\ensuremath{\theta}}
\newunicodechar{λ}{\ensuremath{\lambda}}
\newunicodechar{μ}{\ensuremath{\mu}}
\newunicodechar{σ}{\ensuremath{\sigma}}
\newunicodechar{τ}{\ensuremath{\tau}}
\newunicodechar{φ}{\ensuremath{\varphi}}
\newunicodechar{ψ}{\ensuremath{\psi}}
\newunicodechar{ω}{\ensuremath{\omega}}
\newunicodechar{Σ}{\ensuremath{\Sigma}}
% majuscules produites par \MakeUppercase dans les titres (α-aminés -> Α)
\newunicodechar{Α}{\ensuremath{\alpha}}
\newunicodechar{Β}{\ensuremath{\beta}}
\newunicodechar{Γ}{\ensuremath{\Gamma}}
\newunicodechar{Δ}{\ensuremath{\Delta}}
\newunicodechar{Ε}{\ensuremath{\varepsilon}}
\newunicodechar{Θ}{\ensuremath{\Theta}}
\newunicodechar{Λ}{\ensuremath{\Lambda}}
\newunicodechar{Μ}{\ensuremath{\mu}}
\newunicodechar{Τ}{\ensuremath{\tau}}
\newunicodechar{Φ}{\ensuremath{\Phi}}
\newunicodechar{Ψ}{\ensuremath{\Psi}}
\newunicodechar{Ω}{\ensuremath{\Omega}}
\newunicodechar{↦}{\ensuremath{\mapsto}}
\newunicodechar{∈}{\ensuremath{\in}}
\newunicodechar{∉}{\ensuremath{\notin}}
\newunicodechar{∧}{\ensuremath{\wedge}}
\newunicodechar{⇔}{\ensuremath{\Leftrightarrow}}
\newunicodechar{⟺}{\ensuremath{\Longleftrightarrow}}
\newunicodechar{⇒}{\ensuremath{\Rightarrow}}
\newunicodechar{⟹}{\ensuremath{\Longrightarrow}}
\newunicodechar{∘}{\ensuremath{\circ}}
\newunicodechar{∪}{\ensuremath{\cup}}
\newunicodechar{∩}{\ensuremath{\cap}}
\newunicodechar{≡}{\ensuremath{\equiv}}
\newunicodechar{⊂}{\ensuremath{\subset}}
\newunicodechar{⊥}{\ensuremath{\perp}}
\newunicodechar{∃}{\ensuremath{\exists}}
\newunicodechar{∀}{\ensuremath{\forall}}
\newunicodechar{∛}{\ensuremath{\sqrt[3]{\,}}}
\newunicodechar{∜}{\ensuremath{\sqrt[4]{\,}}}
\newunicodechar{⃗}{\ensuremath{{}^{\to}}}
\newunicodechar{ⁿ}{\ifmmode^{n}\else\textsuperscript{\ensuremath{n}}\fi}
\newunicodechar{ˣ}{\ifmmode^{x}\else\textsuperscript{\ensuremath{x}}\fi}
\newunicodechar{ᵇ}{\ifmmode^{b}\else\textsuperscript{\ensuremath{b}}\fi}
\newunicodechar{ʳ}{\ifmmode^{r}\else\textsuperscript{\ensuremath{r}}\fi}
\newunicodechar{ᵘ}{\ifmmode^{u}\else\textsuperscript{\ensuremath{u}}\fi}
\newunicodechar{ʸ}{\ifmmode^{y}\else\textsuperscript{\ensuremath{y}}\fi}
\newunicodechar{ᵃ}{\ifmmode^{a}\else\textsuperscript{\ensuremath{a}}\fi}
\newunicodechar{ᵉ}{\ifmmode^{e}\else\textsuperscript{\ensuremath{e}}\fi}
\newunicodechar{ᵏ}{\ifmmode^{k}\else\textsuperscript{\ensuremath{k}}\fi}
\newunicodechar{ᵖ}{\ifmmode^{p}\else\textsuperscript{\ensuremath{p}}\fi}
\newunicodechar{ᴺ}{\ifmmode^{N}\else\textsuperscript{\ensuremath{N}}\fi}
\newunicodechar{ᶜ}{\ifmmode^{c}\else\textsuperscript{\ensuremath{c}}\fi}
\newunicodechar{ʰ}{\ifmmode^{h}\else\textsuperscript{\ensuremath{h}}\fi}
\newunicodechar{ᵠ}{\ifmmode^{q}\else\textsuperscript{\ensuremath{q}}\fi}
\newunicodechar{ₙ}{\ifmmode_{n}\else\textsubscript{\ensuremath{n}}\fi}
\newunicodechar{ₐ}{\ifmmode_{a}\else\textsubscript{\ensuremath{a}}\fi}
\newunicodechar{ᵢ}{\ifmmode_{i}\else\textsubscript{\ensuremath{i}}\fi}
\newunicodechar{ᵣ}{\ifmmode_{r}\else\textsubscript{\ensuremath{r}}\fi}
\newunicodechar{ₖ}{\ifmmode_{k}\else\textsubscript{\ensuremath{k}}\fi}
\newunicodechar{ᵦ}{\ifmmode_{\beta}\else\textsubscript{\ensuremath{\beta}}\fi}
\newunicodechar{ₓ}{\ifmmode_{x}\else\textsubscript{\ensuremath{x}}\fi}
\newfontfamily\popdisplay{ArchivoBlack-Regular.ttf}[Path=../../../fonts/]
\newfontfamily\poplogo{SpaceGrotesk-Bold.ttf}[Path=../../../fonts/]
\newfontfamily\popsemi{PlusJakartaSans-SemiBold.ttf}[Path=../../../fonts/]
\newfontfamily\popxbold{PlusJakartaSans-ExtraBold.ttf}[Path=../../../fonts/]
\newfontfamily\popxboldls{PlusJakartaSans-ExtraBold.ttf}[Path=../../../fonts/,LetterSpace=3.0]

\setlist[enumerate]{leftmargin=1.7em,itemsep=5pt,topsep=4pt}
\setlist[itemize]{leftmargin=1.7em,itemsep=4pt,topsep=4pt,label={\color{poporange}\textbullet}}
\setlength{\parindent}{0pt}
\setlength{\parskip}{5pt}
\renewcommand{\arraystretch}{1.25}

% pgfplots : style Pop par défaut
\pgfplotsset{
  every axis/.append style={axis line style={popink,line width=1pt},tick style={popink},
    grid style={popline,line width=0.5pt},label style={font=\small},tick label style={font=\footnotesize}},
  popcurve/.style={poporange,line width=1.8pt,no markers,smooth},
}
% Même style disponible dans un \draw TikZ simple (hors pgfplots)
\tikzset{popcurve/.style={poporange,line width=1.8pt}}
\tikzset{popgrid/.style={popline,very thin}}
\colorlet{popcurve}{poporange} % le nom du style est parfois utilisé comme couleur

% ============================================================
% Pill / squiggle
% ============================================================
\newcommand{\poppill}[3]{%
  \tikz[baseline=-0.55ex]\node[draw=popink,line width=1.2pt,rounded corners=2.6pt,%
    fill=#2,inner xsep=6.5pt,inner ysep=3pt]{%
    \popxboldls\fontsize{7}{7}\selectfont\color{#3}\MakeUppercase{#1}};%
}
\newcommand{\popsquiggle}[1]{%
  \tikz[baseline=-0.4ex]{
    \draw[#1,line width=2.6pt,line cap=round]
      plot [smooth] coordinates {(0,0.09) (0.34,0.01) (0.68,0.21) (1.02,0.01) (1.36,0.10)};
  }%
}
% Mot-clé surligné (marqueur orange sous le texte)
\newcommand{\cle}[1]{{\popxbold\color{popdark}#1}}

% ============================================================
% En-tête / pied
% ============================================================
\pagestyle{fancy}
\fancyhf{}
\renewcommand{\headrulewidth}{0pt}
\renewcommand{\footrulewidth}{0pt}
\lhead{\raisebox{-0.15ex}{\poplogo\fontsize{13}{13}\selectfont\color{popink}OnBuch\color{poporange}+}}
\rhead{\poppill{\DOCMATIERE\ \textbullet\ \DOCNIVEAU\ \textbullet\ Leçon \DOCLECON}{white}{popink}}
\lfoot{\popxbold\fontsize{7.6}{7.6}\selectfont\color{popmuted}\MakeUppercase{Cours signé OnBuch+}\ \textbullet\ {\popsemi\DOCTITRE}}
\rfoot{\tikz[baseline=-0.6ex]\node[rounded corners=8.5pt,fill=popink,minimum height=17pt,minimum width=17pt,inner xsep=5pt,inner sep=0]{\popxbold\fontsize{8}{8}\selectfont\color{popcream}\thepage\,/\,\pageref*{LastPage}};}

% ============================================================
% Titres
% ============================================================
\newcommand{\chaptitre}[2]{%
  \begin{tcolorbox}[enhanced,colback=poporange,colframe=popink,boxrule=2.6pt,arc=11pt,
      drop shadow=popink,left=15pt,right=15pt,top=12pt,bottom=11pt,before skip=3pt,after skip=18pt]
    \poppill{#2}{popink}{popcream}\par\vspace{9pt}
    {\raggedright\hyphenpenalty=10000\exhyphenpenalty=10000\popdisplay\fontsize{22}{24}\selectfont\color{popcream}\MakeUppercase{#1}\par}
  \end{tcolorbox}
}
% Section numérotée : pastille orange avec le numéro + titre Archivo
\newcounter{popsec}
\newcommand{\coursec}[1]{%
  \stepcounter{popsec}\setcounter{popsub}{0}%
  \par\vspace{16pt}\noindent
  \begin{minipage}{\linewidth}
  \tikz[baseline=-0.7ex]\node[draw=popink,line width=1.6pt,fill=poporange,rounded corners=4pt,minimum size=22pt,inner sep=0]{\popdisplay\fontsize{12}{12}\selectfont\color{popcream}\thepopsec};%
  \hspace{8pt}{\popdisplay\fontsize{15}{17}\selectfont\color{popink}\MakeUppercase{#1}}\par
  \vspace{4pt}\noindent{\color{poporange}\rule{36pt}{3.6pt}}
  \end{minipage}\par\nopagebreak\vspace{8pt}
}
\newcounter{popsub}
\newcommand{\courssub}[1]{%
  \stepcounter{popsub}%
  \par\vspace{9pt}\noindent{\popxbold\fontsize{11.5}{13}\selectfont\color{popink}{\color{poporange}\thepopsec.\thepopsub}\hspace{6pt}#1}\par\nopagebreak\vspace{4pt}
}

% ============================================================
% Boîtes pédagogiques — même recette : fond blanc, bordure épaisse colorée,
% ombre dure encre, badge plein. Titre optionnel [#1].
% ============================================================
\tcbset{popbase/.style={breakable,enhanced,colback=white,boxrule=2.3pt,arc=9pt,
  shadow={3pt}{-3pt}{0pt}{popink},left=14pt,right=14pt,top=11pt,bottom=11pt,before skip=11pt,after skip=15pt,
  fontupper=\popsemi\fontsize{10.6}{14.5}\selectfont\color{popink},
  fontlower=\popsemi\fontsize{10.6}{14.5}\selectfont\color{popink}}}
% Coche dessinée (le glyphe ✓ n'existe pas dans les polices du document)
\newcommand{\popcheck}[1]{\tikz[baseline=-0.5ex]\draw[#1,line width=1.6pt,line cap=round,line join=round](-0.13,0.0)--(-0.04,-0.09)--(0.13,0.1);}
\providecommand{\coche}{\popcheck{popgreen}}
\providecommand{\resultat}[1]{\par\smallskip\noindent\fcolorbox{popgreen}{popgreenL}{\parbox{\dimexpr\linewidth-2\fboxsep-2\fboxrule\relax}{\textbf{Résultat.} #1}}\par\smallskip}
\newcommand{\popboxhead}[3]{\poppill{#1}{#2}{white}\ifx\relax#3\relax\else\hspace{6pt}{\popxbold\fontsize{10.6}{12}\selectfont\color{popink}#3}\fi\par\vspace{7pt}}

\newtcolorbox{aretenir}[1][]{popbase,colframe=popgreen,before upper={\popboxhead{À retenir}{popgreen}{#1}}}
\newtcolorbox{exemplebox}[1][]{popbase,colframe=poporange,before upper={\popboxhead{Exemple}{poporange}{#1}}}
\newtcolorbox{definition}[1][]{popbase,colframe=popblue,before upper={\popboxhead{Définition}{popblue}{#1}}}
\newtcolorbox{propriete}[1][]{popbase,colframe=poppurple,before upper={\popboxhead{Propriété}{poppurple}{#1}}}
\newtcolorbox{methode}[1][]{popbase,colframe=popgold,before upper={\popboxhead{Méthode}{popgold}{#1}}}
\newtcolorbox{attention}[1][]{popbase,colframe=poppink,before upper={\popboxhead{Attention}{poppink}{#1}}}
\newtcolorbox{experience}[1][]{popbase,colframe=popblue,colback=popblueL,before upper={\popboxhead{Expérience}{popblue}{#1}}}
% « Le savais-tu ? » : carte violette pleine (contexte, culture, Cameroun)
\newtcolorbox{savaistu}[1][]{popbase,colback=poppurple,colframe=popink,
  fontupper=\popsemi\fontsize{10.4}{14.2}\selectfont\color{white},
  before upper={\poppill{Le savais-tu\,?}{white}{poppurple}\ifx\relax#1\relax\else\hspace{6pt}{\popxbold\color{white}#1}\fi\par\vspace{7pt}}}
% Exercice résolu : énoncé en haut, \tcblower puis la solution sur fond crème
\newcounter{popexo}\newcounter{popexr}
\newtcolorbox{exoresolu}[1][]{popbase,colframe=poporange,colbacklower=popcream,
  segmentation style={popink,line width=1.2pt,dashed},
  before upper={\stepcounter{popexr}\popboxhead{Exercice résolu \thepopexr}{poporange}{#1}},
  before lower={\poppill{Solution}{popink}{popcream}\par\vspace{6pt}}}
% Exercice (non résolu en place) — la correction est dans la partie Corrigés
\newtcolorbox{exercice}[1][]{popbase,colframe=popink,
  before upper={\stepcounter{popexo}\popboxhead{Exercice \thepopexo}{popink}{#1}}}
% Corrigé
\newtcolorbox{corrige}[1][]{popbase,colframe=popgreen,colback=popgreenL,
  before upper={\popboxhead{Corrigé}{popgreen}{#1}}}
% Objectifs / prérequis / bilan
\newtcolorbox{objectifs}{popbase,colframe=popink,colback=poporangeL,
  before upper={\popboxhead{Objectifs de la leçon}{popink}{}}}
\newtcolorbox{prerequis}{popbase,colframe=popink,colback=popblueL,
  before upper={\popboxhead{Prérequis}{popblue}{}}}
\newtcolorbox{bilan}{popbase,colframe=popink,colback=white,boxrule=2.8pt,
  before upper={\popboxhead{Fiche bilan}{poporange}{L'essentiel à connaître pour l'examen}}}
% Figure encadrée avec légende
\newtcolorbox{popfigure}[1][]{enhanced,breakable=false,colback=white,colframe=popline,boxrule=1.4pt,arc=8pt,
  left=10pt,right=10pt,top=10pt,bottom=8pt,before skip=10pt,after skip=12pt,halign=center,
  lower separated=false,halign lower=center,
  fontlower=\popsemi\fontsize{9}{11.5}\selectfont\color{popmuted},
  #1}
\newcommand{\legende}[1]{\par\vspace{6pt}{\popsemi\fontsize{9}{11.5}\selectfont\color{popmuted}#1}}

% ============================================================
% Signature OnBuch+
% ============================================================
\newcommand{\poplogoinline}[1]{{\poplogo\fontsize{#1}{#1}\selectfont\color{popink}OnBuch\color{poporange}+}}
\newcommand{\signatureonbuch}{%
  \begin{tcolorbox}[enhanced,colback=popink,colframe=popink,boxrule=2.2pt,arc=10pt,
      drop shadow=poporange,left=14pt,right=14pt,top=10pt,bottom=10pt,before skip=0pt,after skip=0pt]
    \tikz[baseline=-0.7ex]\node[circle,fill=popgreen,draw=popcream,line width=1.3pt,minimum size=22pt,inner sep=0]{\popcheck{white}};%
    \hspace{9pt}\begin{minipage}[c]{0.62\linewidth}
      {\popxbold\fontsize{12}{13}\selectfont\color{popcream}Signé\ \textbullet\ {\poplogo OnBuch\color{poporange}+}}\par\vspace{2pt}
      {\raggedright\popsemi\fontsize{8}{10.5}\selectfont\color{popline}\MakeUppercase{Équipe pédagogique OnBuch+}\\\MakeUppercase{Conforme au programme officiel MINESEC}\par}
    \end{minipage}\hfill
    {\poplogo\fontsize{22}{22}\selectfont\color{popcream}OnBuch\color{poporange}+}
  \end{tcolorbox}%
}

% ============================================================
% Page de garde
% #1 titre · #2 matière · #3 niveau · #4 module · #5 n° leçon · #6 durée ·
% #7 description / objectifs (texte ou itemize)
% ============================================================
\newcommand{\pagegarde}[7]{%
  \thispagestyle{empty}
  \newgeometry{a4paper,margin=1.7cm,top=1.6cm,bottom=1.4cm}%
  \begin{tikzpicture}[remember picture,overlay]
    \fill[poporange!18] ([xshift=-3.2cm,yshift=-3.6cm]current page.north east) circle (4.6cm);
    \fill[poppurple!14] ([xshift=2.4cm,yshift=7.6cm]current page.south west) circle (3.8cm);
  \end{tikzpicture}%
  {\poplogo\fontsize{30}{30}\selectfont\color{popink}OnBuch\color{poporange}+}\hfill
  \raisebox{0.6ex}{\poppill{Cours signé \textbullet\ Premium}{popink}{popcream}}\par
  \vspace{1pt}\popsquiggle{poppurple}\par
  \vspace{8pt}
  \begin{tcolorbox}[enhanced,colback=poporange,colframe=popink,boxrule=2.8pt,arc=13pt,
      drop shadow=popink,left=18pt,right=18pt,top=12pt,bottom=13pt,after skip=10pt]
    \poppill{#2 \textbullet\ #3}{popink}{popcream}\hspace{5pt}\poppill{#4}{white}{popink}\par\vspace{8pt}
    {\popdisplay\fontsize{14}{14}\selectfont\color{popink}LEÇON #5}\par\vspace{4pt}
    {\raggedright\hyphenpenalty=10000\exhyphenpenalty=10000\popdisplay\fontsize{25}{27}\selectfont\color{popcream}\MakeUppercase{#1}\par}
    \vspace{8pt}
    {\popxbold\fontsize{9.5}{9.5}\selectfont\color{popink}\MakeUppercase{Durée conseillée : #6}}
  \end{tcolorbox}
  \begin{tcolorbox}[enhanced,colback=white,colframe=popink,boxrule=2.2pt,arc=10pt,
      drop shadow=popink,left=16pt,right=16pt,top=10pt,bottom=10pt,after skip=8pt]
    \poppill{Dans cette leçon}{poppurple}{white}\par\vspace{5pt}
    \popsemi\fontsize{9.3}{12.6}\selectfont\color{popink}#7
  \end{tcolorbox}
  \noindent\begin{minipage}[t]{0.487\linewidth}
    \begin{tcolorbox}[enhanced,colback=popgreen,colframe=popink,boxrule=2.2pt,arc=10pt,
        drop shadow=popink,left=11pt,right=11pt,top=10pt,bottom=10pt,height=3.1cm,valign=center]
      \tcbox[colback=white,colframe=popink,boxrule=1.3pt,arc=5pt,boxsep=0pt,nobeforeafter,
        left=3pt,right=3pt,top=3pt,bottom=3pt,valign=center]{\qrcode[height=1.9cm]{https://play.google.com/store/apps/details?id=com.luvvix.onbuch&hl=fr}}%
      \hspace{8pt}\begin{minipage}[c]{0.5\linewidth}
        {\popdisplay\fontsize{11}{12.5}\selectfont\color{white}\MakeUppercase{Télécharge OnBuch}}\par\vspace{4pt}
        {\raggedright\popsemi\fontsize{8.4}{11}\selectfont\color{white}Gratuit \textbullet\ Google Play\\Tuteur IA, annales, concours, résultats\par}
      \end{minipage}
    \end{tcolorbox}
  \end{minipage}\hfill
  \begin{minipage}[t]{0.487\linewidth}
    \begin{tcolorbox}[enhanced,colback=poppurple,colframe=popink,boxrule=2.2pt,arc=10pt,
        drop shadow=popink,left=11pt,right=11pt,top=10pt,bottom=10pt,height=3.1cm,valign=center]
      \tikz[baseline=-0.7ex]\node[circle,fill=white,draw=popink,line width=1.3pt,minimum size=1.6cm,inner sep=0]{\popdisplay\fontsize{24}{24}\selectfont\color{poppurple}+};%
      \hspace{8pt}\begin{minipage}[c]{0.6\linewidth}
        {\popdisplay\fontsize{11}{12.5}\selectfont\color{white}\MakeUppercase{Passe à OnBuch+}}\par\vspace{4pt}
        {\raggedright\popsemi\fontsize{8.4}{11}\selectfont\color{white}Tous les cours signés, Tuteur IA illimité, fiches PDF — onglet OnBuch+ de l'app\par}
      \end{minipage}
    \end{tcolorbox}
  \end{minipage}
  \vspace{8pt}
  \popcontenu
  \vfill
  \signatureonbuch
  \restoregeometry
  \newpage
}

% Rangée « Ce cours contient » (page de garde)
\newcommand{\poptile}[3]{% #1 couleur #2 gros texte #3 légende
  \begin{tcolorbox}[enhanced,colback=#1,colframe=popink,boxrule=2pt,arc=9pt,shadow={2.5pt}{-2.5pt}{0pt}{popink},
      left=6pt,right=6pt,top=9pt,bottom=9pt,halign=center,valign=center,height=1.75cm,width=\linewidth,nobeforeafter]
    {\popdisplay\fontsize{12.5}{13}\selectfont\color{white}\MakeUppercase{#2}}\par\vspace{3pt}
    {\centering\popsemi\fontsize{7.8}{9.5}\selectfont\color{white}#3\par}
  \end{tcolorbox}}
\newcommand{\popcontenu}{%
  \par\noindent{\popxboldls\fontsize{7.5}{7.5}\selectfont\color{popmuted}CE COURS SIGNÉ CONTIENT}\par\vspace{7pt}
  \noindent\begin{minipage}[t]{0.235\linewidth}\poptile{popblue}{Cours}{complet, illustré, pas à pas}\end{minipage}\hfill
  \begin{minipage}[t]{0.235\linewidth}\poptile{poporange}{Méthodes}{et exercices résolus}\end{minipage}\hfill
  \begin{minipage}[t]{0.235\linewidth}\poptile{poppink}{Exercices}{type examen + corrigés}\end{minipage}\hfill
  \begin{minipage}[t]{0.235\linewidth}\poptile{popgreen}{Fiche bilan}{l'essentiel à retenir}\end{minipage}\par}

% Sommaire
\newtcolorbox{sommaire}{enhanced,colback=white,colframe=popink,boxrule=2.2pt,arc=10pt,
  drop shadow=popink,left=18pt,right=18pt,top=13pt,bottom=13pt,before skip=6pt,after skip=20pt,
  before upper={\poppill{Sommaire}{poppurple}{white}\par\vspace{9pt}\popsemi\fontsize{10.6}{14.5}\selectfont\color{popink}}}

% Signature de fin de cours — poussée tout en bas de la dernière page
\newcommand{\findecours}{%
  \par\vfill\nopagebreak
  \begin{center}\popsquiggle{poporange}\end{center}\vspace{4pt}
  \signatureonbuch
}
\AtBeginDocument{\def\llbracket{\mathopen{[\mkern-3.5mu[}}\def\rrbracket{\mathclose{]\mkern-3.5mu]}}} % intervalles d'entiers (glyphe ⟦ absent de la police)
% Glyphes absents de Fira Math : redéfinis à partir de glyphes disponibles
\AtBeginDocument{%
  \def\setminus{\mathbin{\backslash}}\let\smallsetminus\setminus
  \def\vdots{\mathord{\vbox{\baselineskip3.2pt\lineskiplimit0pt\kern4pt\hbox{.}\hbox{.}\hbox{.}}}}%
  \def\ddots{\mathinner{\mkern1mu\raise6.5pt\hbox{.}\mkern2mu\raise3.6pt\hbox{.}\mkern2mu\raise0.7pt\hbox{.}\mkern1mu}}%
  \def\triangle{\mathord{\scalebox{0.9}{$\symup{\Delta}$}}}%
  \def\nabla{\mathord{\raisebox{\depth}{\scalebox{1}[-1]{$\symup{\Delta}$}}}}%
  \def\checkmark{\popcheck{popdark}}%
  \def\star{\mathbin{\ast}}\def\bigstar{\mathord{\ast}}%
  \def\diamond{\mathbin{\scalebox{0.6}{$\Diamond$}}}%
  \def\bigcup{\mathop{\mathchoice{\raisebox{-0.3em}{\scalebox{1.7}{$\cup$}}}{\scalebox{1.25}{$\cup$}}{\cup}{\cup}}\displaylimits}%
  \def\bigcap{\mathop{\mathchoice{\raisebox{-0.3em}{\scalebox{1.7}{$\cap$}}}{\scalebox{1.25}{$\cap$}}{\cap}{\cap}}\displaylimits}%
  \def\lesseqqgtr{\mathrel{\lessgtr}}\def\bigoplus{\mathop{\oplus}\displaylimits}%
}
\providecommand{\ssi}{si et seulement si} % abréviation fréquente
\AtBeginDocument{\providecommand{\wideparen}{\overparen}} % arc de cercle

%% FIGFIX-BEGIN v5 — correctifs globaux des figures (docs/onbuch-generation/tools/figfix)
\usepackage{adjustbox}\usepackage{etoolbox}\tcbuselibrary{raster}
% toute figure TikZ/pgfplots plus large que la ligne est réduite à la largeur disponible
\BeforeBeginEnvironment{tikzpicture}{\begin{adjustbox}{max width=\linewidth}}
\AfterEndEnvironment{tikzpicture}{\end{adjustbox}}
% légendes pgfplots lisibles par-dessus les courbes
\pgfplotsset{every axis/.append style={legend style={fill=white,fill opacity=0.92,text opacity=1,draw=black!15,inner sep=3pt}}}
% étiquettes d'arêtes (syntaxe edge["..."]) avec fond blanc
\tikzset{every edge quotes/.append style={fill=white,inner sep=1.2pt}}
%% FIGFIX-END
\begin{document}

\pagegarde{La Première Guerre mondiale et la crise économique de 1929}{Histoire}{1re Tech}{Histoire – classe de Première technique}{N03}{2 séances (fiche de progression harmonisée)}{Cette leçon étudie les causes, le déroulement et les conséquences de la Première Guerre mondiale (1914-1918), puis les mécanismes et les effets de la crise économique de 1929. Elle montre, documents à l'appui, comment ces deux événements mondiaux ont directement touché le Cameroun et transformé l'ordre international.
\vspace{4pt}
\begin{multicols}{2}\small
\begin{itemize}
\item À la fin de cette leçon, je sais identifier les causes profondes et le déclencheur de la Première Guerre mondiale
\item À la fin de cette leçon, je sais décrire les grandes phases du conflit (guerre de mouvement, guerre de positions, guerre totale)
\item À la fin de cette leçon, je sais localiser sur une carte les fronts, les belligérants et les campagnes du Cameroun (1914-1916)
\item À la fin de cette leçon, je sais analyser les conséquences humaines, territoriales et économiques de la guerre, y compris pour le Cameroun
\item À la fin de cette leçon, je sais expliquer les mécanismes du krach boursier d'octobre 1929 et sa propagation à l'économie mondiale
\item À la fin de cette leçon, je sais interpréter des données chiffrées et des courbes (production, chômage, cours boursiers) pour caractériser la crise
\end{itemize}\end{multicols}}

\begin{sommaire}
\begin{multicols}{2}
\begin{enumerate}
\item Les causes et le déclenchement de la Première Guerre mondiale
\item Le déroulement du conflit : une guerre totale (1914-1918)
\item Les conséquences de la guerre et le nouvel ordre mondial
\item La prospérité américaine des années 1920 et ses fragilités
\item Le krach de 1929 et la grande crise économique mondiale
\item Travaux pratiques : cartes, croquis et analyse de données
\item Activité d'intégration
\item Méthodes et exercices résolus
\item Exercices
\item Corrigés des exercices
\item Fiche bilan
\end{enumerate}
\end{multicols}
\end{sommaire}

\begin{objectifs}
\begin{itemize}
\item À la fin de cette leçon, je sais identifier les causes profondes et le déclencheur de la Première Guerre mondiale
\item À la fin de cette leçon, je sais décrire les grandes phases du conflit (guerre de mouvement, guerre de positions, guerre totale)
\item À la fin de cette leçon, je sais localiser sur une carte les fronts, les belligérants et les campagnes du Cameroun (1914-1916)
\item À la fin de cette leçon, je sais analyser les conséquences humaines, territoriales et économiques de la guerre, y compris pour le Cameroun
\item À la fin de cette leçon, je sais expliquer les mécanismes du krach boursier d'octobre 1929 et sa propagation à l'économie mondiale
\item À la fin de cette leçon, je sais interpréter des données chiffrées et des courbes (production, chômage, cours boursiers) pour caractériser la crise
\item À la fin de cette leçon, je sais rédiger et justifier une conclusion argumentée à partir de documents historiques
\item À la fin de cette leçon, je sais relier ces événements mondiaux à des situations concrètes de la vie courante au Cameroun (prix du cacao, chômage, mémoire des anciens combattants)
\end{itemize}
\end{objectifs}

\begin{prerequis}
\begin{itemize}
\item Les grandes puissances européennes et la colonisation de l'Afrique à la fin du XIXe siècle (classe de 4e)
\item La partage du Cameroun : protectorat allemand (1884), puis mandats français et britannique
\item Notions économiques de base : production, échange, monnaie, Bourse
\item Lecture de cartes historiques et de graphiques statistiques simples
\item La notion de crise économique (inflation, chômage) abordée en début d'année
\end{itemize}
\end{prerequis}

\newpage

\coursec{Les causes et le déclenchement de la Première Guerre mondiale}

En 1916, des milliers de jeunes Camerounais sont enrôlés de force comme tirailleurs et porteurs dans une guerre qui oppose la France, la Grande-Bretagne et l'Allemagne sur leur propre sol : Douala et Yaoundé deviennent des champs de bataille. Quinze ans plus tard, en 1931, le prix du cacao s'effondre sur le marché de Douala et les planteurs camerounais ne vendent plus leurs récoltes, ruinés par une crise née à des milliers de kilomètres, à la Bourse de New York. Comment une guerre européenne et un krach américain ont-ils bouleversé la vie quotidienne des Camerounais ? Cette leçon explique les mécanismes de ces deux catastrophes mondiales et leurs conséquences concrètes sur le Cameroun colonial.

\textbf{Problématique :} comment l'Europe, pourtant prospère et organisée en alliances censées garantir la paix, bascule-t-elle en quelques semaines, à l'été 1914, dans la plus grande guerre que le monde ait connue jusqu'alors ?

\courssub{Les rivalités impérialistes entre les puissances européennes}

Pour comprendre la Première Guerre mondiale, il faut d'abord comprendre pourquoi les grandes puissances européennes se regardaient en chiens de fusil bien avant 1914. L'idée intuitive est simple : imagine plusieurs commerçants qui se disputent le même marché, le même quartier, les mêmes clients. Chacun veut vendre plus que l'autre, et chacun finit par se sentir menacé. À l'échelle des États, cette concurrence porte un nom : l'\cle{impérialisme}.

\begin{definition}[L'impérialisme]
L'\cle{impérialisme} est la politique d'expansion territoriale et économique d'un État au détriment d'autres peuples. Un État impérialiste cherche à conquérir des colonies, à contrôler des marchés, des matières premières et des routes commerciales, afin d'accroître sa puissance et sa richesse.
\end{definition}

À la fin du XIX\textsuperscript{e} siècle et au début du XX\textsuperscript{e} siècle, les puissances européennes se partagent le monde, notamment l'Afrique (conférence de Berlin, 1884-1885) et l'Asie. Mais ce partage crée des frustrations et des conflits permanents :

\begin{itemize}
\item \textbf{La France contre l'Allemagne.} La France, vaincue en 1871, a perdu l'Alsace-Lorraine au profit de l'Allemagne. Elle nourrit un désir de \emph{revanche}. De plus, les deux pays se disputent le Maroc : les crises marocaines de 1905 (coup de Tanger) et 1911 (coup d'Agadir) les amènent au bord de la guerre. L'Allemagne, arrivée tardivement au partage colonial, estime ne pas avoir reçu sa « part du gâteau » : elle possède le Togo, le Kamerun, le Sud-Ouest africain et l'Afrique orientale allemande, mais juge cet empire trop modeste face à ceux de la France et de la Grande-Bretagne.
\item \textbf{La Grande-Bretagne contre l'Allemagne.} La Grande-Bretagne possède le plus vaste empire colonial du monde (Inde, Canada, Australie, une grande partie de l'Afrique). Elle domine les mers grâce à sa flotte, la \emph{Royal Navy}. Or l'Allemagne de l'empereur Guillaume II construit à partir de 1898 une flotte de guerre moderne, la \emph{Kaiserliche Marine} (marine impériale), ce que Londres vit comme une menace directe pour sa suprématie navale et son commerce.
\item \textbf{La Russie contre l'Autriche-Hongrie.} Ces deux empires se disputent l'influence dans les Balkans, région du sud-est de l'Europe où l'Empire ottoman recule.
\end{itemize}

Ces rivalités économiques et coloniales constituent des \textbf{causes profondes} de la guerre : elles créent un climat de méfiance générale où chaque puissance se prépare au conflit.

\courssub{La course aux armements et la formation des deux alliances}

La méfiance entre puissances se traduit d'abord par une \cle{course aux armements} : chaque État augmente ses dépenses militaires, perfectionne ses canons, ses cuirassés (comme le \emph{Dreadnought} britannique lancé en 1906), agrandit ses armées et allonge la durée du service militaire. L'Allemagne et la France portent leur armée à plusieurs millions d'hommes mobilisables. Chaque avancée militaire d'un camp provoque une réponse de l'autre : c'est une spirale qui rend la guerre de plus en plus probable, car les états-majors élaborent des plans de mobilisation rapide (comme le plan Schlieffen allemand) qu'il sera difficile d'arrêter une fois lancés.

Ensuite, les puissances se regroupent en deux blocs armés, censés se dissuader mutuellement.

\begin{definition}[La Triplice et la Triple-Entente]
La \cle{Triplice} (ou Triple-Alliance), formée en 1882, réunit l'\textbf{Allemagne}, l'\textbf{Autriche-Hongrie} et l'\textbf{Italie}. C'est une alliance défensive : si l'un des membres est attaqué, les autres doivent lui venir en aide.

La \cle{Triple-Entente}, constituée entre 1893 et 1907, réunit la \textbf{France}, la \textbf{Russie} et le \textbf{Royaume-Uni}. Ce n'est pas une alliance militaire aussi stricte, mais un ensemble d'accords qui rapprochent les trois pays face au bloc allemand.
\end{definition}

L'Europe se retrouve donc divisée en deux camps armés jusqu'aux dents. Le paradoxe est le suivant : ces alliances étaient censées maintenir la paix par la dissuasion, mais elles vont transformer un conflit local en guerre générale, car chaque État est lié à ses alliés par des promesses d'assistance.

\begin{popfigure}
\begin{center}
\begin{tikzpicture}[scale=1, every node/.style={font=\small}]
% Contours schématiques de l'Europe
\draw[thick, popmuted, rounded corners=4pt] (0,0) -- (0,4.6) -- (2.2,5.2) -- (4.5,5.0) -- (6.5,5.4) -- (8.6,4.8) -- (9.2,3.4) -- (8.4,1.8) -- (7.2,0.8) -- (5.0,0.2) -- (2.5,0.6) -- cycle;
% Pays Triple-Entente (bleu)
\fill[popblue!35, rounded corners=3pt] (0.4,1.6) -- (0.4,3.4) -- (1.8,3.6) -- (2.0,2.2) -- (1.2,1.4) -- cycle; % France
\fill[popblue!35] (0.3,4.0) -- (1.0,4.9) -- (1.7,4.4) -- (1.2,3.8) -- cycle; % GB
\fill[popblue!35, rounded corners=3pt] (6.6,2.2) -- (6.8,4.6) -- (8.4,4.4) -- (8.8,3.2) -- (8.0,2.0) -- cycle; % Russie
% Pays Triplice (orange)
\fill[poporange!45, rounded corners=3pt] (2.2,2.6) -- (2.4,4.2) -- (4.0,4.4) -- (4.2,2.8) -- (3.2,2.4) -- cycle; % Allemagne
\fill[poporange!45, rounded corners=3pt] (4.4,2.2) -- (4.4,3.4) -- (6.0,3.4) -- (6.2,2.4) -- (5.2,1.9) -- cycle; % Autriche-Hongrie
\fill[poporange!45, rounded corners=3pt] (3.0,0.5) -- (3.4,2.0) -- (4.0,1.8) -- (4.4,0.7) -- (3.8,0.3) -- cycle; % Italie
% Noms
\node at (1.2,2.5) {\textbf{France}};
\node[font=\footnotesize] at (1.0,4.4) {\textbf{R.-U.}};
\node at (7.6,3.3) {\textbf{Russie}};
\node at (3.3,3.4) {\textbf{Allemagne}};
\node[font=\footnotesize] at (5.3,2.7) {\textbf{Autriche-}};
\node[font=\footnotesize] at (5.3,2.35) {\textbf{Hongrie}};
\node[font=\footnotesize] at (3.7,1.15) {\textbf{Italie}};
% Balkans
\fill[poppink!50] (5.6,1.0) circle (0.28);
\node[font=\footnotesize, popdark] at (6.9,0.9) {Balkans};
\draw[-{Stealth}, popdark] (6.55,0.95) -- (5.95,1.02);
% Sarajevo
\fill[popink] (5.55,1.35) circle (0.07);
\node[font=\footnotesize, popink, anchor=west] at (5.7,1.45) {Sarajevo};
% Légende
\fill[popblue!35] (0.4,0.35) rectangle (0.8,0.6);
\node[anchor=west, font=\footnotesize] at (0.9,0.47) {Triple-Entente};
\fill[poporange!45] (3.0,-0.15) rectangle (3.4,0.1);
\node[anchor=west, font=\footnotesize] at (3.5,-0.03) {Triplice};
% Nord
\draw[-{Stealth}, thick, popdark] (8.9,5.0) -- (8.9,5.7);
\node[font=\footnotesize, popdark] at (8.9,5.9) {N};
\end{tikzpicture}
\end{center}
\legende{Croquis schématique de l'Europe en 1914 : les deux blocs d'alliances. En bleu, la Triple-Entente (France, Royaume-Uni, Russie) ; en orange, la Triplice (Allemagne, Autriche-Hongrie, Italie). Le point noir marque Sarajevo, dans les Balkans, la « poudrière de l'Europe ». Croquis simplifié, sans précision cartographique.}
\end{popfigure}

\courssub{Les nationalismes exacerbés et la « poudrière des Balkans »}

Troisième cause profonde : le \cle{nationalisme}, c'est-à-dire la volonté d'un peuple d'affirmer sa nation, son identité et sa puissance, souvent au détriment des autres. Au début du XX\textsuperscript{e} siècle, les nationalismes sont \emph{exacerbés} (poussés à l'extrême) : la presse attise les haines, on glorifie l'armée, on enseigne aux écoliers que la guerre serait « fraîche et joyeuse ».

Le nationalisme est particulièrement explosif dans les \textbf{Balkans}, région du sud-est de l'Europe que les contemporains surnomment la \emph{« poudrière de l'Europe »}. Pourquoi ce surnom ?

\begin{itemize}
\item L'\textbf{Empire ottoman}, qui dominait la région depuis des siècles, recule et abandonne des territoires (guerres balkaniques de 1912-1913).
\item Les peuples slaves de la région (Serbes, Bosniaques, Croates...) veulent leur indépendance ou leur regroupement. La \textbf{Serbie}, devenue indépendante, rêve de réunir tous les Slaves du Sud dans une « Grande Serbie ».
\item Or l'\textbf{Autriche-Hongrie}, empire multinational, a annexé en 1908 la Bosnie-Herzégovine, peuplée en partie de Serbes. Elle craint que le nationalisme serbe ne fasse éclater son empire.
\item La \textbf{Russie} se présente comme la protectrice des peuples slaves (panslavisme) et soutient la Serbie.
\end{itemize}

La mèche est donc posée : il ne manque qu'une étincelle pour faire exploser la poudrière.

\begin{methode}[Analyser les causes d'un conflit]
Pour analyser les causes d'un conflit dans une dissertation ou un commentaire de document, il faut toujours distinguer deux niveaux :
\begin{enumerate}
\item \textbf{Les causes profondes} (ou structurelles) : ce sont les tensions installées depuis longtemps. Pour 1914 : les rivalités impérialistes, la course aux armements, le système des alliances, les nationalismes.
\item \textbf{La cause immédiate} (ou déclencheur) : c'est l'événement précis qui met le feu aux poudres. Pour 1914 : l'attentat de Sarajevo.
\end{enumerate}
Une bonne copie montre que le déclencheur n'explique rien à lui seul : sans les causes profondes, l'étincelle n'aurait pas provoqué d'incendie général.
\end{methode}

\courssub{Le déclencheur : l'attentat de Sarajevo (28 juin 1914)}

Le 28 juin 1914, l'archiduc \cle{François-Ferdinand}, héritier du trône d'Autriche-Hongrie, effectue une visite officielle à Sarajevo, capitale de la Bosnie. Un jeune nationaliste serbe de Bosnie, \textbf{Gavrilo Princip}, membre du groupe révolutionnaire « Jeune Bosnie » soutenu par des réseaux nationalistes serbes, tire à deux reprises sur le couple princier : l'archiduc et son épouse Sophie sont tués.

Cet assassinat est un choc pour l'Europe. Mais rappelons-le : des attentats, il y en a eu d'autres avant 1914 sans provoquer de guerre mondiale. C'est le contexte de tensions accumulées qui transforme ce crime en catastrophe.

\begin{attention}[L'attentat n'est pas la seule cause de la guerre]
Erreur fréquente dans les copies : écrire que « la Première Guerre mondiale est causée par l'attentat de Sarajevo ». C'est faux. L'attentat n'est que \textbf{l'étincelle}, la cause immédiate. Les vraies causes sont profondes : rivalités impérialistes, course aux armements, alliances, nationalismes. Sans elles, l'attentat serait resté un drame local. Pense à l'image de la poudrière : l'étincelle n'est dangereuse que parce que la poudre est là.
\end{attention}

\courssub{Le jeu des alliances : les déclarations de guerre en chaîne (juillet-août 1914)}

Après l'attentat, les événements s'enchaînent en quelques semaines, comme des dominos qui tombent les uns après les autres. C'est ce qu'on appelle la \cle{crise de juillet} :

\begin{enumerate}
\item \textbf{23 juillet 1914} : l'Autriche-Hongrie, soutenue par l'Allemagne (le fameux « chèque en blanc » du 5 juillet), adresse un \textbf{ultimatum} humiliant à la Serbie, avec des exigences presque impossibles à accepter.
\item \textbf{28 juillet 1914} : jugeant la réponse serbe insuffisante, l'Autriche-Hongrie \textbf{déclare la guerre à la Serbie} et bombarde Belgrade.
\item \textbf{30 juillet} : la Russie, protectrice des Slaves, ordonne la \textbf{mobilisation générale} de son armée.
\item \textbf{1\textsuperscript{er} août} : l'Allemagne déclare la guerre à la \textbf{Russie}.
\item \textbf{3 août} : l'Allemagne déclare la guerre à la \textbf{France} ; le lendemain, elle envahit la \textbf{Belgique} neutre pour contourner les défenses françaises (plan Schlieffen).
\item \textbf{4 août} : la violation de la neutralité belge provoque l'entrée en guerre du \textbf{Royaume-Uni} contre l'Allemagne.
\end{enumerate}

En moins de dix jours, toutes les grandes puissances européennes sont en guerre. Le mécanisme des alliances, censé garantir la paix, a fonctionné comme une machine infernale : chaque déclaration en entraîne une autre. L'Italie, membre de la Triplice, refuse d'abord d'entrer en guerre (elle rejoindra la Triple-Entente en 1915), mais le conflit est déjà général.

\begin{popfigure}
\begin{center}
\begin{tikzpicture}[
  node distance=0.55cm and 1.1cm,
  etape/.style={draw=popink, thick, rounded corners=3pt, fill=popcream, text width=3.6cm, align=center, font=\footnotesize, inner sep=4pt},
  date/.style={font=\footnotesize\bfseries, text=popblue},
  fleche/.style={-{Stealth[length=2.5mm]}, very thick, poporange}
]
\node[etape, fill=poppinkL, align=center] (att){\textbf{28 juin 1914}\\ Attentat de Sarajevo : l'archiduc François-Ferdinand est assassiné};
\node[etape, below=of att, align=center] (ult){\textbf{23 juillet}\\ Ultimatum de l'Autriche-Hongrie à la Serbie};
\node[etape, below=of ult, align=center] (ah){\textbf{28 juillet}\\ L'Autriche-Hongrie déclare la guerre à la Serbie};
\node[etape, below left=0.55cm and 0.2cm of ah, align=center] (rus){\textbf{30 juillet}\\ Mobilisation générale de la Russie};
\node[etape, below right=0.55cm and 0.2cm of ah, align=center] (all){\textbf{1\textsuperscript{er} août}\\ L'Allemagne déclare la guerre à la Russie};
\node[etape, below=0.55cm of all, align=center] (fra){\textbf{3 août}\\ L'Allemagne déclare la guerre à la France et envahit la Belgique};
\node[etape, below=0.55cm of fra, align=center] (gb){\textbf{4 août}\\ Le Royaume-Uni déclare la guerre à l'Allemagne};
\draw[fleche] (att) -- (ult);
\draw[fleche] (ult) -- (ah);
\draw[fleche] (ah.south west) -- (rus.north);
\draw[fleche] (ah.south east) -- (all.north);
\draw[fleche] (rus.south) |- (fra.west);
\draw[fleche] (all) -- (fra);
\draw[fleche] (fra) -- (gb);
\end{tikzpicture}
\end{center}
\legende{De l'attentat de Sarajevo aux déclarations de guerre en chaîne (28 juin -- 4 août 1914). Chaque étape provoque la suivante par le jeu des alliances : la Russie soutient la Serbie, l'Allemagne soutient l'Autriche-Hongrie, la France est liée à la Russie, et le Royaume-Uni réagit à l'invasion de la Belgique.}
\end{popfigure}

\begin{exoresolu}[Distinguer cause profonde et cause immédiate]
\textbf{Énoncé.} Un élève écrit : « La Première Guerre mondiale a éclaté parce que Gavrilo Princip a tué l'archiduc François-Ferdinand à Sarajevo. » Corrige cette affirmation en distinguant les causes profondes et la cause immédiate du conflit.
\tcblower
\textbf{Solution détaillée.} L'affirmation est incomplète et donc fausse dans sa formulation. L'attentat de Sarajevo du 28 juin 1914 est bien l'événement déclencheur, mais il ne suffit pas à expliquer une guerre mondiale. Il faut distinguer :
\begin{enumerate}
\item \textbf{Les causes profondes} : les rivalités impérialistes entre la France, l'Allemagne et la Grande-Bretagne pour les colonies et les marchés ; la course aux armements ; la division de l'Europe en deux alliances (Triple-Entente contre Triplice) ; les nationalismes exacerbés, notamment dans les Balkans.
\item \textbf{La cause immédiate} : l'attentat de Sarajevo, qui déclenche la crise de juillet et les déclarations de guerre en chaîne.
\end{enumerate}
\textbf{Formulation correcte :} « La Première Guerre mondiale résulte de tensions profondes accumulées depuis des décennies (impérialisme, armements, alliances, nationalismes) ; l'attentat de Sarajevo n'a été que l'étincelle qui a mis le feu à cette poudrière. »
\end{exoresolu}

\courssub{Situation concrète : le Cameroun allemand pris en étau dès août 1914}

La guerre déclarée en Europe atteint immédiatement l'Afrique, car les empires coloniaux sont les prolongements des métropoles. Le cas du \cle{Kamerun} (Cameroun allemand depuis 1884) est exemplaire.

Observe la situation géographique de la colonie allemande en août 1914 :

\begin{itemize}
\item Au \textbf{nord et à l'ouest} : le \textbf{Nigeria}, colonie britannique ;
\item à l'\textbf{est et au sud} : l'\textbf{Afrique équatoriale française} (Tchad, Moyen-Congo, Gabon), colonie française ;
\item seule la \textbf{Guinée espagnole} (Río Muni), territoire d'une puissance neutre, offre une issue au sud.
\end{itemize}

Le Kamerun est donc \textbf{pris en étau} entre les possessions des deux principaux ennemis de l'Allemagne. Dès août 1914, les troupes françaises et britanniques attaquent la colonie depuis le Nigeria, le Tchad et le Congo. Douala, le grand port, tombe en septembre 1914 ; Yaoundé résiste jusqu'en janvier 1916 ; les derniers défenseurs allemands se replient vers la Guinée espagnole neutre en février 1916.

Pour mener cette guerre, les deux camps mobilisent massivement les populations locales : des milliers de Camerounais sont enrôlés comme \textbf{tirailleurs} (soldats) et surtout comme \textbf{porteurs}, chargés de transporter munitions, vivres et équipements à travers la forêt équatoriale. Les conditions sont terribles : épuisement, maladies, famines. Les pertes parmi les porteurs se comptent en milliers. La guerre européenne devient ainsi, concrètement, une tragédie camerounaise.

\begin{savaistu}[Le Cameroun, champ de bataille de la guerre mondiale]
Le Cameroun est l'un des rares territoires africains où des combats directs ont eu lieu entre armées européennes pendant la Première Guerre mondiale (1914-1916). Les forces françaises et britanniques y affrontent les troupes allemandes de la \emph{Schutztruppe}. Après la défaite allemande, la colonie est partagée : la France obtient la plus grande partie (le Cameroun oriental), la Grande-Bretagne une bande occidentale (le \emph{British Cameroons}). Ce partage, confirmé par un mandat de la Société des Nations en 1922, explique encore aujourd'hui la coexistence des héritages francophone et anglophone au Cameroun.
\end{savaistu}

\begin{aretenir}[L'essentiel de la section]
\begin{itemize}
\item La Première Guerre mondiale a des \textbf{causes profondes} : rivalités impérialistes (France, Allemagne, Grande-Bretagne), course aux armements, système des deux alliances (\cle{Triple-Entente} : France, Royaume-Uni, Russie ; \cle{Triplice} : Allemagne, Autriche-Hongrie, Italie) et nationalismes exacerbés, surtout dans les Balkans.
\item La \textbf{cause immédiate} est l'\textbf{attentat de Sarajevo} (28 juin 1914) contre l'archiduc François-Ferdinand.
\item Le jeu des alliances transforme la crise en guerre générale : déclarations en chaîne du 28 juillet au 4 août 1914.
\item Dès août 1914, le \textbf{Kamerun}, pris en étau entre les colonies françaises et britanniques, devient un champ de bataille ; des milliers de Camerounais sont enrôlés comme tirailleurs et porteurs.
\end{itemize}
\end{aretenir}

\coursec{Le déroulement du conflit : une guerre totale (1914-1918)}

La section précédente a montré comment le système d'alliances, les nationalismes exacerbés et la course aux armements ont transformé un attentat à Sarajevo en embrasement général. Nous allons maintenant suivre le conflit pour comprendre comment une guerre que tous croyaient courte est devenue une \cle{guerre totale} qui a broyé des nations entières et redessiné la carte du monde — y compris celle du Cameroun.

\courssub{la guerre de mouvement et l'échec du plan Schlieffen}

Au début août 1914, les états-majors européens appliquent des plans d'offensive conçus pour une guerre éclair. L'Allemagne met en œuvre le \cle{plan Schlieffen} : écraser la France en six semaines par une vaste manœuvre d'enveloppement par la Belgique, avant de reporter ses forces vers l'Est contre la Russie, jugée plus lente à mobiliser.

\begin{popfigure}
\begin{tikzpicture}[scale=0.9, every node/.style={font=\small}]
% Fond
\fill[popcream] (-0.5,-0.5) rectangle (15.5,5.5);
% Frontières simplifiées
\draw[popline, thick] (0,0) -- (14,0) -- (14,5) -- (0,5) -- cycle;
% Mer du Nord
\fill[popblueL] (0,0) rectangle (14,0.8);
\node[popdark] at (7,0.4) {Mer du Nord};
% Belgique
\fill[poporangeL] (2,0.8) rectangle (6,2.2);
\node[popdark] at (4,1.5) {Belgique};
% France
\fill[popblueL] (6,0.8) rectangle (14,5);
\node[popdark] at (10,3) {France};
% Allemagne
\fill[popredL] (0,2.2) rectangle (2,5);
\node[popdark, rotate=90] at (0.5,3.6) {Allemagne};
% Flèches plan Schlieffen
\draw[popred, very thick, ->] (1,3.5) .. controls (3,3) and (5,2) .. (7,1.5);
\node[popred, above] at (3,3.5) {Aile droite};
\draw[popred, thick, ->] (1,4) -- (2.5,4);
\node[popred, above] at (1.5,4.2) {Aile gauche};
% Contre-attaque française (Marne)
\draw[popblue, very thick, ->] (9,2) .. controls (8,2.5) and (7,3) .. (6.5,3.5);
\node[popblue, right] at (9,2) {Contre-offensive};
% Ligne de front stabilisée
\draw[popdark, dashed, thick] (6.5,1.5) -- (6.5,4.5);
\node[popdark, right] at (6.5,4.5) {Front stabilisé};
% Villes
\foreach \x/\y/\name in {3/1.2/Bruxelles, 5/1.5/Liège, 8/1.2/Reims, 10/1.5/Paris, 12/2/Verdun} {
  \fill[popdark] (\x,\y) circle (2pt);
  \node[popdark, anchor=north] at (\x,\y-0.3) {\name};
}
% Marne
\node[popblue, font=\bfseries] at (8.5,2.5) {Marne};
% Légende
\begin{scope}[xshift=0.5cm, yshift=-1.2cm]
\draw[popred, very thick, ->] (0,0) -- (2,0) node[right] {Plan Schlieffen (all.)};
\draw[popblue, very thick, ->] (0,-0.5) -- (2,-0.5) node[right] {Contre-offensive (fr.)};
\draw[popdark, dashed, thick] (0,-1) -- (2,-1) node[right] {Front 1915-1917};
\end{scope}
\end{tikzpicture}
\legende{Schéma simplifié du plan Schlieffen et de la bataille de la Marne (septembre 1914). L'aile droite allemande passe par la Belgique, mais est stoppée sur la Marne.}
\end{popfigure}

L'aile droite allemande progresse vite, mais elle s'éloigne de ses bases de ravitaillement. Le 3 septembre, le général Joffre ordonne la contre-offensive : c'est la \cle{bataille de la Marne} (5-12 septembre 1914). Les troupes françaises, renforcées par le corps expéditionnaire britannique, exploitent une brèche entre la 1\textsuperscript{re} et la 2\textsuperscript{de} armée allemandes. Les taxis de la Marne deviennent le symbole de la mobilisation civile : 600 véhicules réquisitionnés transportent environ 6 000 soldats vers le front. L'armée allemande recule vers l'Aisne et s'y retranche. La guerre de mouvement s'achève ; commence la guerre de positions.

\begin{definition}[Guerre de mouvement / Guerre de positions]
La \textbf{guerre de mouvement} se caractérise par des déplacements rapides de grandes masses de troupes, des fronts fluides et des batailles d'encerclement. La \textbf{guerre de positions} (ou guerre de tranchées) s'installe quand les deux camps s'enterrent face à face : le front se fige sur des centaines de kilomètres, l'artillerie domine, et toute offensive coûte des vies par milliers pour quelques mètres de terrain.
\end{definition}

\courssub{1915-1917 : la guerre de positions et des tranchées}

De la mer du Nord à la frontière suisse, une ligne continue de tranchées, de fils de fer barbelés et de blockhaus barre l'Europe. Le soldat vit dans la boue, les rats, la vermine, sous le pilonnage permanent. L'artillerie cause environ 70 \% des pertes. Pour percer, les états-majors lancent des offensives massives, précédées de préparations d'artillerie de plusieurs jours qui détruisent l'effet de surprise.

\begin{exemplebox}[Verdun et la Somme : l'usure industrielle]
\textbf{Verdun (février-décembre 1916).} L'Allemagne attaque pour « saigner la France à blanc ». La bataille dure 300 jours. Environ 700 000 hommes (morts, blessés, disparus) des deux côtés pour un front de 20 km. Le village de Fleury-devant-Douaumont change 16 fois de main. La \cle{Voie sacrée} (route Bar-le-Duc -- Verdun) assure le ravitaillement : un camion toutes les 14 secondes à son apogée.

\textbf{La Somme (juillet-novembre 1916).} Offensive franco-britannique pour soulager Verdun. Premier jour : 58 000 Britanniques hors de combat (dont 20 000 morts). Premiers chars britanniques (Mark I) engagés en septembre : lents, peu fiables, mais effet psychologique fort. Gain territorial : 10 km pour 1,2 million de pertes totales.
\end{exemplebox}

\begin{aretenir}[Nouvelles armes, nouvelle horreur]
\begin{itemize}
\item \textbf{Gaz de combat :} première utilisation massive à Ypres (avril 1915, chlore) ; puis moutarde (ypérite, 1917) : brûlures, cécité, suffocation lente.
\item \textbf{Chars :} britanniques (1916), français Schneider et Saint-Chamond (1917), Renault FT (1918, premier char moderne avec tourelle rotative).
\item \textbf{Aviation :} reconnaissance, puis chasse (Guynemer, Richthofen), bombardement stratégique naissant.
\item \textbf{Artillerie lourde sur voie ferrée :} canons de 370 mm, 520 mm tirant à 30 km.
\end{itemize}
\end{aretenir}

\begin{attention}[Piège : confondre percée et rupture]
Une \textbf{percée} est une brèche dans la ligne ennemie ; une \textbf{rupture} suppose l'exploitation de cette brèche par des réserves mobiles (chars, cavalerie, camions) avant que l'ennemi ne referme le front. En 1915-1917, il y a des percées (ex. offensive Nivelle, avril 1917) mais jamais de rupture : les réserves manquent de mobilité, les communications sont coupées, l'artillerie adverse tire à vue. Au Probatoire, ne dites pas « l'offensive a percé le front » si elle n'a pas débouché sur une manœuvre en profondeur.
\end{attention}

\courssub{La guerre totale : mobilisation générale des sociétés}

La guerre de positions exige des quantités industrielles de munitions, d'obus, de fusils, de chaussures, de pain. L'économie bascule en \cle{économie de guerre} : l'État réquisitionne usines, matières premières, main-d'œuvre. C'est la définition même de la guerre totale.

\begin{definition}[Guerre totale]
Conflit qui \textbf{mobilise toutes les ressources humaines, économiques, scientifiques et morales d'un pays}, civils compris. La distinction entre front et arrière s'efface : l'usine devient une seconde ligne de front, la femme remplace l'ouvrier parti au combat, l'enfant ramasse des herbes pour l'alimentation, la propagande conditionne les esprits.
\end{definition}

\begin{experience}[Analyse de document : l'affiche « Emprunt de la Libération » (1918)]
\textbf{Document :} Affiche française montrant une femme au travail dans une usine d'obus, avec la légende : « Souscrivez pour la victoire. L'emprunt de la Libération. »

\textbf{Démarche :}
\begin{enumerate}
\item \textbf{Identifier la source :} affiche de propagande officielle, 1918, France.
\item \textbf{Repérer les indices visuels :} femme en bleu de travail, tour à métaux, obus en arrière-plan, visage déterminé, drapeau tricolore.
\item \textbf{Interpréter :} la femme remplace l'homme mobilisé ; l'effort financier (emprunt) est présenté comme le prolongement de l'effort industriel ; la « Libération » lie guerre et paix future.
\item \textbf{Conclure :} l'État mobilise l'épargne populaire et le travail féminin pour soutenir l'effort de guerre. L'arrière est militarisé.
\end{enumerate}
\end{experience}

\begin{center}
\begin{tabularx}{\linewidth}{|l|X|X|}\hline
\rowcolor{popblueL} \textbf{Acteur} & \textbf{Mobilisation} & \textbf{Exemple chiffré (France)} \\ \hline
Hommes & Conscription universelle, classes 1887-1919 & 8,5 millions mobilisés (sur 40 millions d'habitants) \\ \hline
Femmes & Remplacement dans l'industrie, l'agriculture, les services & 430 000 « munitionnettes » dans les usines d'armement (1918) \\ \hline
Colonies & Tirailleurs combattants, porteurs, travailleurs & 180 000 tirailleurs sénégalais ; 50 000 Indochinois ; main-d'œuvre en métropole \\ \hline
Usines & Réquisition, conversion, travail aux pièces, 12 h/jour & Production annuelle d'obus : ~300 000 (1914) $\to$ ~50 millions (1918) \\ \hline
Finances & Emprunts nationaux, impôts, inflation & 4 emprunts nationaux (1915-1918) : 55 milliards de francs \\ \hline
Propagande & Censure, bourrage de crânes, union sacrée & « La guerre sera courte » (août 1914) $\to$ « Tenir jusqu'à la victoire » (1917) \\ \hline
\end{tabularx}
\end{center}

\legende{Tableau récapitulatif de la mobilisation totale en France (données comparables pour les autres belligérants). Source : synthèse manuels MINESEC.}

\begin{savaistu}[Les porteurs camerounais, force invisible]
Pendant la campagne du Cameroun (1914-1916), pour chaque soldat combattant, il faut 3 à 4 porteurs pour acheminer vivres, munitions, matériel médical sur des pistes inexistantes. Des dizaines de milliers de Camerounais sont réquisitionnés (souvent de force via la \cle{corvée}) par les Allemands, puis par les Franco-Britanniques. Beaucoup meurent de dysenterie, de paludisme, d'épuisement, loin de tout hôpital. Leur sacrifice est rarement commémoré.
\end{savaistu}

\courssub{1917, année charnière : entrée des États-Unis, sortie de la Russie}

Deux événements inversent le rapport de forces.

\begin{itemize}
\item \textbf{Avril 1917 : entrée en guerre des États-Unis.} Provocations : guerre sous-marine à outrance (torpillage du \emph{Lusitania}, 1915 ; reprise février 1917), télégramme Zimmermann (alliance proposée au Mexique contre les USA). Le président Wilson obtient du Congrès la déclaration de guerre le 6 avril. Les États-Unis apportent : hommes frais (2 millions en 1918), puissance industrielle, crédit financier, moralité nouvelle (« guerre pour la démocratie »).
\item \textbf{Novembre 1917 : révolution bolchévique en Russie.} L'épuisement, la faim, les défaites (Tannenberg 1914, offensive Broussilov 1916) font basculer le régime. Lénine signe le \textbf{traité de Brest-Litovsk} (mars 1918) : la Russie perd Pologne, Ukraine, Pays baltes, Finlande. L'Allemagne récupère des divisions de l'Est pour l'offensive finale à l'Ouest.
\end{itemize}

\begin{methode}[Construire une frise chronologique de la Première Guerre mondiale]
\begin{enumerate}
\item \textbf{Choisir l'échelle :} 1 cm = 1 mois (soit 52 cm pour 1914-1918) ou 1 cm = 3 mois sur une page A4.
\item \textbf{Tracer l'axe horizontal :} gradué de août 1914 à novembre 1918.
\item \textbf{Placer les dates repères (événements structurants) :}
  \begin{itemize}
  \item Août 1914 : déclaration de guerre, plan Schlieffen, Marne.
  \item Février 1916 : Verdun ; Juillet 1916 : Somme.
  \item Avril 1917 : entrée USA ; Novembre 1917 : révolution russe.
  \item Mars 1918 : offensive Michael ; Juillet 1918 : seconde Marne.
  \item 11 novembre 1918 : armistice.
  \end{itemize}
\item \textbf{Hiérarchiser les événements :} au-dessus de l'axe : front occidental ; au-dessous : fronts secondaires, diplomatie, vie civile. Utiliser des codes couleur (rouge = offensives allemandes, bleu = alliées, vert = événements politiques).
\item \textbf{Légender :} titre, échelle, auteur, source.
\end{enumerate}
\end{methode}

\begin{popfigure}
\begin{tikzpicture}[scale=0.85, every node/.style={font=\small}]
% Axe principal
\draw[popdark, thick] (0,0) -- (14,0);
% Graduations principales (années)
\foreach \x/\year in {0/1914, 3.5/1915, 7/1916, 10.5/1917, 14/1918} {
  \draw[popdark, thick] (\x,-0.15) -- (\x,0.15);
  \node[popdark, anchor=north] at (\x,-0.3) {\year};
}
% Mois (ticks fins)
\foreach \x in {0.29,0.58,0.87,1.17,1.46,1.75,2.04,2.33,2.63,2.92,3.21,3.5,
                3.79,4.08,4.38,4.67,4.96,5.25,5.54,5.83,6.13,6.42,6.71,7.0,
                7.29,7.58,7.88,8.17,8.46,8.75,9.04,9.33,9.63,9.92,10.21,10.5,
                10.79,11.08,11.38,11.67,11.96,12.25,12.54,12.83,13.13,13.42,13.71} {
  \draw[popmuted] (\x,-0.08) -- (\x,0.08);
}
% Événements au-dessus (front ouest)
\foreach \x/\txt/\col in {
  0.15/Marne/popblue,
  3.5/Verdun/popred,
  5.5/Somme/popblue,
  7.5/Chemin des Dames/popred,
  10.5/Entrée USA/popblue,
  11.2/Russie sort/popred,
  12.2/Off. Michael/popred,
  13.0/2\textsuperscript{e} Marne/popblue,
  13.8/Armistice/popdark} {
  \draw[\col, thick] (\x,0.2) -- (\x,1.2);
  \node[\col, anchor=south west, rotate=45] at (\x,1.25) {\txt};
}
% Événements en-dessous (politique/colonies)
\foreach \x/\txt/\col in {
  0.5/Togoland pris/popgreen,
  1.0/Douala pris/popgreen,
  3.0/Yaoundé/popgreen,
  4.5/Retraite all. Guinée esp./popgreen,
  8.0/Brest-Litovsk/poppurple,
  11.5/14 pts Wilson/poppurple} {
  \draw[\col, thick] (\x,-0.2) -- (\x,-1.2);
  \node[\col, anchor=north west, rotate=-45] at (\x,-1.25) {\txt};
}
% Légende
\begin{scope}[yshift=-2.2cm]
\node[popblue] at (0,-0.2) {\textbullet\ Front ouest (alliés)};
\node[popred] at (4,-0.2) {\textbullet\ Front ouest (allemands)};
\node[popgreen] at (8,-0.2) {\textbullet\ Campagne Cameroun};
\node[poppurple] at (12,-0.2) {\textbullet\ Diplomatie / Politique};
\end{scope}
\end{tikzpicture}
\legende{Frise chronologique simplifiée de la Première Guerre mondiale (1914-1918). Au-dessus : grands tournants du front occidental. En dessous : campagne du Cameroun et événements politiques majeurs.}
\end{popfigure}

\courssub{l'offensive allemande échoue, l'armistice du 11 novembre 1918}

Fort de ses divisions revenues de l'Est, Ludendorff lance l'\cle{offensive du printemps} (opération Michael, 21 mars 1918) : percée de 60 km en quelques jours grâce aux \cle{troupes d'assaut} (Stosstruppen) et aux tirs de barrage roulants. Mais l'avance dépasse les lignes de ravitaillement, les chars alliés contre-attaquent, les Américains arrivent par centaines de milliers.

Le 18 juillet 1918, Foch (généralissime allié depuis avril) lance la \cle{seconde bataille de la Marne} : contre-offensive franco-américaine (chars Renault FT, aviation massive). L'armée allemande recule, ses alliés s'effondrent (Bulgarie sept., Ottoman oct., Autriche-Hongrie nov.). La marine allemande se mutine (Kiel, 3 nov.), la révolution éclate à Berlin (9 nov.), l'empereur Guillaume II abdique et s'enfuit aux Pays-Bas.

\begin{definition}[Armistice]
Accord \textbf{mettant fin aux combats sans être un traité de paix}. Il suspend les hostilités, fixe des conditions militaires (évacuation des territoires occupés, remise de matériel, occupation de la Rhénanie) mais ne règle pas les questions territoriales, réparations, responsabilités : celles-ci seront traitées par les traités de paix (Versailles, Saint-Germain, Trianon, Neuilly, Sèvres).
\end{definition}

\begin{attention}[Erreur fréquente : armistice $\neq$ traité de paix]
L'\textbf{armistice du 11 novembre 1918} (signé à 5 h 15 dans le wagon de Rethondes, entrée en vigueur 11 h) arrête les tirs. Le \textbf{traité de Versailles} est signé le \textbf{28 juin 1919} (5 ans jour pour jour après Sarajevo). Entre les deux : conférence de la paix, blocus allié maintenu, révolution allemande, traité de Brest-Litovsk annulé. Au Probatoire, si le sujet demande « les conséquences de la guerre », ne vous arrêtez pas à l'armistice : les traités de paix en sont la traduction juridique et politique.
\end{attention}

\courssub{La campagne du Cameroun (1914-1916)}

Le Cameroun, colonie allemande depuis 1884, devient un théâtre d'opérations dès août 1914. Les forces franco-britanniques (depuis le Nigeria, le Congo français, le Tchad) convergent vers le centre.

\begin{popfigure}
\begin{tikzpicture}[scale=0.9, every node/.style={font=\small}]
% Contour très simplifié du Cameroun (polygone)
\fill[popcream] (0,0) -- (10,0) -- (11,2) -- (10.5,5) -- (8,6.5) -- (5,6) -- (2,5) -- (1,2.5) -- cycle;
\draw[popline, thick] (0,0) -- (10,0) -- (11,2) -- (10.5,5) -- (8,6.5) -- (5,6) -- (2,5) -- (1,2.5) -- cycle;
% Villes
\foreach \x/\y/\name/\col in {
  2.5/1.2/Douala/popblue,
  6.5/3.5/Yaoundé/popred,
  8.5/4.5/Ngaoundéré/popdark,
  1.5/3.5/Édéa/popmuted,
  4.5/1.5/Kribi/popmuted} {
  \fill[\col] (\x,\y) circle (3pt);
  \node[\col, anchor=west] at (\x+0.2,\y) {\name};
}
% Axes d'invasion
% Depuis le sud (Congo français) -> Kribi -> Édéa -> Douala
\draw[popblue, thick, ->] (1.5,0.5) -- (4.5,1.5);
\node[popblue, font=\tiny] at (3,0.8) {Colonne fr. (Congo)};
% Depuis l'ouest (Nigeria britannique) -> Douala
\draw[poporange, thick, ->] (0.5,1.2) -- (2.5,1.2);
\node[poporange, font=\tiny] at (0.5,0.8) {Colonne br. (Nigeria)};
% Depuis le nord (Tchad) -> Ngaoundéré -> Yaoundé
\draw[popgreen, thick, ->] (9,5.5) -- (8.5,4.5) -- (6.5,3.5);
\node[popgreen, font=\tiny] at (9.5,5) {Colonne fr. (Tchad)};
% Retraite allemande vers la Guinée espagnole (Rio Muni)
\draw[popred, dashed, thick, ->] (6.5,3.5) .. controls (5,2) and (3,1) .. (1.5,0.5);
\node[popred, font=\tiny] at (4,2.5) {Retraite allemande};
% Flèche Nord
\begin{scope}[xshift=9.5cm, yshift=5.5cm]
\draw[popdark, thick, ->] (0,0) -- (0,0.8);
\node[popdark, font=\tiny, anchor=south] at (0,0.9) {N};
\end{scope}
% Échelle indicative
\draw[popdark] (0.5,-0.5) -- (2.5,-0.5);
\node[popdark, font=\tiny, anchor=north] at (1.5,-0.6) {$\sim$ 200 km};
\end{tikzpicture}
\legende{Croquis schématique de la campagne du Cameroun (1914-1916). Flèches bleues/oranges/vertes : colonnes d'invasion franco-britanniques. Flèche rouge pointillée : retraite allemande vers la Guinée espagnole (Rio Muni). Croquis non contractuel, orientation et distances approximatives.}
\end{popfigure}

\begin{itemize}
\item \textbf{Septembre 1914 :} prise de \cle{Douala} (capitale économique, port principal) par une colonne britannique venue du Nigeria et une colonne française venue du Congo. Le gouverneur allemand Karl Ebermaier retire son administration à Yaoundé.
\item \textbf{Janvier 1916 :} \cle{bataille de Yaoundé}. Les colonnes convergent. Les Allemands, à court de munitions et de vivres, évacuent la ville dans la nuit du 13 au 14 janvier.
\item \textbf{Février 1916 :} dernière résistance à Mora (extrême-nord). Le 18 février, les derniers soldats allemands franchissent la frontière vers la \cle{Guinée espagnole} (Rio Muni) où ils sont internés. Le Cameroun est partagé : mandat français (4/5) et mandat britannique (1/5, deux bandes au nord et au sud) confirmé par la SDN en 1922.
\end{itemize}

\begin{savaistu}[Le sort des porteurs camerounais]
Les porteurs (souvent recrutés de force par la \cle{corvée}) transportent charges de 50-60 kg sur des dizaines de kilomètres, pieds nus, sous la pluie, sans soins. La mortalité atteint 20-30 \% sur certaines colonnes. Après la guerre, aucun monument ne leur est dédié à Yaoundé ou Douala. Leur mémoire est portée par l'histoire orale et quelques travaux d'historiens contemporains (ex. Marc Michel, \emph{Les Africains et la Grande Guerre}).
\end{savaistu}

\courssub{Le bilan humain : une hécatombe sans précédent}

La Première Guerre mondiale tue environ \textbf{10 millions de militaires} et \textbf{7 à 10 millions de civils} (famines, maladies, déportations, génocide arménien 1915-1916, grippe espagnole 1918-1919). Des millions de « \cle{gueules cassées} » (blessés au visage, défigurés) rentrent chez eux, porteurs d'un stigmate visible de la guerre industrielle.

\begin{popfigure}
\begin{tikzpicture}
\begin{axis}[
  ybar,
  bar width=0.6cm,
  width=\linewidth,
  height=8cm,
  ymin=0,
  ymax=2.2,
  ylabel={Pertes militaires (millions)},
  ylabel style={font=\small},
  xtick=data,
  xticklabels={Allemagne,Russie,France,Autriche-Hongrie,Royaume-Uni},
  xticklabel style={rotate=30, anchor=east, font=\small},
  ytick={0,0.5,1,1.5,2},
  yticklabel style={font=\small},
  enlarge x limits=0.2,
  grid=major,
  grid style={popline},
  axis line style={popline},
  tick style={popline},
]
\addplot[popblue, fill=popblueL] coordinates {
  (0,2.0)   % Allemagne ~2,0 millions
  (1,1.8)   % Russie ~1,8 millions
  (2,1.4)   % France ~1,4 millions
  (3,1.2)   % Autriche-Hongrie ~1,2 millions
  (4,0.9)   % Royaume-Uni ~0,9 million
};
\end{axis}
\end{tikzpicture}
\legende{Pertes militaires (morts) des principaux belligérants (estimations historiques arrondies). Source : synthèse d'historiens (Chickering, Prost, Winter). Les pertes russes incluent la guerre civile 1917-1922 ; les pertes autrichiennes incluent tous les fronts.}
\end{popfigure}

\begin{exoresolu}[Analyse quantitative : comparer les taux de pertes]
\textbf{Énoncé :} La France compte 40 millions d'habitants en 1914. L'Allemagne en compte 67 millions. Calculer le taux de morts militaires pour 1 000 habitants pour chaque pays. Que conclure ?

\tcblower
\textbf{Solution détaillée :}
\begin{align*}
\text{France : } & \frac{1,4 \times 10^6}{40 \times 10^6} \times 1000 = 35 \text{ morts pour 1 000 hab.} \\
\text{Allemagne : } & \frac{2,0 \times 10^6}{67 \times 10^6} \times 1000 \approx 30 \text{ morts pour 1 000 hab.}
\end{align*}
\textbf{Conclusion :} Bien que l'Allemagne ait plus de morts en valeur absolue, la France paie un tribut proportionnellement plus lourd (35 ‰ vs 30 ‰). Ce saignement démographique explique en partie la volonté française de sécurité absolue en 1919 (Rhénanie, réparations, alliance avec la Pologne, ligne Maginot).
\end{exoresolu}

\begin{aretenir}[L'héritage de la guerre totale]
\begin{itemize}
\item \textbf{Démographique :} classes creuses (1914-1918), déficit de naissances $\approx$ 1,4 million en France ; déséquilibre sexes (veuves, célibataires).
\item \textbf{Psychologique :} traumatisme collectif, pacifisme (« plus jamais ça »), mais aussi revanchisme (Allemagne, Italie).
\item \textbf{Politique :} effondrement de 4 empires (allemand, austro-hongrois, russe, ottoman) ; naissance de la SDN ; émergence des États-Unis comme puissance mondiale.
\item \textbf{Économique :} l'Europe endettée (vis-à-vis des USA), industries reconverties, inflation galopante (Allemagne 1923).
\item \textbf{Colonial :} mandats de la SDN (ex-Cameroun, Togo, Tanganyika, Rwanda-Urundi, Syrie, Liban, Irak, Palestine) ; promesses non tenues d'autodétermination $\to$ nationalismes coloniaux.
\end{itemize}
\end{aretenir}

\begin{exercice}[Entraînement Probatoire : document d'époque]
\textbf{Document :} Extrait du discours de Clemenceau à la Chambre des députés, 20 novembre 1918 : « Nous avons gagné la guerre... Il nous reste à gagner la paix, et ce sera peut-être le plus difficile. »

\textbf{Questions :}
\begin{enumerate}
\item Situer le document (auteur, date, contexte immédiat).
\item Expliquer pourquoi « gagner la paix » est présenté comme plus difficile que gagner la guerre.
\item En vous appuyant sur vos connaissances, citez trois points de désaccord entre les vainqueurs lors de la conférence de la paix (1919).
\end{enumerate}
\end{exercice}

\begin{corrige}[Exercice n°1]
\begin{enumerate}
\item \textbf{Contexte :} Georges Clemenceau, président du Conseil français, s'adresse à la Chambre 9 jours après l'armistice. La France est victorieuse mais exsangue ; les troupes alliées occupent la Rhénanie ; la conférence de la paix va s'ouvrir à Paris (janvier 1919).
\item \textbf{Analyse :} « Gagner la paix » signifie transformer la victoire militaire en un ordre durable : traiter les frontières, les réparations, le désarmement allemand, la sécurité française, les aspirations nationales (Pologne, Tchécoslovaquie, Yougoslavie), les mandats coloniaux, la création de la SDN. Les intérêts divergent : France (sécurité, affaiblissement durable de l'Allemagne), Royaume-Uni (équilibre européen, reprise du commerce), USA (14 points Wilson, SDN, pas d'annexions), Italie (terres promises), Japon (Pacifique). L'opinion publique réclame « payer les Boches » ; l'économie exige des réparations rapides.
\item \textbf{Trois désaccords majeurs :}
  \begin{itemize}
  \item \textbf{Rhénanie :} Foch/Clemenceau veulent détachement ou occupation permanente ; Lloyd George et Wilson refusent (violation du droit des peuples) $\to$ compromis : occupation 15 ans, zone démilitarisée permanente.
  \item \textbf{Réparations :} France veut somme fixe élevée et paiement en nature (charbon, bois) ; USA veulent capacité de paiement réelle ; UK veut relance du commerce allemand. Montant final fixé seulement en 1921 (132 milliards marks-or).
  \item \textbf{Colonies allemandes :} France/UK veulent annexions déguisées ; Wilson impose le système des mandats (tutelle SDN) $\to$ mandats de classe A (Proche-Orient), B (Afrique), C (Pacifique).
  \end{itemize}
\end{enumerate}
\end{corrige}

La guerre totale a donc produit un monde nouveau, instable, où les vainqueurs sont divisés, les vaincus humiliés mais non détruits, et les colonies promises à une émancipation différée. C'est dans ce contexte que les États-Unis, premiers créanciers du monde, connaissent une prospérité spectaculaire dans les années 1920 — avant l'effondrement de 1929 que nous étudierons dans la section suivante.

\coursec{Les conséquences de la guerre et le nouvel ordre mondial}

La guerre totale qui vient de s'achever laisse l'Europe exsangue. Après avoir analysé le déroulement du conflit — de la guerre de mouvement aux tranchées, de l'extension mondiale à l'effondrement des empires centraux — nous devons maintenant mesurer l'ampleur du choc. Le silence des canons ne marque pas la fin de l'histoire : il ouvre une période de négociations diplomatiques intenses, de redécoupage territorial radical et de transformations économiques et sociales profondes. C'est dans les traités de paix, signés entre 1919 et 1920, que se dessine le visage du monde de l'entre-deux-guerres, avec ses espoirs de paix durable et, déjà, les germes d'un nouveau conflit mondial.

\courssub{Les traités de paix : un redécoupage radical de l'Europe et du monde}

\courssub{Le traité de Versailles et le système des traités de 1919-1920}

La conférence de la paix s'ouvre à Paris en janvier 1919. Les vainqueurs — la France (Clemenceau), le Royaume-Uni (Lloyd George), les États-Unis (Wilson), l'Italie (Orlando) — dictent leurs conditions aux vaincus, exclus des négociations. Le traité le plus célèbre, celui de \textbf{Versailles}, est signé le \textbf{28 juin 1919} dans la galerie des Glaces, lieu symbolique où l'Empire allemand avait été proclamé en 1871. Il ne concerne que l'Allemagne. Quatre autres traités complètent le dispositif : Saint-Germain (Autriche, sept. 1919), Trianon (Hongrie, juin 1920), Neuilly (Bulgarie, nov. 1919) et Sèvres (Empire ottoman, août 1920, révisé à Lausanne en 1923).

\begin{definition}[Traité de Versailles]
Traité de paix signé le 28 juin 1919 entre les Alliés et l'Allemagne. Il impose à l'Allemagne la responsabilité de la guerre, des pertes territoriales lourdes, un désarmement strict et le paiement de réparations financières massives.
\end{definition}

\begin{exemplebox}[Les clauses territoriales majeures de Versailles]
\begin{itemize}
    \item \textbf{Alsace-Lorraine} rendue à la France (perdue en 1871).
    \item \textbf{Posnanie et corridor de Dantzig} cédés à la Pologne (accès à la mer), Dantzig ville libre sous SDN.
    \item \textbf{Sarre} placée sous administration de la SDN pour 15 ans (charbon à la France).
    \item \textbf{Rhénanie} démilitarisée et occupée 15 ans.
    \item \textbf{Perte de toutes les colonies} allemandes (reparties sous mandats).
    \item \textbf{Anschluss} (union avec l'Autriche) interdite.
\end{itemize}
\end{exemplebox}

\begin{popfigure}
\begin{tikzpicture}[scale=0.9, every node/.style={font=\small}]
    % Titre
    \node[align=center, font=\bfseries\large, text=popdark] at (0, 5.5) {L'Europe avant 1914 et après les traités de 1919-1920};
    % Légende
    \begin{scope}[yshift=-3.5cm]
        \draw[popline, thick] (-5,0) -- (-4.5,0);
        \node[right] at (-4.5,0) {Frontières 1914 (empires)};
        \draw[popblue, thick, fill=popblue!20] (-5,-0.6) rectangle (-4.5,-0.1);
        \node[right] at (-4.5,-0.35) {Nouvelles frontières 1920};
        \draw[poporange, thick, fill=poporange!20] (-5,-1.2) rectangle (-4.5,-0.7);
        \node[right] at (-4.5,-0.95) {Pays créés / agrandis};
        \fill[popgreen!30] (-5,-1.8) rectangle (-4.5,-1.3);
        \node[right] at (-4.5,-1.55) {Territoires perdus par les vaincus};
    \end{scope}
    % CARTE AVANT 1914 (gauche)
    \begin{scope}[xshift=-7cm, yshift=0.5cm]
        \node[align=center, font=\bfseries, text=popdark] at (0, 4) {EUROPE EN 1914};
        % Empires : contours simplifiés
        % Empire allemand
        \draw[popline, thick, fill=popblue!10] (-2.5,2) .. controls (-3,1) and (-3,-1) .. (-2.5,-2) .. controls (-1.5,-2.5) and (0.5,-2.5) .. (1.5,-1.5) .. controls (2.5,-0.5) and (2.5,1.5) .. (1.5,2.5) .. controls (0,3) and (-1,2.5) .. (-2.5,2) -- cycle;
        \node[popdark, font=\small\bfseries] at (-0.5,0.5) {EMPIRE ALLEMAND};
        % Empire austro-hongrois
        \draw[popline, thick, fill=poporange!10] (1.5,-1.5) .. controls (2.5,-2) and (4,-2) .. (4.5,0) .. controls (4,1.5) and (3,2) .. (2,1.5) .. controls (1.5,1) and (1,-1) .. (1.5,-1.5) -- cycle;
        \node[popdark, font=\small\bfseries] at (3,0) {AUTRICHE-HONGRIE};
        % Empire russe (partie ouest)
        \draw[popline, thick, fill=poppurple!10] (-2.5,-2) .. controls (-3,-3) and (-1,-3.5) .. (2,-3) .. controls (3,-2) and (2.5,0) .. (1.5,-1.5) .. controls (0.5,-2.5) and (-1.5,-2.5) .. (-2.5,-2) -- cycle;
        \node[popdark, font=\small\bfseries] at (-1.5,-2.5) {EMPIRE RUSSE};
        % Empire ottoman (partie europeenne)
        \draw[popline, thick, fill=poppink!10] (4.5,0) .. controls (5,-1) and (5,-2) .. (3.5,-2.5) .. controls (3,-1.5) and (4,0) .. (4.5,0) -- cycle;
        \node[popdark, font=\tiny\bfseries] at (4.2,-1) {OTTOMAN};
        % France / UK / Italie (simplifiés)
        \draw[popline, thick, fill=popgreen!10] (-2.5,2) -- (-4,2) -- (-4,1) -- (-3,1) -- (-2.5,0.5) -- cycle;
        \node[popdark, font=\tiny] at (-3.5,1.5) {FRANCE};
        \draw[popline, thick, fill=popgreen!10] (-4,2) -- (-4,3.5) -- (-2,3.5) -- (-2,2) -- cycle;
        \node[popdark, font=\tiny] at (-3,3) {UK};
    \end{scope}
    % CARTE APRES 1920 (droite)
    \begin{scope}[xshift=7cm, yshift=0.5cm]
        \node[align=center, font=\bfseries, text=popdark] at (0, 4) {EUROPE EN 1920};
        % Allemagne reduite
        \draw[popblue, thick, fill=popblue!20] (-2,1.5) .. controls (-2.5,0.5) and (-2.5,-1) .. (-2,-1.8) .. controls (-1,-2.2) and (1,-2.2) .. (1.5,-1.2) .. controls (2,-0.2) and (2,1) .. (1,1.8) .. controls (0,2.2) and (-1,1.8) .. (-2,1.5) -- cycle;
        \node[popblue, font=\small\bfseries] at (-0.3,0) {ALLEMAGNE};
        % France avec Alsace-Lorraine
        \draw[popblue, thick, fill=popblue!20] (-2,1.5) -- (-4,1.5) -- (-4,2.5) -- (-2.5,2.5) -- (-2,1.8) -- cycle;
        \node[popdark, font=\tiny] at (-3.2,2) {FRANCE};
        % Pologne (corridor)
        \draw[popgreen, thick, fill=popgreen!20] (1.5,-1.2) .. controls (2.5,-1.5) and (3,-0.5) .. (2.5,0.5) .. controls (1.5,1) and (0.5,0) .. (0.5,-1) -- cycle;
        \node[popgreen, font=\small\bfseries] at (1.8,-0.3) {POLOGNE};
        % Tchecoslovaquie
        \draw[poporange, thick, fill=poporange!20] (0.5,0.5) .. controls (1.5,1.5) and (3,1.5) .. (3.5,0) .. controls (3,-0.5) and (1.5,-0.5) .. (0.5,0) -- cycle;
        \node[poporange, font=\small\bfseries] at (1.8,0.5) {TCHÉCO.};
        % Autriche
        \draw[poporange, thick, fill=poporange!20] (3.5,0) .. controls (4,0.5) and (4.5,-0.5) .. (4,-1.5) .. controls (3.5,-1) and (3,-0.5) .. (3.5,0) -- cycle;
        \node[popdark, font=\tiny] at (3.8,-0.3) {AUTR.};
        % Hongrie
        \draw[poppurple, thick, fill=poppurple!20] (3.5,0) .. controls (4.5,0.5) and (5.5,-0.5) .. (5,-2) .. controls (4,-2) and (3,-1) .. (3.5,0) -- cycle;
        \node[poppurple, font=\small\bfseries] at (4.2,-1) {HONGRIE};
        % Yougoslavie / Roumanie (simplifié)
        \draw[poppink, thick, fill=poppink!20] (3,-1.5) .. controls (4,-2.5) and (2,-3) .. (0.5,-1.5) .. controls (1.5,-1) and (3,-1) .. (3,-1.5) -- cycle;
        \node[poppink, font=\small\bfseries] at (2.5,-2) {YOUGO. / ROUM.};
        % Pays Baltes / Finlande (haut droite)
        \draw[popline, thick, fill=popyellow!30] (0.5,2.5) -- (3,2.5) -- (3,3.5) -- (0.5,3.5) -- cycle;
        \node[popdark, font=\tiny] at (1.7,3) {PAYS BALTES / FINLANDE};
    \end{scope}
    % Flèches de transformation
    \draw[popdark, very thick, ->] (-0.5, 0) -- (0.5, 0) node[midway, above, font=\small\bfseries] {TRAITÉS 1919-1920};
\end{tikzpicture}
\legende{Figure 1 : Schéma comparatif du redécoupage européen. Les empires multinationaux (allemand, austro-hongrois, russe, ottoman) disparaissent au profit d'États-nations (Pologne, Tchécoslovaquie, Yougoslavie, Pays baltes) et d'États agrandis (France, Roumanie, Italie). L'Allemagne perd 13\% de son territoire et 10\% de sa population.}
\end{popfigure}

\courssub{Les « dommages de guerre » et le fardeau des réparations}

L'article 231 du traité de Versailles (la « clause de culpabilité ») déclare l'Allemagne et ses alliés « responsables de tous les pertes et dommages » subis par les Alliés. Cette base juridique permet d'exiger des \textbf{réparations}. Le montant final est fixé en 1921 à \textbf{132 milliards de marks-or} (équivalent à environ 33 milliards de dollars de l'époque), une somme colossale, payable en or, en nature (charbon, bois, bétail, navires) et en devises.

\begin{aretenir}[Conséquences des réparations sur l'Allemagne]
\begin{itemize}
    \item \textbf{Crise monétaire :} Pour payer, l'Allemagne imprime de la monnaie $\rightarrow$ hyperinflation de 1923 (1 dollar = 4 200 milliards de marks en nov. 1923). Épargne de la classe moyenne anéantie.
    \item \textbf{Rancœur nationaliste :} Le « diktat » de Versailles devient le moteur de la propagande nationaliste (Hitler : « honte de Versailles »).
    \item \textbf{Dépendance financière :} Plan Dawes (1924) puis Young (1929) \textbf{rééchelonnent} la dette via des prêts américains $\rightarrow$ l'économie allemande devient dépendante du capital US.
\end{itemize}
\end{aretenir}

\begin{attention}[Piège fréquent : Confondre « réparations » et « dommages de guerre »]
Les \textbf{dommages de guerre} sont les destructions matérielles réelles (usines, fermes, infrastructures) subies par la France, la Belgique, la Serbie surtout. Les \textbf{réparations} sont la compensation financière/juridique exigée du vaincu. Le montant des réparations (132 Md marks-or) dépasse largement la valeur estimée des dommages physiques réels. Au Bac, ne dites pas « l'Allemagne a payé pour tout reconstruire » : elle a payé une indemnité politique, en grande partie par emprunts américains qu'elle ne remboursera jamais totalement.
\end{attention}

\courssub{La Société des Nations (SDN) : l'espoir d'une paix par le droit}

Né de l'idéalisme wilsonien (14 points, point 14), la SDN voit le jour le \textbf{10 janvier 1920} (siège à Genève). C'est la première organisation internationale universelle visant à prévenir la guerre par l'arbitrage, le désarmement et la sécurité collective.

\begin{definition}[Société des Nations (SDN)]
Organisation internationale créée par le traité de Versailles (1919), entrée en vigueur en 1920. Objectif : maintenir la paix par la négociation, l'arbitrage, la sécurité collective (sanctions contre l'agresseur) et la coopération technique. Ancêtre direct de l'ONU (1945).
\end{definition}

\begin{popfigure}
\begin{tikzpicture}[scale=0.9, every node/.style={font=\small}]
    \node[align=center, font=\bfseries\large, text=popdark] at (0, 5) {Organisation et objectifs de la Société des Nations (SDN)};
    % Centre : Assemblée / Conseil / Secrétariat
    \node[draw=popblue, thick, fill=popblue!15, rounded corners, minimum width=3.5cm, minimum height=1.8cm, align=center] (centre) at (0, 1.5) {\textbf{ORGANES PRINCIPAUX}\\\textbf{Assemblée} (tous les États, 1 voix)\\\textbf{Conseil} (4 permanents + 4-9 élus)\\\textbf{Secrétariat} (administration permanente)};
    % Organes auxiliaires
    \node[draw=popgreen, thick, fill=popgreen!15, rounded corners, minimum width=3cm, minimum height=1.2cm, align=center] (bit) at (-5, -1) {\textbf{BIT}\\(Bureau Intl du Travail)\\\textit{Conditions de travail}};
    \node[draw=popgreen, thick, fill=popgreen!15, rounded corners, minimum width=3cm, minimum height=1.2cm, align=center] (cour) at (5, -1) {\textbf{CIJ}\\(Cour Intl de Justice)\\\textit{Arbitrage judiciaire}};
    \node[draw=poporange, thick, fill=poporange!15, rounded corners, minimum width=3.5cm, minimum height=1.2cm, align=center] (mandats) at (0, -2.8) {\textbf{COMMISSION DES MANDATS}\\\textit{Contrôle des territoires confiés (ex-colonies allemandes/ottomanes)}};
    % Objectifs (cercles externes)
    \foreach \ang/\txt/\col in {
        90/{Sécurité collective\\ (sanctions, Art. 16)}/popred,
        150/{Désarmement\\ (Conf. 1932 échec)}/poppurple,
        210/{Règlement pacifique\\ (arbitrage, CIJ)}/popblue,
        330/{Coopération technique\\ (santé, travail, réfugiés)}/popgreen,
        30/{Protection minorités\\ (traites spéciaux)}/poporange} {
        \node[draw=\col, thick, fill=\col!15, circle, minimum size=2.2cm, align=center] at (\ang:3.8) {\txt};
        \draw[\col, thick, ->] (\ang:2.8) -- (\ang:1.8);
    }
    % Flèches vers organes auxiliaires
    \draw[popblue, thick, ->] (centre.south west) -- (bit.north east);
    \draw[popblue, thick, ->] (centre.south east) -- (cour.north west);
    \draw[popblue, thick, ->] (centre.south) -- (mandats.north);
    % Membres absents (encadré)
    \node[draw=popred, thick, fill=popred!10, rounded corners, align=center, font=\small] at (-5, 4.5) {\textbf{GRANDS ABSENTS}\\\textbf{USA} (Sénat refuse)\\\textbf{Allemagne} (admise 1926, sort 1933)\\\textbf{URSS} (admise 1934, exclue 1939)\\\textbf{Japon/Italie} (sortent années 30)};
\end{tikzpicture}
\legende{Figure 2 : Structure de la SDN. Malgré une architecture complète (Assemblée, Conseil, CIJ, BIT, Mandats), l'absence des USA, le vote à l'unanimité (droit de veto de fait) et l'absence de force militaire propre la paralysent face aux crises des années 1930.}
\end{popfigure}

\courssub{Le Cameroun sous mandats : un laboratoire du « mandat » de la SDN}

Le sort des colonies allemandes illustre l'hypocrisie du « droit des peuples à disposer d'eux-mêmes » prôné par Wilson. L'article 22 du Pacte de la SDN crée le \textbf{système des mandats} : les territoires arrachés aux vaincus ne sont pas annexés, mais « confiés » à des puissances mandataires (France, UK, Japon, Belgique, Afrique du Sud, Australie, NZ) qui doivent les administrer « au nom de la SDN » et rendre compte annuellement.

\begin{definition}[Mandat de la SDN]
Statut juridique international (art. 22 Pacte SDN) : un territoire (ex-colonie allemande ou ottomane) est placé sous l'administration temporaire d'une puissance victorieuse (mandataire), qui a l'obligation de favoriser le bien-être et le développement des populations et de rendre compte chaque année à la Commission des mandats de la SDN. Ce n'est \textbf{pas} une annexion ni une colonie classique.
\end{definition}

\begin{attention}[Erreur fréquente au Probatoire : « Le Cameroun devient français en 1919 »]
\textbf{FAUX.} En 1919, le Kamerun allemand est \textbf{occupé militairement} par les Franco-Britanniques. Ce n'est qu'en \textbf{juillet 1922} que la SDN officialise le partage en \textbf{deux mandats distincts} :
\begin{itemize}
    \item \textbf{Mandat français (Cameroun oriental) :} $\approx$ 4/5 du territoire (420 000 km$^2$), capitale Yaoundé.
    \item \textbf{Mandat britannique (Cameroons) :} $\approx$ 1/5 du territoire (88 000 km$^2$), deux bandes frontalières du Nigeria (Nord et Sud), administré depuis Lagos.
\end{itemize}
Dire « colonie française » occulte l'obligation de rendre compte à la SDN et la différence de statut juridique (pas de droit de naturalisation massive, pas d'annexion pure).
\end{attention}

\begin{exemplebox}[Chiffres clés du partage du Cameroun (1922)]
\begin{center}
\begin{tabularx}{\linewidth}{|l|X|X|}\hline
\textbf{Critère} & \textbf{Mandat français (Est)} & \textbf{Mandat britannique (Ouest)} \\ \hline
Superficie & $\approx$ 420 000 km$^2$ (4/5) & $\approx$ 88 000 km$^2$ (1/5) \\ \hline
Population (est. 1922) & $\approx$ 1 500 000 hab. & $\approx$ 400 000 hab. \\ \hline
Capitale administrative & Yaoundé (dès 1921) & Buéa (Sud) / Garoua (Nord) — géré depuis Lagos \\ \hline
Politique d'administration & Assimilationniste (droit français, code de l'indigénat) & Indirect Rule (chefferies traditionnelles, droit coutumier) \\ \hline
Économie & Plantations (cacao, café, banane), ports (Douala), chemin de fer & Intégré à l'économie nigériane, cultures vivrières, peu d'infrastructures lourdes \\ \hline
\end{tabularx}
\end{center}
\end{exemplebox}

\begin{popfigure}
\begin{tikzpicture}[scale=1.1, every node/.style={font=\small}]
    % Contour Cameroun simplifié (polygone approximatif)
    \coordinate (N) at (0, 4.5);
    \coordinate (NE) at (2.5, 3.5);
    \coordinate (E) at (3.5, 1);
    \coordinate (SE) at (3, -1);
    \coordinate (S) at (1, -2.5);
    \coordinate (SW) at (-1.5, -2);
    \coordinate (W) at (-3, 0);
    \coordinate (NW) at (-2.5, 2.5);
    % AVANT 1914 (Gauche)
    \begin{scope}[xshift=-6cm]
        \node[align=center, font=\bfseries, text=popdark] at (0, 5) {KAMERUN ALLEMAND (1914)};
        \draw[popline, very thick, fill=poporange!20] (N) -- (NE) -- (E) -- (SE) -- (S) -- (SW) -- (W) -- (NW) -- cycle;
        \node[poporange, font=\bfseries\large] at (0, 1) {KAMERUN};
        \node[popdark] at (0, -0.5) {Gouvernorat : Buea puis Yaoundé};
        \node[popdark] at (0, -1) {Unité administrative unique};
        % Villes
        \fill[popdark] (0.5, 2) circle (1.5pt) node[right] {Douala};
        \fill[popdark] (0, 0.5) circle (1.5pt) node[right] {Yaoundé};
        \fill[popdark] (-1.5, 1.5) circle (1.5pt) node[left] {Dschang};
        \fill[popdark] (1.5, -0.5) circle (1.5pt) node[right] {Nkongsamba};
        \fill[popdark] (-0.5, -1.5) circle (1.5pt) node[right] {Ebolowa};
        % Nord
        \fill[popdark] (0, 3.5) circle (1.5pt) node[above] {Garoua};
        \fill[popdark] (-1, 3) circle (1.5pt) node[left] {Maroua};
        % Flèche Nord
        \draw[popdark, thick, ->] (-3.5, 4) -- (-3, 4.3) node[above] {N};
    \end{scope}
    % APRES 1922 (Droite)
    \begin{scope}[xshift=6cm]
        \node[align=center, font=\bfseries, text=popdark] at (0, 5) {MANDATS FRANÇAIS ET BRITANNIQUE (1922)};
        % Mandat Français (Est) - 4/5
        \draw[popblue, very thick, fill=popblue!15] (N) -- (NE) -- (E) -- (SE) -- (S) -- (0, -0.5) -- (0, 2.5) -- cycle;
        \node[popblue, font=\bfseries] at (1, 1.5) {MANDAT FRANÇAIS};
        \node[popblue, font=\small] at (1, 1) {(Cameroun Oriental)};
        \node[popblue, font=\small] at (1, 0.5) {$\approx$ 4/5 du territoire};
        \node[popblue, font=\small] at (1, 0) {Capitale : Yaoundé};
        % Mandat Britannique Sud (Ouest)
        \draw[popgreen, very thick, fill=popgreen!15] (0, 2.5) -- (0, -0.5) -- (S) -- (SW) -- (W) -- (NW) -- (N) -- cycle;
        \node[popgreen, font=\bfseries] at (-1.5, 1.5) {MANDAT BRITANNIQUE};
        \node[popgreen, font=\small] at (-1.5, 1) {(Cameroons du Sud)};
        \node[popgreen, font=\small] at (-1.5, 0.5) {$\approx$ 1/5 du territoire};
        \node[popgreen, font=\small] at (-1.5, 0) {Admin. depuis Lagos};
        % Ligne de partage (Milner-Simon 1919)
        \draw[popred, dashed, very thick] (0, 2.5) -- (0, -0.5);
        \node[popred, font=\small\bfseries, rotate=90] at (0.3, 1) {LIGNE MILNER-SIMON (1919)};
        % Villes (mêmes coords)
        \fill[popdark] (0.5, 2) circle (1.5pt) node[right] {Douala (Fr)};
        \fill[popdark] (0, 0.5) circle (1.5pt) node[right] {Yaoundé (Fr)};
        \fill[popdark] (-1.5, 1.5) circle (1.5pt) node[left] {Dschang (Fr)};
        \fill[popdark] (-0.5, -1.5) circle (1.5pt) node[right] {Ebolowa (Fr)};
        \fill[popdark] (0, 3.5) circle (1.5pt) node[above] {Garoua (Fr)};
        \fill[popdark] (-1, 3) circle (1.5pt) node[left] {Maroua (Fr)};
        % Nord britannique (exclave)
        \draw[popgreen, very thick, fill=popgreen!15] (-3.5, 5) rectangle (-2.5, 3.5);
        \node[popgreen, font=\tiny, rotate=90] at (-3, 4.25) {NORD BRITANNIQUE};
        \fill[popdark] (-3, 4.5) circle (1pt) node[left] {Yola};
        % Flèche Nord
        \draw[popdark, thick, ->] (3.5, 4) -- (4, 4.3) node[above] {N};
    \end{scope}
    % Légende commune
    \begin{scope}[yshift=-4.5cm]
        \draw[poporange, thick, fill=poporange!20] (-4,0) rectangle (-3.5,0.4) node[right=0.2cm] {Kamerun 1914 (unifié)};
        \draw[popblue, thick, fill=popblue!15] (-4,-0.6) rectangle (-3.5,-0.2) node[right=0.2cm] {Mandat français (Est)};
        \draw[popgreen, thick, fill=popgreen!15] (-4,-1.2) rectangle (-3.5,-0.8) node[right=0.2cm] {Mandat britannique (Ouest, 2 zones)};
        \draw[popred, dashed, thick] (-4,-1.8) -- (-3.5,-1.8) node[right=0.2cm] {Frontière de partage 1919};
    \end{scope}
\end{tikzpicture}
\legende{Figure 3 : Le Kamerun allemand (1914) et sa partition en mandats français (Est, bleu) et britannique (Ouest, vert) en 1922. La ligne Milner-Simon (pointillés rouges) coupe des groupes ethniques (Bamileke, Bassa, Douala). Le Nord britannique est une exclave administrée depuis le Nigeria. Ces frontières dessinent le Cameroun actuel.}
\end{popfigure}

\begin{savaistu}[Les frontières du Cameroun actuel naissent en 1919/1922]
C'est le tracé de la ligne Milner-Simon (1919), confirmée par la SDN (1922), qui dessine l'essentiel des frontières internationales du Cameroun actuel (avec le Nigeria, le Tchad, la RCA, le Congo, le Gabon, la Guinée équatoriale). La réunification de 1961 (indépendance + plébiscite) recolle les deux morceaux du mandat, mais la frontière interne Est/Ouest structure encore l'administration, l'éducation (sous-systèmes francophone/anglophone) et le contentieux anglophone actuel.
\end{savaistu}

\courssub{Conséquences économiques : l'Europe endettée, les États-Unis hégémoniques}

La guerre a fonctionné comme un formidable transfert de richesse et de puissance de l'Europe vers les États-Unis.

\begin{propriete}[Bilan économique de la guerre pour les grandes puissances]
\begin{itemize}
    \item \textbf{Europe (France, UK, Allemagne) :} Destruction du capital productif (usines, mines, fermes), perte de marchés coloniaux, dette publique explosée (France : dette $\times$6, UK : dette $\times$11). L'Europe passe de créancière nette à débitrice nette vis-à-vis des USA.
    \item \textbf{États-Unis :} Production industrielle +20\% (1914-1918), or afflue (réserves d'or $\times$ 3), dollar devient monnaie de réserve mondiale. NY remplace Londres comme place financière n°1.
    \item \textbf{Japon :} Comble le vide laissé par l'Europe en Asie (textile, navale), expansion industrielle rapide.
\end{itemize}
\end{propriete}

\begin{exoresolu}[Analyse de la dépendance financière européenne vis-à-vis des USA (1919-1929)]
\textbf{Document :} Extrait des flux financiers internationaux (milliards de dollars constants 1913).
\begin{center}
\begin{tabular}{|l|c|c|c|}\hline
\textbf{Flux} & \textbf{1914-1918 (Guerre)} & \textbf{1919-1923 (Réparation/Dawes)} & \textbf{1924-1929 (Prospérité)} \\ \hline
Prêts US $\rightarrow$ Alliés (France, UK) & +10 (avances guerre) & +2 (stabilisation) & +15 (investissements) \\ \hline
Réparations Allemagne $\rightarrow$ Alliés & 0 & +2 (réel payé) & +8 (via plan Dawes/Young) \\ \hline
Remboursement dettes Alliés $\rightarrow$ USA & 0 & -0,5 (défaut partiel) & -2 (service dette) \\ \hline
Investissements US $\rightarrow$ Allemagne & 0 & +4 (plan Dawes) & +12 (usines, obligations) \\ \hline
\end{tabular}
\end{center}
\textbf{Question :} Montrez que le système financier de l'entre-deux-guerres repose sur une « pompe » circulaire fragile.
\textbf{Solution :}
\begin{enumerate}
    \item \textbf{Identification du cycle :} USA $\xrightarrow{\text{prêts}}$ Allemagne $\xrightarrow{\text{réparations}}$ France/UK $\xrightarrow{\text{remboursement dettes guerre}}$ USA.
    \item \textbf{Fragilité :} L'Allemagne ne paie ses réparations qu'avec l'argent emprunté aux USA (plan Dawes 1924, Young 1929). La France/UK ne rembourse ses dettes de guerre aux USA qu'avec les réparations allemandes.
    \item \textbf{Rupture 1929 :} Krach de Wall Street $\rightarrow$ arrêt des prêts US $\rightarrow$ effondrement allemand $\rightarrow$ fin réparations $\rightarrow$ défaut Alliés. Le système s'effondre car il n'y a pas de création de richesse réelle en Allemagne, seulement de la dette roulée.
\end{enumerate}
\end{exoresolu}

\courssub{Conséquences sociales : deuil, mutation du travail et émancipation féminine}

\begin{itemize}
    \item \textbf{Hécatombe démographique :} $\approx$ 10 millions de morts militaires, 20 millions de blessés (dont 6-7 millions d'invalides permanents : « gueules cassées », gazés, amputés). France : 1,4 M morts (10\% de la population masculine active), 600 000 veuves, 1 million d'orphelins.
    \item \textbf{Féminisation du travail :} Femmes remplacent les hommes aux usines (munitions), champs, transports, services publics. 1914-1918 : taux d'activité féminin bondit (France : 35\% $\rightarrow$ 45\%).
    \item \textbf{Émancipation partielle :} Droit de vote accordé aux femmes \textbf{après} la guerre dans plusieurs pays (Allemagne 1918, Autriche 1918, UK 1918 [>30 ans], USA 1920, URSS 1917). \textbf{Pas en France} (vote 1944). Retour forcé au foyer après 1918 (« réparation démographique »), mais le mythe de la « femme au foyer » est fissuré.
    \item \textbf{État-providence naissant :} Pensions d'invalidité, allocations veuves/orphelins, offices des mutilés $\rightarrow$ première grande intervention sociale de l'État dans la vie privée.
\end{itemize}

\begin{methode}[Comment analyser une statistique de pertes humaines au Bac]
\begin{enumerate}
    \item \textbf{Contextualiser :} Donnée brute (ex: 1,4 M morts français) $\rightarrow$ rapport à la population totale (40 M) et à la classe d'âge (hommes 20-40 ans).
    \item \textbf{Comparer :} Taux de mortalité par pays (Serbie 16\% pop., France 3,5\%, UK 2\%, USA 0,1\%).
    \item \textbf{Conséquences structurelles :} Creux des classes creuses (1915-1919), déficit de naissances ($\approx$ 1,4 M naissances « manquantes » en France), vieillissement accéléré.
    \item \textbf{Lier au politique :} Pacifisme des années 20, refus de la guerre en 1938-39 (Munich), natalisme d'État (code de la famille 1939).
\end{enumerate}
\end{methode}

\courssub{Les germes de la Seconde Guerre mondiale dans les traités de 1919}

\begin{aretenir}[Pourquoi Versailles prépare 1939 (synthèse Bac)]
\begin{enumerate}
    \item \textbf{Humiliation allemande :} Clause de culpabilité (art. 231), démantèlement territorial, désarmement unilatéral $\rightarrow$ mythe du « coup de poignard dans le dos » (Dolchstoßlegende), révisionnisme structurel.
    \item \textbf{Frontières intenables :} Corridor de Dantzig (sépare Prusse-Orient), minorités allemandes en Tchécoslovaquie (Sudètes) et Pologne $\rightarrow$ prétextes à l'expansionnisme hitlérien (1938-39).
    \item \textbf{Effondrement économique programmé :} Réparations impossibles $\rightarrow$ hyperinflation 1923 $\rightarrow$ dépendance aux prêts US $\rightarrow$ crash 1929 $\rightarrow$ chute de la République de Weimar $\rightarrow$ Hitler chancelier (1933).
    \item \textbf{SDN impuissante :} Pas d'armée, vote à l'unanimité, grands absents (USA, puis Allemagne, Japon, Italie, URSS) $\rightarrow$ échec face à Mandchourie (1931), Éthiopie (1935), Rhénanie (1936).
    \item \textbf{Colonies/mandats :} Promesses de Wilson trahies $\rightarrow$ nationalismes coloniaux (Indochine, Algérie, Cameroun/UPC) $\rightarrow$ décolonisation violente post-1945.
\end{enumerate}
\end{aretenir}

\begin{exercice}[Sujet type Probatoire — Analyse de conséquences]
\textbf{Document 1 :} Extrait du discours de Clemenceau à la Chambre (sept. 1919) : « Le traité n'est pas parfait... mais il assure à la France la sécurité... la Rhénanie, la Sarre, l'Alsace-Lorraine... »
\textbf{Document 2 :} Caricature allemande 1919 : l'Allemagne nue, enchaînée, saignée par les Alliés (Versailles).
\textbf{Document 3 :} Carte du Cameroun 1922 (mandats).
\textbf{Questions :}
\begin{enumerate}
    \item En vous appuyant sur les documents et vos connaissances, montrez que le traité de Versailles répond aux objectifs de sécurité de la France mais crée les conditions d'une instabilité durable en Europe. (6 pts)
    \item Expliquez en quoi le statut de mandat au Cameroun diffère de la colonisation classique et quelles en sont les conséquences administratives jusqu'en 1960. (4 pts)
    \item « La crise de 1929 est la conséquence lointaine des traités de 1919. » Discutez cette affirmation en structurant votre réponse. (10 pts)
\end{enumerate}
\end{exercice}

\begin{corrige}[Exercice n°3 - Éléments de corrigé]
\textbf{1.} \textit{Sécurité française :} Récupération Alsace-Lorraine, Rhénanie démilitarisée/occupée (glacis), Sarre (charbon), limitation armée allemande (100 000 hommes). \textit{Instabilité :} Humiliation allemande (culpabilité, réparations, pertes territoriales) $\rightarrow$ révisionnisme ; nouvelles frontières ethniquement hétérogènes (Tchécoslovaquie, Pologne, Yougoslavie) $\rightarrow$ conflits minoritaires ; SDN sans force $\rightarrow$ impuissance.
\textbf{2.} \textit{Mandat $\neq$ Colonie :} Mandat = tutelle internationale (SDN), obligation de rapport annuel, interdiction annexion/naturalisation forcée, préparation à l'autonomie (théorique). \textit{Conséquences Cameroun :} Double administration (française assimilationniste / britannique Indirect Rule), double système éducatif/juridique, frontière interne structurante, réunification 1961 par plébiscite (Nord $\rightarrow$ Nigeria, Sud $\rightarrow$ Cameroun).
\textbf{3.} \textit{Thèse :} Lien direct via chaîne : Réparations $\rightarrow$ Dawes/Young (dette US) $\rightarrow$ Krach 1929 $\rightarrow$ arrêt prêts US $\rightarrow$ effondrement Allemagne/Autriche $\rightarrow$ crise mondiale. \textit{Antithèse :} Causes propres à 1929 (surproduction US, bulle spéculative boursière, dérégulation bancaire, protectionnisme Hawley-Smoot). \textit{Synthèse :} Les traités créent la \textbf{fragilité structurelle} (dette, déséquilibres) qui transforme une crise boursière US en \textbf{Grande Dépression mondiale}. Sans Versailles/Dawes, la contagion aurait été moins violente en Europe.
\end{corrige}

\coursec{La prospérité américaine des années 1920 et ses fragilités}

La section précédente a mis en lumière les conséquences politiques et territoriales de la Grande Guerre : l'Europe sort exsangue, endettée, tandis que les États-Unis affirment leur rôle de première puissance économique et financière mondiale. C'est dans ce contexte de bascule du centre de gravité de l'économie mondiale vers l'Atlantique Nord que s'inscrit la prospérité américaine des « années folles ». Pourtant, derrière l'éclat de la croissance et de la modernité, se creusent des failles structurelles qui prépareront l'effondrement de 1929.

\courssub{L'essor économique et les « années folles » : modernité et production de masse}

Au lendemain de la guerre, les États-Unis connaissent une expansion économique sans précédent. Le produit national brut (PNB) augmente d'environ 40 \% entre 1921 et 1929. Cette prospérité repose sur une révolution industrielle de deuxième vague : l'électricité, l'automobile, la chimie et les nouvelles méthodes d'organisation du travail transforment en profondeur l'appareil productif et la vie quotidienne.

L'automobile devient le symbole de cette ère nouvelle. Henry Ford applique à grande échelle les principes du \cle{taylorisme} (organisation scientifique du travail : décomposition des tâches, chronométrage, suppression des gestes inutiles) et du \cle{fordisme} (chaîne de montage mobile, standardisation des pièces, salaires élevés pour créer une clientèle solvable). L'usine de Highland Park, puis celle de River Rouge, produisent la Ford T à un rythme effréné : un châssis sort de la chaîne toutes les 93 minutes en 1914, contre 12 heures auparavant. Le prix de la voiture passe de 850 dollars en 1908 à 260 dollars en 1925, la rendant accessible aux classes moyennes. En 1929, les États-Unis comptent plus de 23 millions d'automobiles, soit une pour 5 habitants.

\begin{exemplebox}[L'effet d'entraînement de l'automobile]
L'industrie automobile tire tout le secteur secondaire : acier (75 \% de la production), caoutchouc, verre, pétrole (essence), construction routière. Elle stimule aussi l'urbanisme (lotissements pavillonnaires en périphérie des villes) et le commerce (supermarchés accessibles en voiture).
\end{exemplebox}

L'électrification des foyers et des usines change le rythme de la production et de la vie domestique. La part de l'électricité dans l'énergie industrielle passe de 30 \% (1914) à 70 \% (1929). Les appareils ménagers (réfrigérateurs, machines à laver, aspirateurs, radios) se diffusent, libérant du temps domestique — surtout pour les femmes — et créant de nouveaux marchés. Le cinéma (Hollywood) et la radio (premières chaînes nationales comme NBC en 1926) forgent une culture de masse commune, vecteur du « rêve américain ».

\begin{definition}[Taylorisme et fordisme]
Le \textbf{taylorisme} (Frederick Winslow Taylor, 1911) est une méthode d'organisation du travail fondée sur la division verticale (conception/exécution) et horizontale (parcèlement des tâches), le chronométrage et la prime au rendement. Le \textbf{fordisme} (Henry Ford, 1913-1914) associe le taylorisme à la chaîne de montage mobile, à la standardisation extrême des pièces et à une politique salariale volontariste (5 dollars/jour en 1914) pour stabiliser la main-d'œuvre et lui permettre d'acheter les produits qu'elle fabrique.
\end{definition}

\courssub{Les mécanismes de la spéculation boursière : l'illusion d'une richesse sans limite}

La croissance réelle de l'économie réelle (biens et services) s'accompagne, à partir de 1924-1925, d'une déconnexion progressive des marchés financiers. La Bourse de New York (Wall Street) devient le lieu d'une spéculation effrénée, alimentée par le crédit facile et la conviction partagée que « les actions ne font que monter ».

\begin{definition}[Spéculation boursière]
La \textbf{spéculation} consiste à acheter des titres (actions, obligations) non pas pour percevoir des dividendes ou participer à la gouvernance de l'entreprise, mais dans l'espoir exclusif de les revendre plus cher à court terme, réalisant une plus-value sans création de richesse productive nouvelle. Elle repose sur l'anticipation d'une hausse future des cours, souvent déconnectée des fondamentaux économiques (bénéfices, actifs réels).
\end{definition}

Le moteur de cette bulle est l'\cle{achat sur marge} (\emph{buying on margin}). L'investisseur ne verse qu'une fraction de la valeur de l'action (souvent 10 \% seulement, voir \emph{Le savais-tu ?}) ; le courtier avance le reste, les actions achetées servant de garantie (collatéral). Tant que les cours montent, la valeur du portefeuille augmente, permettant d'emprunter davantage pour acheter encore plus d'actions : c'est le cercle vertueux de la spéculation, qui devient vicieux dès que les cours fléchissent.

\begin{popfigure}
\begin{tikzpicture}[
    node distance=2.2cm and 2.5cm,
    box/.style={draw=popblue, thick, rounded corners, fill=popblueL, minimum width=3.2cm, minimum height=1.6cm, align=center, font=\small},
    arrow/.style={-Stealth, thick, popdark}
]
\node[box, align=center] (credit){Crédit facile\\et taux bas};
\node[box, right=of credit, align=center] (achat){Achat d'actions\\à crédit (marge 10 \%)};
\node[box, below=of achat, align=center] (hausse){Hausse artificielle\\des cours à Wall Street};
\node[box, left=of hausse, align=center] (emprunt){Nouvel emprunt\\garanti par les plus-values latentes};

\draw[arrow] (credit) -- node[above, font=\tiny, pos=0.5] {1. Liquidités abondantes} (achat);
\draw[arrow] (achat) -- node[right, font=\tiny, pos=0.5] {2. Demande massive} (hausse);
\draw[arrow] (hausse) -- node[below, font=\tiny, pos=0.5] {3. Patrimoine gonflé} (emprunt);
\draw[arrow] (emprunt) -- node[left, font=\tiny, pos=0.5] {4. Réinvestissement} (credit);

% Flèche centrale pour montrer le cercle
\draw[arrow, poporange, dashed, bend left=15] (emprunt) to node[left=2mm, font=\tiny, poporange] {Boucle spéculative} (credit);
\end{tikzpicture}
\legende{Schéma du cercle de la spéculation boursière (1924-1929) : le crédit alimente l'achat d'actions, qui fait monter les cours, ce qui augmente la valeur des garanties et permet un nouvel emprunt.}
\end{popfigure}

\begin{attention}[Piège fréquent : confusion cours boursiers et richesse réelle]
\emph{Erreur :} « La Bourse monte, donc le pays s'enrichit. »
\emph{Réalité :} La hausse des cours ne crée pas d'usines, ne forme pas d'ingénieurs, n'augmente pas la productivité. Elle ne fait que transférer du pouvoir d'achat futur (crédit) vers le présent, en gonflant artificiellement la valeur nominale du capital. Quand la bulle éclate, la richesse « papier » disparaît, mais les dettes contractées pour l'acheter restent bien réelles.
\end{attention}

\begin{popfigure}
\begin{tikzpicture}
\begin{axis}[
    width=\linewidth,
    height=8cm,
    xmin=1921, xmax=1933,
    ymin=0, ymax=400,
    xtick={1921,1923,1925,1927,1929,1931,1933},
    ytick={0,100,200,300,380},
    xlabel={Année},
    ylabel={Indice Dow Jones (base 100 en 1921)},
    grid=major,
    grid style={popline},
    axis line style={popdark},
    tick style={popdark},
    label style={font=\small, popdark},
    tick label style={font=\footnotesize, popdark},
    legend style={at={(0.02,0.98)}, anchor=north west, font=\footnotesize, fill=popcream, draw=popblue}
]
% Données approximatives DJIA (normalisé 1921=100)
% 1921: 71 -> 100, 1922: 98, 1923: 95, 1924: 120, 1925: 156, 1926: 157, 1927: 202, 1928: 300, 1929: 381 (pic sept), 1930: 164, 1931: 78, 1932: 41, 1933: 60
\addplot[popcurve, very thick, popblue, mark=*, mark size=2pt] coordinates {
    (1921,100) (1922,138) (1923,134) (1924,169) (1925,220) (1926,221) (1927,284) (1928,422) (1929,536) (1930,231) (1931,110) (1932,58) (1933,84)
};
\addlegendentry{Indice boursier (échelle relative)}

% Annotations
\node[anchor=south, font=\tiny, popdark] at (axis cs:1929,536) {Pic sept. 1929};
\node[anchor=north, font=\tiny, popdark] at (axis cs:1932,58) {Creux juil. 1932};
\draw[poporange, thick, -Stealth] (axis cs:1929.5,400) -- (axis cs:1929.8,500);
\draw[poporange, thick, -Stealth] (axis cs:1931.5,150) -- (axis cs:1932.2,80);
\end{axis}
\end{tikzpicture}
\legende{Évolution de l'indice boursier de Wall Street (Dow Jones Industrial Average, base 100 en 1921) de 1921 à 1933. La phase d'euphorie (1924-1929) voit l'indice multiplier par 5, suivie de l'effondrement brutal de 1929-1932 (perte de 90 \% de la valeur). Source : données historiques DJIA.}
\end{popfigure}

\courssub{Les déséquilibres structurels cachés : surproduction, salaires et endettement}

Derrière la façade de la prospérité, l'économie américaine accumule des contradictions majeures que la spéculation masque temporairement.

\begin{definition}[Surproduction]
La \textbf{surproduction} (ou surcapacité) est une situation où le volume de biens produits excède la demande solvable des consommateurs. Elle entraîne l'accumulation de stocks invendus, la baisse des prix, la réduction de la production et, in fine, le chômage. Elle signale un déséquilibre entre l'appareil productif (offre) et le pouvoir d'achat (demande).
\end{definition}

\noindent\textbf{1. Surproduction agricole (dès 1920-1921).} Pendant la guerre, les fermiers américains se sont suréquipés (tracteurs, engrais) pour nourrir l'Europe. Après 1918, la demande européenne s'effondre (retour à la production locale, protectionnisme). Les cours agricoles s'effondrent (blé : -50 \% entre 1919 et 1921). Le revenu agricole stagne ou baisse pendant toute la décennie, alors que la productivité augmente. Les fermiers s'endettent pour acheter du matériel ; en 1929, 40 \% des fermiers sont locataires ou très endettés.

\noindent\textbf{2. Surproduction industrielle (à partir de 1927-1928).} La productivité horaire augmente de 40 \% entre 1919 et 1929 (électrification, taylorisme), mais les salaires réels ne progressent que de 11 \%. L'écart se creuse : la part des salaires dans la valeur ajoutée baisse, celle des profits monte. Les ménages les plus aisés s'équipent (automobile, radio, réfrigérateur) ; le marché se sature. Les stocks s'accumulent (automobiles, acier, textile). La production automobile plafonne dès 1928.

\noindent\textbf{3. Faiblesse du pouvoir d'achat et inégalités.} En 1929, les 5 \% des foyers les plus riches perçoivent 30 \% du revenu national ; les 60 \% les plus pauvres en perçoivent à peine 25 \%. La consommation de masse repose sur le crédit à la consommation (achat à tempérament), qui explose : en 1929, 75 \% des automobiles, 80 \% des radios et 60 \% des meubles sont achetés à crédit. L'endettement des ménages double entre 1920 et 1929.

\begin{savaistu}[L'effet de levier de 10 \% : la ruine en quelques heures]
En 1929, la pratique courante du \emph{margin buying} permettait d'acheter 100 dollars d'actions en ne déposant que 10 dollars (10 \% de marge). Le courtier prêtait les 90 dollars restants, les actions servant de gage.
\begin{itemize}
    \item Si l'action monte à 110 dollars : gain de 10 dollars pour 10 dollars misés = \textbf{+100 \% de rendement}.
    \item Si l'action baisse à 90 dollars : perte de 10 dollars pour 10 dollars misés = \textbf{-100 \% de rendement} (capital propre anéanti).
\end{itemize}
Le 24 octobre 1929 (\emph{Jeudi noir}), puis le 29 octobre (\emph{Mardi noir}), les cours chutent si vite que les courtiers lancent des \emph{margin calls} (appels de marge) : l'investisseur doit remettre du cash immédiatement ou ses actions sont vendues de force au plus bas, le ruinant définitivement. Ce mécanisme a transformé une correction boursière en catastrophe financière systémique.
\end{savaistu}

\courssub{Les liens financiers transatlantiques : le cycle de la dette et des réparations}

La prospérité américaine ne s'arrête pas aux frontières. Les États-Unis, devenus premier créancier mondial, financent la reconstruction européenne et le règlement des réparations allemandes, créant un système circulaire fragile.

\begin{center}
\begin{tabularx}{\linewidth}{|l|X|X|}\hline
\rowcolor{popblueL} \textbf{Acteur} & \textbf{Flux financier (sortie)} & \textbf{Flux financier (entrée)} \\ \hline
\textbf{États-Unis} & Prêts privés (obligations) vers l'Allemagne, l'Autriche, la Hongrie, la Pologne, etc. & Remboursement des dettes de guerre (France, UK) + Intérêts des prêts allemands \\ \hline
\textbf{Allemagne} & Paiement des réparations (Traité de Versailles, Plan Dawes 1924, Plan Young 1929) & Prêts américains (obligations Dawes/Young) \\ \hline
\textbf{France / Royaume-Uni} & Remboursement des dettes de guerre aux USA & Réparations reçues de l'Allemagne \\ \hline
\end{tabularx}
\end{center}

Ce « cycle de la dette » (ou « carrousel financier ») ne fonctionne que tant que les capitaux américains affluent vers l'Europe. Or, à partir de 1928, la hausse des taux d'intérêt aux USA (pour freiner la spéculation) et l'attrait de la Bourse détournent l'épargne américaine des prêts à l'étranger. L'Allemagne, privée de nouveaux crédits, ne peut plus payer ses réparations ; la France et le Royaume-Uni ne peuvent plus rembourser leurs dettes de guerre aux USA. Le système se grippe bien avant le krach boursier d'octobre 1929.

\begin{methode}[Lire et interpréter une courbe boursière historique]
\textbf{Objectif :} Repérer les phases de tendance, les ruptures et les seuils psychologiques pour comprendre la dynamique spéculative.
\begin{enumerate}
    \item \textbf{Identifier l'échelle et l'unité :} Axe des ordonnées (indice, points, ou base 100), axe des abscisses (temps : mois, années). Vérifier si l'échelle est linéaire ou logarithmique (une chute de 50 \% a la même amplitude visuelle en log).
    \item \textbf{Repérer la tendance de fond (long terme) :} Haussière (creux et sommets ascendants), baissière (creux et sommets descendants), ou latérale (range). Tracer une droite de tendance si pertinent.
    \item \textbf{Localiser les ruptures brutales (krachs, pics) :} Variations verticales soudaines sur quelques séances (ex. : -12 \% le 28 oct. 1929, -11 \% le 29 oct.). Noter les dates exactes.
    \item \textbf{Analyser les volumes (si disponibles) :} Un pic de volume accompagne souvent un sommet de marché (distribution par les « initiés ») ou un creux de panique (capitulation).
    \item \textbf{Contextualiser :} Relier les points de retournement aux événements politiques, économiques, monétaires (hausse des taux FED, faillites bancaires, lois protectionnistes).
    \item \textbf{Conclure sur la nature du mouvement :} Bulle spéculative (hausse exponentielle, déconnexion des fondamentaux, euphorie médiatique) vs tendance fondamentale soutenue par les bénéfices des entreprises.
\end{enumerate}
\end{methode}

\courssub{Exercice résolu : Analyse d'un document statistique sur la production automobile}

\begin{exoresolu}[Analyse de la production automobile américaine (1921-1929)]
\textbf{Document :} Tableau de la production annuelle d'automobiles aux États-Unis (en millions d'unités).
\begin{center}
\begin{tabular}{|c|c|c|c|c|c|c|c|c|c|}\hline
Année & 1921 & 1922 & 1923 & 1924 & 1925 & 1926 & 1927 & 1928 & 1929 \\ \hline
Production (M) & 1,9 & 2,6 & 3,7 & 3,6 & 4,3 & 4,3 & 3,5 & 4,4 & 5,3 \\ \hline
\end{tabular}
\end{center}

\textbf{Questions :}
\begin{enumerate}
    \item Représenter graphiquement l'évolution de la production (choisir le graphique adapté).
    \item Calculer le taux de variation moyen annuel de la production sur la période 1921-1929.
    \item Identifier l'année de rupture de la croissance continue et proposer une explication économique.
    \item En déduire si le marché automobile était saturé en 1929.
\end{enumerate}

\tcblower

\textbf{Solution détaillée :}

\begin{enumerate}
    \item \textbf{Graphique :} Un diagramme en barres (histogramme) ou une courbe en points reliés est adapté pour une série chronologique annuelle discrète. L'axe des abscisses porte les années, l'axe des ordonnées la production en millions. Le titre : « Production automobile aux États-Unis (1921-1929) ».
    \item \textbf{Taux de variation moyen annuel (TVMA) :}
    Formule : $TVMA = \left( \frac{V_{final}}{V_{initial}} \right)^{\frac{1}{n}} - 1$ avec $n$ = nombre d'années.
    $V_{1921} = 1{,}9$ ; $V_{1929} = 5{,}3$ ; $n = 8$.
    $TVMA = \left( \frac{5{,}3}{1{,}9} \right)^{\frac{1}{8}} - 1 = (2{,}789)^{0,125} - 1 \approx 1{,}136 - 1 = 0{,}136$.
    \textbf{La production a cru en moyenne de 13,6 \% par an entre 1921 et 1929.}
    \item \textbf{Rupture :} La production baisse en 1927 (3,5 M contre 4,3 M en 1926), soit une chute de -18,6 \%. C'est l'année du changement de modèle chez Ford (arrêt de la Ford T en mai 1927 pour préparer la Ford A, lancée en décembre). Cette transition industrielle provoque un arrêt temporaire de la production chez le leader du marché (50 \% des parts). Hors cet aléa technique, la tendance reste haussière, mais le taux de croissance ralentit : +23 \% (1928/1927) puis +20 \% (1929/1928) contre +40-50 \% au début des années 20.
    \item \textbf{Saturation ?} En 1929, le parc atteint 23 millions de véhicules pour 122 millions d'habitants (1 voiture pour 5,3 habitants). Le taux d'équipement des ménages urbains est élevé. Le marché de \emph{premier équipement} montre des signes de saturation (ralentissement de la croissance), mais le marché de \emph{remplacement} (seconde voiture, renouvellement) n'est pas encore mature. La surproduction latente existe : les usines tournent à 80-85 \% de leurs capacités maximales, les stocks de voitures invendues s'accumulent chez les concessionnaires fin 1929. \textbf{Conclusion :} Le marché n'est pas saturé en valeur absolue, mais la demande solvable (sans crédit excessif) ne suit plus l'offre potentielle.
\end{enumerate}
\end{exoresolu}

\begin{aretenir}[Synthèse : Une prospérité bâtie sur le sable]
\begin{itemize}
    \item \textbf{Moteurs réels :} Innovations technologiques (électricité, chimie), organisation du travail (fordisme), marché intérieur vaste, position de créancier mondial.
    \item \textbf{Fragilités majeures :} Surproduction agricole structurelle ; surproduction industrielle naissante (productivité > salaires) ; inégalités de revenus record ; endettement massif des ménages (crédit conso) et des fermiers ; spéculation boursière déconnectée de l'économie réelle (effet de levier 10 \%) ; système financier international dépendant du flux de capitaux US.
    \item \textbf{Le point de bascule :} Dès 1928, la FED (Réserve fédérale) relève ses taux pour calmer la spéculation $\rightarrow$ crédit plus cher $\rightarrow$ ralentissement investissement/consommation $\rightarrow$ afflux capitaux US vers Wall Street $\rightarrow$ assèchement prêts à l'Europe $\rightarrow$ fragilisation Allemagne/France/UK. L'économie réelle marque le pas (production acier, auto, construction en baisse dès mi-1929) alors que la Bourse bat encore des records. Le krach d'octobre 1929 n'est que le révélateur brutal de cette divergence.
\end{itemize}
\end{aretenir}

\coursec{Le krach de 1929 et la grande crise économique mondiale}

La « prospérité » des années 1920, que nous avons analysée dans la section précédente, reposait sur des fondations fragiles : surproduction agricole et industrielle, endettement massif des ménages et des entreprises, spéculation boursière déconnectée de l'économie réelle. Cette bulle spéculative ne pouvait qu'éclater, plongeant le monde dans la plus grave crise économique de l'histoire contemporaine.

\courssub{Le « Jeudi noir » : l'effondrement de la Bourse de New York}

Le \textbf{24 octobre 1929}, « jeudi noir », marque le point de rupture. À la Bourse de New York (\emph{Wall Street}), la panique s'empare des investisseurs : près de \num{12,9} millions d'actions sont mises en vente en une seule séance, sans trouver suffisamment d'acheteurs. Les cours s'effondrent. Le mardi 29 octobre (« mardi noir »), environ \num{16,4} millions de titres changent de mains ; en quelques semaines, la capitalisation boursière a fondu de plusieurs dizaines de milliards de dollars.

\begin{definition}[Krach boursier]
Effondrement brutal et généralisé des cours des actions en Bourse, provoqué par une vente massive et panique des titres, entraînant la ruine des spéculateurs et la paralysie du marché financier.
\end{definition}

\begin{attention}[Krach et crise : ne pas confondre cause et déclencheur]
Le krach d'octobre 1929 n'est pas la \emph{cause profonde} de la crise : les causes sont la surproduction, le sous-emploi des capacités, l'endettement et la spéculation à crédit. Le krach est le \emph{déclencheur} et le \emph{révélateur} de la crise. Il met brutalement fin à l'illusion d'une croissance sans fin et détruit la richesse « papier » qui soutenait la consommation et l'investissement. Au Probatoire, cette distinction est attendue dans l'introduction d'un devoir.
\end{attention}

\courssub{La chaîne de la crise : de la finance à l'économie réelle}

Le krach enclenche un cercle vicieux dévastateur. La perte de valeur des actions ruine les spéculateurs qui avaient acheté à crédit ; incapables de rembourser, ils entraînent les banques dans leur chute. Les banques font faillite, l'épargne des ménages disparaît, le crédit s'assèche. Les entreprises, privées de financement et confrontées à une demande en chute libre, ferment. Le chômage explose, réduisant encore la consommation, ce qui entraîne de nouvelles faillites : c'est une \cle{boucle de rétroaction négative}.

\begin{popfigure}
\begin{tikzpicture}[
    node distance=1.1cm and 1.5cm,
    box/.style={rectangle, draw=popdark, thick, fill=popcream, rounded corners, minimum width=3.4cm, minimum height=1.1cm, align=center, font=\small},
    arrow/.style={-Stealth, thick, popdark},
    link/.style={font=\footnotesize, text=popdark, midway, fill=popcream, inner sep=2pt}
]
\node[box, align=center] (krach){Krach boursier\\ (octobre 1929)};
\node[box, below=of krach, align=center] (banques){Faillites bancaires\\ et assèchement du crédit};
\node[box, below=of banques, align=center] (entreprises){Faillites d'entreprises\\ et arrêt de la production};
\node[box, below=of entreprises, align=center] (chomage){Chômage de masse\\ (13 millions aux USA en 1933)};
\node[box, below=of chomage, align=center] (conso){Effondrement de la\\ consommation des ménages};
\node[box, below=of conso, align=center] (nouvelles){Nouvelles faillites\\ et aggravation de la crise};
\draw[arrow] (krach) -- node[link, left, align=center]{Perte de confiance,\\ retraits massifs} (banques);
\draw[arrow] (banques) -- node[link, left, align=center]{Crédit coupé,\\ appels de fonds} (entreprises);
\draw[arrow] (entreprises) -- node[link, left, align=center]{Licenciements,\\ fermetures} (chomage);
\draw[arrow] (chomage) -- node[link, left, align=center]{Perte de revenus,\\ pouvoir d'achat nul} (conso);
\draw[arrow] (conso) -- node[link, left, align=center]{Demande nulle,\\ stocks invendus} (nouvelles);
\draw[arrow, bend right=40] (nouvelles.east) to node[link, right=1.5cm, align=left] {Boucle de\\ rétroaction\\ négative} (banques.east);
\end{tikzpicture}
\legende{Schéma de la chaîne de causalité : du krach boursier à la dépression économique par boucles de rétroaction.}
\end{popfigure}

\begin{aretenir}[Chiffres clés de la Grande Dépression (1929--1933)]
\begin{itemize}
    \item \textbf{Production industrielle mondiale} : chute d'environ \textbf{un tiers} (environ \SI{-30}{\%}).
    \item \textbf{Commerce mondial} (en valeur) : chute de près de \textbf{deux tiers} (environ \SI{-65}{\%}).
    \item \textbf{Chômage aux États-Unis} : environ \num{1,5} million (1929) $\to$ près de \num{13} millions (1933), soit environ \SI{25}{\%} de la population active.
    \item \textbf{Chômage en Allemagne} : environ \num{1,3} million (1929) $\to$ environ \num{6} millions (1932).
    \item \textbf{Prix des matières premières} : cacao, café, coton perdent \SI{50}{\%} à \SI{70}{\%} de leur valeur.
\end{itemize}
\end{aretenir}

\courssub{Une crise mondiale : propagation et effondrement du commerce}

La crise ne reste pas cantonnée aux États-Unis. Elle se propage au monde entier par trois canaux majeurs :
\begin{enumerate}
    \item \textbf{Le canal financier} : les banques américaines rappellent leurs prêts à court terme accordés à l'Europe (notamment à l'Allemagne, qui dépendait des capitaux américains depuis le plan Dawes de 1924 puis le plan Young de 1929). Privée de liquidités, l'Europe centrale voit son système bancaire s'effondrer : faillite du \emph{Creditanstalt} autrichien (mai 1931), puis crise bancaire allemande (juillet 1931).
    \item \textbf{Le canal commercial} : les États-Unis relèvent leurs droits de douane (\emph{Hawley-Smoot Tariff Act}, 1930). Les pays européens ripostent par des mesures protectionnistes. Le protectionnisme généralisé asphyxie les échanges internationaux.
    \item \textbf{Le canal monétaire} : l'abandon de l'étalon-or par le Royaume-Uni (septembre 1931), puis la dévaluation du dollar (1933), créent une instabilité monétaire qui paralyse le commerce international.
\end{enumerate}

\begin{popfigure}
\begin{tikzpicture}
\begin{axis}[
    width=\linewidth, height=8cm,
    xlabel={Année}, ylabel={Indice (1929 = 100)},
    xmin=1928.5, xmax=1933.5, ymin=20, ymax=110,
    xtick={1929,1930,1931,1932,1933},
    ytick={20,40,60,80,100},
    grid=major, grid style={popline},
    legend style={at={(0.02,0.98)}, anchor=north west, fill=popcream, draw=popdark}
]
\addplot[popblue, very thick, mark=*, mark size=2.5] coordinates {
    (1929,100) (1930,78) (1931,59) (1932,38) (1933,35)
};
\addlegendentry{Valeur du commerce mondial}
\addplot[poporange, very thick, mark=square*, mark size=2.5] coordinates {
    (1929,100) (1930,85) (1931,65) (1932,42) (1933,38)
};
\addlegendentry{Prix des matières premières (indicatif)}
\end{axis}
\end{tikzpicture}
\legende{Effondrement de la valeur du commerce mondial et des prix des matières premières (base 100 en 1929). Source : données historiques de la Société des Nations, valeurs arrondies.}
\end{popfigure}

\begin{popfigure}
\begin{tikzpicture}
\begin{axis}[
    width=\linewidth, height=8cm,
    ybar, bar width=6mm,
    xlabel={Année}, ylabel={Chômeurs (millions)},
    symbolic x coords={1929,1930,1931,1932,1933},
    xtick=data,
    ymin=0, ymax=15,
    ytick={0,2,4,6,8,10,12,14},
    ymajorgrids=true, grid style={popline},
    legend style={at={(0.5,1.02)}, anchor=south, legend columns=2, fill=popcream, draw=popdark},
    nodes near coords, nodes near coords align={vertical}, every node near coord/.append style={font=\tiny, popdark}
]
\addplot[fill=popblue!70, draw=popblue] coordinates {
    (1929,1.5) (1930,4.3) (1931,8.0) (1932,12.1) (1933,12.8)
};
\addplot[fill=poporange!70, draw=poporange] coordinates {
    (1929,1.3) (1930,3.0) (1931,4.5) (1932,6.0) (1933,5.5)
};
\legend{États-Unis, Allemagne}
\end{axis}
\end{tikzpicture}
\legende{Évolution du nombre de chômeurs aux États-Unis et en Allemagne (1929--1933), valeurs arrondies. En Allemagne, le chômage officiel recule à partir de 1933 avec les grands travaux et le réarmement engagés par le régime nazi.}
\end{popfigure}

\begin{attention}[Erreur fréquente : une crise « américaine » ?]
Ne pas limiter la crise de 1929 aux États-Unis. C'est une \textbf{crise mondiale}. La structure de l'économie mondiale de l'époque (dépendance européenne aux capitaux américains, étalon-or, spécialisation coloniale dans les matières premières) en fait le vecteur de transmission. Les colonies africaines, productrices de matières premières, sont touchées de plein fouet par l'effondrement des cours et la contraction de la demande des métropoles.
\end{attention}

\courssub{Le Cameroun face à la crise : une colonie sous tension}

Le Cameroun, sous mandat français (pour l'essentiel du territoire) et britannique (deux bandes occidentales), est une économie d'exportation : cacao, café, coton, bois, huile de palme, caoutchouc. La chute des cours mondiaux frappe durement les revenus d'exportation et, par ricochet, le budget colonial et les planteurs africains.

\begin{itemize}
    \item \textbf{Chute des cours} : le prix du cacao payé au planteur s'effondre entre 1929 et 1932, perdant souvent plus de la moitié de sa valeur. Le café suit la même trajectoire.
    \item \textbf{Baisse des revenus des planteurs} : les paysans africains, qui avaient investi dans les cultures d'exportation, voient leur pouvoir d'achat s'effondrer. Beaucoup réduisent les plantations ou se replient sur les cultures vivrières.
    \item \textbf{Hausse de la pression fiscale} : pour compenser la baisse des recettes douanières, l'administration coloniale augmente l'impôt de capitation (impôt par tête) et en durcit le recouvrement.
    \item \textbf{Recours au travail forcé} : face à la pénurie de main-d'œuvre volontaire (due à la baisse des salaires et au repli vers les cultures vivrières), l'administration intensifie le recours aux \emph{prestations} (travail forcé légal pour les travaux dits d'intérêt public) et au travail forcé sur les chantiers, notamment ferroviaires (poursuite du Transcamerounais).
\end{itemize}

\begin{savaistu}[Le Cameroun, économie de plantation en crise]
Au début des années 1930, le prix du cacao payé au planteur camerounais chute si bas qu'il ne couvre plus les frais de production. De nombreuses plantations sont délaissées au profit des cultures vivrières (igname, manioc, banane plantain). Ce repli oblige l'administration à réagir : création de l'Office du cacao (1931) pour encadrer et soutenir la filière, mais aussi renforcement du contrôle social et fiscal pour maintenir la production d'exportation. La crise révèle ainsi la dépendance de l'économie camerounaise aux cours mondiaux des matières premières.
\end{savaistu}

\courssub{Les réponses politiques : New Deal et montée des régimes autoritaires}

Face à l'ampleur du désastre, les réponses divergent selon les régimes politiques.

\textbf{1. Aux États-Unis : le New Deal (1933).}\par
Élu en novembre 1932, Franklin D. Roosevelt lance dès mars 1933 une politique de relance massive : intervention de l'État dans l'économie, assainissement et régulation bancaires (\emph{Emergency Banking Act}), grands travaux publics (\emph{Tennessee Valley Authority}, \emph{Public Works Administration}), protection sociale (\emph{Social Security Act}, 1935), soutien aux prix agricoles. C'est la naissance d'une forme d'État interventionniste, inspirée plus tard par les théories de l'économiste John Maynard Keynes.

\begin{definition}[New Deal]
Politique de relance économique et de réformes structurelles menée par le président Franklin D. Roosevelt à partir de 1933 aux États-Unis. Elle repose sur l'intervention de l'État (grands travaux, régulation bancaire, protection sociale, soutien aux prix agricoles) pour résorber le chômage et relancer la demande globale.
\end{definition}

\textbf{2. En Europe : la tentation autoritaire.}\par
L'incapacité des régimes parlementaires à résoudre la crise (politiques de déflation et d'austérité budgétaire) discrédite la démocratie libérale.
\begin{itemize}
    \item \textbf{Allemagne} : la crise est le terreau de l'arrivée d'Adolf Hitler au pouvoir \textbf{légalement} le 30 janvier 1933. Le chômage de masse (environ \num{6} millions de chômeurs) et la ruine des classes moyennes détruisent la République de Weimar. Hitler promet « du pain et du travail » grâce aux grands travaux (autoroutes) et au réarmement.
    \item \textbf{Autres pays} : renforcement du régime fasciste en Italie (Mussolini au pouvoir depuis 1922), guerre d'Espagne (1936--1939), régimes autoritaires en Europe centrale et orientale, militarisation du Japon.
\end{itemize}

\begin{aretenir}[Lien entre la crise de 1929 et la Seconde Guerre mondiale]
La crise de 1929 est un \textbf{facteur majeur} de la Seconde Guerre mondiale :
\begin{itemize}
    \item elle détruit la stabilité économique de l'entre-deux-guerres ;
    \item elle favorise l'arrivée au pouvoir de régimes révisionnistes et expansionnistes (Allemagne nazie, Japon militariste) ;
    \item elle pousse les puissances démocratiques (France, Royaume-Uni) à la politique d'apaisement (accords de Munich, 1938), par faiblesse économique et pacifisme de l'opinion ;
    \item elle justifie l'autarcie et la conquête d'« espaces vitaux » (\emph{Lebensraum} pour l'Allemagne, « sphère de coprospérité » pour le Japon).
\end{itemize}
\end{aretenir}

\courssub{Travaux pratiques : analyse de documents et croquis}

\begin{exoresolu}[Analyse de l'impact de la crise sur le commerce mondial]
\textbf{Document :} d'après la Société des Nations, la valeur du commerce mondial est passée d'environ 30 milliards de dollars-or en 1929 à environ 10,5 milliards en 1933, tandis que le volume des échanges reculait d'environ \SI{25}{\%}.

\textbf{Questions :}
\begin{enumerate}
    \item Calculer le taux de variation de la valeur du commerce mondial entre 1929 et 1933.
    \item Expliquer pourquoi la valeur chute plus vite que le volume (\SI{-65}{\%} contre \SI{-25}{\%}).
    \item Citer deux conséquences de cet effondrement sur les économies coloniales.
\end{enumerate}

\tcblower

\textbf{Solution détaillée :}
\begin{enumerate}
    \item Taux de variation $= \dfrac{V_{1933} - V_{1929}}{V_{1929}} \times 100 = \dfrac{10,5 - 30}{30} \times 100 = \dfrac{-19,5}{30} \times 100 = \mathbf{-65\,\%}$.
    \item La valeur chute plus que le volume car les \textbf{prix} se sont effondrés (déflation généralisée) : les prix des produits échangés ont baissé d'environ \SI{50}{\%}. On échange donc un peu moins de marchandises (volume : \SI{-25}{\%}), mais surtout des marchandises beaucoup moins chères, d'où une chute cumulée de la valeur d'environ \SI{-65}{\%}.
    \item Conséquences coloniales : (1) effondrement des recettes d'exportation $\to$ déficits budgétaires coloniaux $\to$ hausse des impôts locaux (capitation) ; (2) chute des revenus des producteurs africains $\to$ insécurité alimentaire, abandon des cultures d'exportation, repli vers les cultures vivrières, et durcissement du travail forcé pour maintenir les infrastructures d'évacuation (routes, voies ferrées).
\end{enumerate}
\end{exoresolu}

\begin{methode}[Rédiger une conclusion argumentée (Probatoire / Baccalauréat technique)]
Pour conclure un devoir sur la crise de 1929 :
\begin{enumerate}
    \item \textbf{Rappeler la problématique} : ex. « En quoi la crise de 1929 marque-t-elle une rupture dans l'histoire économique et politique du XX\textsuperscript{e} siècle ? »
    \item \textbf{Mobiliser 2 à 3 faits précis et datés} : ex. « Le krach d'octobre 1929 déclenche une dépression mondiale (environ \num{13} millions de chômeurs aux USA, \num{6} millions en Allemagne). L'effondrement du commerce mondial (\SI{-65}{\%} en valeur) ruine les économies coloniales, comme celle du cacao camerounais. »
    \item \textbf{Synthétiser les conséquences structurelles} : remise en cause du libéralisme économique pur, naissance de l'interventionnisme d'État (New Deal), montée des régimes totalitaires (Hitler au pouvoir en 1933).
    \item \textbf{Ouvrir sur une conséquence durable} : « Cette crise structure le monde d'après 1945 : accords de Bretton Woods (1944), GATT, construction européenne et États-providence, pour éviter la répétition de 1929. »
\end{enumerate}
\end{methode}

\begin{exercice}[Entraînement : croquis de la propagation de la crise]
Réaliser un croquis schématique (format paysage) intitulé « La propagation mondiale de la crise de 1929 (1929--1933) ».
\begin{itemize}
    \item Représenter les 3 pôles : États-Unis, Europe (Allemagne, Royaume-Uni, France), Colonies (Afrique, Asie, Amérique latine).
    \item Utiliser des flèches de couleurs et d'épaisseurs différentes pour les 3 canaux : \textbf{financier} (retrait des capitaux américains vers l'Europe), \textbf{commercial} (protectionnisme, baisse des échanges), \textbf{monétaire} (abandon de l'étalon-or, dévaluations).
    \item Légender les impacts locaux : chômage (USA, Allemagne), chute des cours des matières premières (Afrique, Amérique latine), faillites bancaires (Europe centrale).
    \item Ajouter une flèche « retour » : Colonies $\to$ Métropoles (baisse des commandes, troubles sociaux).
\end{itemize}
\end{exercice}

\begin{corrige}[Exercice : croquis de la propagation]
Le croquis doit montrer :
\begin{itemize}
    \item \textbf{Épicentre américain} : krach $\to$ rappel des prêts $\to$ Hawley-Smoot (1930).
    \item \textbf{Flèche financière (rouge, épaisse)} : USA $\to$ Allemagne (capitaux du plan Young stoppés) $\to$ faillites bancaires en Autriche puis en Allemagne (1931) $\to$ contagion au Royaume-Uni et à la France.
    \item \textbf{Flèche commerciale (bleue)} : protectionnisme généralisé $\to$ effondrement du volume et de la valeur des échanges mondiaux.
    \item \textbf{Flèche monétaire (verte)} : le Royaume-Uni quitte l'étalon-or (septembre 1931) $\to$ dévaluation du dollar (1933) $\to$ instabilité monétaire persistante.
    \item \textbf{Impact sur les colonies (flèches pointillées vers l'Afrique, l'Asie, l'Amérique latine)} : chute des cours (cacao, café, coton, caoutchouc) $\to$ crise budgétaire coloniale $\to$ pression fiscale et travail forcé.
    \item \textbf{Retour (flèche orange)} : baisse des importations des colonies $\to$ baisse des exportations des métropoles $\to$ aggravation du chômage en Europe.
\end{itemize}
La légende doit être organisée par type de flux. Une rose des vents et une échelle qualitative (épaisseur proportionnelle à l'intensité) sont attendues.
\end{corrige}

\coursec{Travaux pratiques : cartes, croquis et analyse de données}

Après avoir étudié le krach de Wall Street et la grande crise économique mondiale, il est temps de passer du récit à la pratique. L'historien ne se contente pas de lire : il \cle{construit} des cartes, il \cle{analyse} des documents, il \cle{calcule} à partir de données chiffrées. Cette séance de travaux pratiques vous entraîne à ces trois gestes professionnels, en restant dans les deux savoirs de la leçon : la Première Guerre mondiale et la crise de 1929. Chaque TP donne lieu à une production évaluée : carte, croquis, courbe, commentaire ou conclusion rédigée.

\courssub{TP 1 — Construire une carte : les deux blocs d'alliances en Europe en 1914}

\textbf{Objectif.} Réaliser sur fond de carte muet de l'Europe une carte intitulée « Les deux blocs d'alliances en Europe en 1914 ».

\textbf{Consignes.} Coloriez en deux couleurs distinctes les pays de la \cle{Triple-Alliance} (Allemagne, Autriche-Hongrie, Italie) et ceux de la \cle{Triple-Entente} (France, Royaume-Uni, Russie). Laissez les pays neutres (Espagne, Suisse, Pays-Bas, Norvège, Suède) en gris. Placez les capitales par un point nommé.

\begin{methode}[Construire un croquis ou une carte géographique]
\begin{enumerate}
\item \textbf{Tracer le fond de carte} : contour simplifié du continent ou du pays, frontières principales. Inutile d'être exact au millimètre : le croquis doit être reconnaissable.
\item \textbf{Choisir les éléments essentiels seulement} : chaque élément doit servir le sujet. Ici : les deux alliances, rien d'autre.
\item \textbf{Organiser la légende} : couleurs, symboles et traits classés par thème, chacun avec sa signification. Tout ce qui est sur la carte doit figurer dans la légende, et rien dans la légende ne doit manquer sur la carte.
\item \textbf{Compléter le cartouche} : titre précis (thème + lieu + date), orientation (flèche du nord), échelle indicative.
\end{enumerate}
\end{methode}

\begin{attention}[L'erreur classique du croquis surchargé]
La faute la plus fréquente est de vouloir tout représenter : relief, rivières, climats, villes secondaires\ldots{} Un bon croquis est \cle{sélectif} et \cle{lisible} : trois à cinq éléments bien choisis valent mieux que quinze éléments illisibles. Le correcteur juge la pertinence du choix, pas la quantité.
\end{attention}

\courssub{TP 2 — Croquis de la campagne du Cameroun (1914-1916)}

\textbf{Objectif.} Sur le fond de carte du Cameroun ci-dessous, construire le croquis « La conquête du Cameroun allemand, 1914-1916 ».

\textbf{Consignes.} Placez les villes de \cle{Douala}, \cle{Yaoundé} et \cle{Garoua}. Tracez par des flèches pleines les axes d'invasion : les troupes britanniques venues du Nigeria vers le nord (Garoua, Mora) et l'ouest ; les troupes françaises remontant du Congo (Brazzaville) vers le sud et le centre (Yaoundé). Indiquez par des flèches pointillées les replis allemands : la garnison du nord vers le camp retranché de Mora, la colonne principale (gouverneur Ebermaier) depuis Yaoundé vers la frontière orientale (Rio Muni, Guinée espagnole). Indiquez la reddition de Mora (février 1916). Complétez la légende organisée par thèmes (invasion alliée, repli allemand, lieux clés) et le cartouche (titre, nord, échelle).

\begin{popfigure}
\begin{center}
\begin{tikzpicture}[scale=1.05]
% Contour très simplifié du Cameroun
\draw[thick,popdark,fill=popcream] (0,0) -- (2.2,-0.2) -- (4.6,0.3) -- (5.2,1.6) -- (4.8,3.2) -- (4.2,4.6) -- (3.4,5.6) -- (2.6,6.4) -- (2.2,7.4) -- (1.4,7.8) -- (0.9,6.8) -- (1.2,5.4) -- (0.4,4.2) -- (-0.6,3.0) -- (-0.4,1.4) -- cycle;
% Villes repères
\fill[popink] (0.6,1.2) circle (0.07) node[right,font=\small]{Douala};
\fill[popink] (2.0,2.4) circle (0.07) node[right,font=\small]{Yaoundé};
\fill[popink] (2.3,5.9) circle (0.07) node[right,font=\small]{Garoua};
% Nord
\draw[->,thick,popdark] (5.6,6.6) -- (5.6,7.6) node[above,font=\small]{N};
% Échelle indicative
\draw[thick,popdark] (4.6,-0.6) -- (5.6,-0.6) node[midway,below,font=\small]{200 km (indicatif)};
% Indication pays limitrophes (texte seul)
\node[font=\tiny,popmuted,align=center] at (-1.5,4.0) {Nigeria\\(R.-U.)};
\node[font=\tiny,popmuted,align=center] at (6.5,2.5) {Congo\\(France)};
\node[font=\tiny,popmuted,align=center] at (3.5,8.0) {Tchad\\(France)};
\node[font=\tiny,popmuted,align=center] at (-0.5,-0.5) {Guinée esp.};
\end{tikzpicture}
\end{center}
\legende{Fond de carte simplifié du Cameroun (1914) avec villes repères, orientation et échelle. L'élève ajoute les flèches d'invasion (pleines), les replis allemands (pointillées), les symboles de batailles/redditions et rédige la légende complète.}
\end{popfigure}

\begin{exemplebox}[Un titre précis vaut des points]
Mauvais titre : « Le Cameroun ». Bon titre : « La campagne du Cameroun et la défaite allemande (1914-1916) ». Le titre doit contenir le \textbf{thème}, le \textbf{lieu} et la \textbf{date}.
\end{exemplebox}

\courssub{TP 3 — Tracer et commenter la courbe de l'indice de Wall Street (1921-1933)}

\textbf{Objectif.} À partir du tableau de données ci-dessous, tracer la courbe de l'indice boursier américain (base 100 en 1921), puis la commenter en trois phrases : tendance, rupture, conséquence.

\begin{center}
\begin{tabularx}{\linewidth}{|l|X|X|X|X|X|X|X|X|}\hline
\rowcolor{popblueL} \textbf{Année} & 1921 & 1923 & 1925 & 1927 & 1929 & 1930 & 1932 & 1933 \\ \hline
\textbf{Indice} & 100 & 130 & 180 & 245 & 380 & 295 & 85 & 100 \\ \hline
\end{tabularx}
\end{center}

\begin{popfigure}
\begin{center}
\begin{tikzpicture}
\begin{axis}[
width=0.85\linewidth,height=6.5cm,
xlabel={Ann\'ee},ylabel={Indice (base 100 en 1921)},
xmin=1920.5,xmax=1933.5,ymin=0,ymax=420,
xtick={1921,1923,1925,1927,1929,1931,1933},
ytick={0,100,200,300,400},
grid=both,grid style={popline!40},
legend style={at={(0.02,0.97)},anchor=north west,font=\small}
]
\addplot[popcurve,mark=*,mark size=1.5pt] coordinates {
(1921,100) (1923,130) (1925,180) (1927,245) (1929,380) (1930,295) (1932,85) (1933,100)
};
\addlegendentry{Indice de Wall Street}
\draw[dashed,poporange] (axis cs:1929,0) -- (axis cs:1929,380) node[pos=0.85,right,font=\small]{krach d'octobre 1929};
\end{axis}
\end{tikzpicture}
\end{center}
\legende{Gabarit de courbe (corrigé) : l'élève reporte les 8 points du tableau, les relie, puis annote la rupture de 1929. Source : données reconstituées, indice base 100 en 1921.}
\end{popfigure}

\begin{exoresolu}[Calculer le pourcentage de baisse entre 1929 et 1932]
Énoncé : l'indice passe de 380 en 1929 à 85 en 1932. Calculer le pourcentage de baisse, puis rédiger une phrase de conclusion.
\tcblower
Solution détaillée. Le taux de variation se calcule par :
\[
t = \frac{\text{valeur finale} - \text{valeur initiale}}{\text{valeur initiale}} \times 100
= \frac{85 - 380}{380} \times 100
= \frac{-295}{380} \times 100 \approx -77{,}6\,\%.
\]
L'indice a donc perdu environ \textbf{78\,\%} de sa valeur entre 1929 et 1932. Conclusion rédigée : « En trois ans, la Bourse de New York a perdu près des quatre cinquièmes de sa valeur : le krach de 1929 n'est pas une simple baisse passagère, mais un effondrement durable qui ruine les épargnants et précipite les banques et les entreprises en faillite. »
\end{exoresolu}

\begin{attention}[Piège de rédaction]
Ne confondez pas « baisse de 295 points » et « baisse de 77,6\,\% » : le correcteur attend le calcul en pourcentage, seul comparable. Autre piège : écrire « l'indice remonte à 100 en 1933, donc la crise est finie » — non : 100, c'est le niveau de 1921, soit douze ans d'expansion effacés.
\end{attention}

\courssub{TP 4 — Dossier documentaire : tranchées, propagande et témoignage camerounais}

Le dossier comporte trois documents : une photographie de tranchée sur le front occidental (1916), une affiche de propagande française appelant à l'effort de guerre, et le témoignage d'un ancien combattant camerounais racontant son engagement dans les tirailleurs et son retour au pays.

\begin{methode}[Analyser une photographie historique]
\begin{enumerate}
\item \textbf{Décrire} objectivement : que voit-on ? (personnages, objets, décor, premier plan / arrière-plan).
\item \textbf{Identifier le contexte} : date, lieu, auteur si connu, événement auquel l'image se rattache.
\item \textbf{Dégager l'intention de l'auteur} : informer, dénoncer, glorifier, mobiliser ? Une photographie de guerre n'est jamais neutre : elle sert un message.
\item \textbf{Conclure} sur l'intérêt et les limites du document pour l'historien (point de vue partiel, censure, mise en scène possible).
\end{enumerate}
\end{methode}

\textbf{Questions guidées.} Pour la photographie : quelles conditions de vie des soldats devine-t-on (boue, promiscuité, attente) ? Pour l'affiche : quels mots et quelles images cherchent à convaincre, et qui est visé ? Pour le témoignage : que révèle-t-il de la participation des Africains à une guerre européenne, et de la façon dont ils furent traités au retour ?

\begin{savaistu}[Les monuments aux morts du Cameroun]
À Douala et à Yaoundé, les monuments aux morts rappellent la participation des Camerounais aux deux guerres mondiales : des milliers de tirailleurs et de porteurs camerounais ont servi, souvent sous la contrainte, et beaucoup ne sont jamais revenus. Ces monuments sont eux-mêmes des documents historiques : ils montrent comment une société choisit de se souvenir.
\end{savaistu}

\courssub{TP 5 — Analyse croisée de données chiffrées : rédiger une conclusion justifiée}

\textbf{Objectif.} Croiser trois séries statistiques pour montrer que la crise de 1929 atteint jusqu'aux planteurs camerounais.

\begin{center}
\begin{tabularx}{\linewidth}{|l|X|X|X|X|}\hline
\rowcolor{popblueL} \textbf{Indicateur} & 1928 & 1930 & 1932 & Variation 1928-1932 \\ \hline
Chômage aux États-Unis (\% de la population active) & 4 & 9 & 24 & $+20$ points \\ \hline
Production industrielle mondiale (indice base 100) & 100 & 86 & 62 & $-38\,\%$ \\ \hline
Prix du cacao à l'exportation (indice base 100) & 100 & 70 & 45 & $-55\,\%$ \\ \hline
\end{tabularx}
\end{center}

\begin{exemplebox}[De la donnée à la conclusion]
Raisonnement en chaîne : le krach détruit l'emploi industriel aux États-Unis et en Europe (chômage en hausse) ; les consommateurs appauvris achètent moins de produits coloniaux ; la demande de cacao s'effondre, donc son prix chute de plus de moitié ; le planteur camerounais, payé moins cher sa récolte, voit ses revenus s'effondrer alors que l'administration maintient l'impôt. Conclusion justifiée : « La crise de 1929 est une crise \cle{mondiale} : née à Wall Street, elle frappe le paysan camerounais par le biais des cours mondiaux des matières premières. »
\end{exemplebox}

\courssub{Restitution : exposé collectif}

Par groupes de quatre, préparez un exposé de cinq minutes reliant \textbf{guerre}, \textbf{crise} et \textbf{vie quotidienne au Cameroun}. Plan conseillé : (1) rappel d'un fait de la guerre 1914-1918 touchant le Cameroun (campagne de 1914-1916, tirailleurs) ; (2) présentation de votre courbe de Wall Street et du calcul de baisse ; (3) explication du lien entre krach, chute du prix du cacao et vie des planteurs ; (4) conclusion personnelle. Chaque groupe présente au moins un support construit pendant les TP (carte, croquis ou courbe). La grille d'évaluation porte sur : exactitude des faits et dates, qualité de la légende et des calculs, clarté de l'argumentation, justification de la conclusion — deux savoir-faire exigibles au Probatoire technique.

\begin{aretenir}[L'essentiel des TP]
Une carte ou un croquis se construit en quatre gestes : fond de carte, éléments essentiels, légende organisée, titre précis. Une photographie s'analyse en trois temps : décrire, contextualiser, dégager l'intention. Une donnée chiffrée ne parle qu'à travers un calcul (taux de variation, pourcentage) et une phrase de conclusion justifiée. Ces outils transforment le récit de la guerre et de la crise en démonstration historique.
\end{aretenir}

\newpage

\coursec{Activité d'intégration}

\courssub{Objectifs de l'activité}
\begin{itemize}
    \item Mobiliser les connaissances sur la crise de 1929 et sa propagation mondiale pour expliquer une situation locale concrète : l'effondrement du revenu cacaoyer au Cameroun.
    \item Exploiter un dossier documentaire hétérogène (carte, courbes, tableau chiffré, témoignage) en croisant les informations.
    \item Calculer un taux de variation (baisse des prix) et l'interpréter historiquement.
    \item Rédiger une synthèse argumentée (10-15 lignes) structurée par une chaîne de causalité claire : Krach boursier $\rightarrow$ contraction du crédit/du commerce $\rightarrow$ effondrement de la demande mondiale $\rightarrow$ chute des cours des matières premières $\rightarrow$ ruine des producteurs coloniaux.
    \item Justifier sa démarche en citant explicitement au moins trois faits précis extraits des documents fournis.
\end{itemize}

\courssub{Situation : Enquête historique — Pourquoi le cacao de grand-père ne se vendait plus en 1932 ?}
\textbf{Contexte.} Nous sommes en 1932, dans un village de la région du Centre (zone de forêt dense), sous mandat français. Votre grand-père, \emph{Ngono}, planteur de cacao depuis 1925, possède une plantation de 5 hectares. En 1928, il vendait sa récolte à un comptoir de la \emph{Société Commerciale de l'Ouest Africain} (SCOA) à Douala et vivait décemment. En 1932, l'acheteur lui annonce que le cours a effondré : il ne peut plus lui avancer d'argent pour les outils, ni payer les taxes coloniales (impôt de capitation, prestations en nature). La famille doit réduire ses achats de tissu, de pétrole, de sel. Grand-père murmure : « C'est la faute des Blancs de New York qui ont joué à l'argent ». 

\vspace{0.5em}
\noindent\textbf{Votre mission.} En historien, vous devez valider ou infirmer l'intuition de grand-père en établissant le lien rigoureux entre le krach de la Bourse de New York (octobre 1929) et la misère du planteur camerounais en 1932. Vous disposez du \textbf{Dossier Documentaire} ci-après.

\vspace{1em}
\noindent\textbf{Document 1 : Carte schématique — Zones cacaoyères du Cameroun sous mandats (vers 1930).}
\begin{popfigure}
\begin{tikzpicture}[scale=0.9]
    % Contour très simplifié du Cameroun (polygone approximatif)
    \draw[popdark, thick] (0,0) -- (8,0) -- (9,2) -- (8.5,5) -- (6,6) -- (3,5.5) -- (1,3.5) -- (0.5,1.5) -- cycle;
    % Frontière mandats (ligne pointillée verticale approximative)
    \draw[popblue, dashed, thick] (4.5,0) -- (4.5,6);
    % Fleche Nord
    \draw[->, popdark] (7.5,5.5) -- (7.5,5.8) node[above] {N};
    % Zones cacaoyères (hachures)
    \pattern[pattern=north east lines, pattern color=poporange!60] (1,1) -- (3,0.5) -- (4,2) -- (3.5,3) -- (1.5,2.5) -- cycle; % Ouest/Sud-Ouest (Brit)
    \pattern[pattern=north west lines, pattern color=popgreen!60] (3.5,1.5) -- (6,1) -- (7,3) -- (5.5,4.5) -- (4,3.5) -- cycle; % Centre/Sud/Est (Fr)
    % Villes
    \node[circle, fill=popdark, inner sep=1.5pt] (douala) at (3.2,1.2) {};
    \node[anchor=west] at (douala.east) {\small Douala (Port)};
    \node[circle, fill=popdark, inner sep=1.5pt] (yaounde) at (4.5,3) {};
    \node[anchor=west] at (yaounde.east) {\small Yaoundé};
    \node[circle, fill=popdark, inner sep=1.5pt] (dschang) at (2.5,2.5) {};
    \node[anchor=east] at (dschang.west) {\small Dschang};
    % Légende manuelle
    \node[anchor=south west, text width=5cm, font=\small] at (0,-0.8) {
        \textbf{Légende :} 
        \tikz\fill[pattern=north east lines, pattern color=poporange!60] (0,0) rectangle (0.3,0.2); Zone cacaoyère principale (Mandat britannique - Ouest) \\
        \tikz\fill[pattern=north west lines, pattern color=popgreen!60] (0,0) rectangle (0.3,0.2); Zone cacaoyère principale (Mandat français - Centre, Sud, Est) \\
        \tikz\draw[popblue, dashed, thick] (0,0) -- (0.3,0); Frontière des mandats (Traité de Versailles, 1919) \\
        \tikz\draw[->, popdark] (0,0) -- (0.3,0); Port d'exportation (Douala)
    };
\end{tikzpicture}
\legende{Carte schématique du Cameroun vers 1930. Les zones hachurées indiquent les bassins de production cacaoyère. La frontière en pointillés sépare le mandat français (Est) du mandat britannique (Ouest). Tout le cacao sort par le port de Douala.}
\end{popfigure}

\vspace{1em}
\noindent\textbf{Document 2 : Évolution comparée de l'indice du cours du cacao et de l'indice Dow Jones (1925-1935).}
\begin{popfigure}
\begin{tikzpicture}
    \begin{axis}[
        width=\linewidth, height=8cm,
        xmin=1925, xmax=1935,
        ymin=0, ymax=220,
        axis y line*=left,
        ylabel={Indice cours cacao (1925=100)},
        ylabel style={text width=3cm, align=center},
        xtick={1925,1926,1927,1928,1929,1930,1931,1932,1933,1934,1935},
        xticklabel style={rotate=45, anchor=east, font=\tiny},
        ytick={0,50,100,150,200},
        grid=major, grid style={popline},
        legend style={at={(0.02,0.98)}, anchor=north west, font=\small, fill=popcream, draw=popdark},
        title style={font=\small, align=center},
        title={Évolution comparée : Cacao (indice) et Bourse de New York (points)}
    ]
    % Données Cacao : Indice base 100 en 1925 (valeurs arrondies)
    \addplot[poporange, thick, mark=*] coordinates {
        (1925, 100) (1926, 150) (1927, 175) (1928, 133) (1929, 108) (1930, 71) (1931, 46) (1932, 35) (1933, 40) (1934, 46) (1935, 50)
    };
    \addlegendentry{Cacao (échelle de gauche, indice 1925=100)}
    \end{axis}
    % Axe droit pour Dow Jones (points réels approximatifs)
    \begin{axis}[
        xmin=1925, xmax=1935,
        ymin=0, ymax=400,
        axis y line*=right,
        ylabel={Indice Dow Jones (points)},
        ylabel style={text width=3cm, align=center},
        ytick={0,50,100,150,200,250,300,350,400},
        hide x axis,
        legend style={at={(0.98,0.98)}, anchor=north east, font=\small, fill=popcream, draw=popdark},
    ]
    % Données Dow Jones (pic sept 1929 ~381, fond 1932 ~41)
    \addplot[popblue, thick, mark=square*] coordinates {
        (1925, 150) (1926, 180) (1927, 220) (1928, 280) (1929, 350) (1930, 200) (1931, 100) (1932, 45) (1933, 80) (1934, 100) (1935, 130)
    };
    \addlegendentry{Dow Jones (échelle de droite, points)}
    \end{axis}
\end{tikzpicture}
\legende{Évolution comparée. L'indice du cacao (orange, base 100 en 1925) chute dès 1928 (surproduction), puis s'effondre avec le krach boursier (bleu, oct. 1929). Les deux courbes atteignent leur point bas en 1932. Sources : League of Nations statistics, Historical Statistics of the USA.}
\end{popfigure}

\vspace{1em}
\noindent\textbf{Document 3 : Prix payé au planteur au comptoir de Douala (Francs par kg de cacao commercial).}
\begin{center}
\begin{tabularx}{\linewidth}{|c|X|c|}\hline
\rowcolor{popblueL}
\textbf{Année} & \textbf{Contexte économique majeur} & \textbf{Prix au planteur (F/kg)} \\ \hline
1928 & Prospérité relative, cours mondiaux hauts & 2,50 \\ \hline
1929 & Krach d'octobre (Wall Street) & 2,10 \\ \hline
1930 & Début Grande Dépression, protectionnisme (Hawley-Smoot) & 1,40 \\ \hline
1931 & Crise bancaire européenne (Kreditanstalt), abandon étalon-or (UK) & 0,95 \\ \hline
1932 & Point bas de la dépression, chômage massif USA/Allemagne & \textbf{0,60} \\ \hline
1933 & Début New Deal (USA), conférence de Londres (échec), reprise lente & 0,75 \\ \hline
\end{tabularx}
\end{center}
\legende{Prix bord champ / comptoir Douala. Source : Rapports annuels du Gouverneur du Cameroun français, Archives Nationales Yaoundé.}

\vspace{1em}
\noindent\textbf{Document 4 : Témoignage reconstitué de Ngono, planteur à Nkolbisson (région du Centre, 1932).}
\begin{definition}[Témoignage oral collecté par l'administration en 1933, rapport de circonscription]
« Mon fils, en 1928, quand j'apportais mes sacs à Douala, le commis de la SCOA me donnait 2 francs 50 le kilo. Avec 5 hectares, je faisais 3000 kg. Je touchais 7500 francs. Je payais l'impôt (12 francs/homme, nous sommes 6 hommes valides = 72 francs), j'achetais le pétrole, le tissu, le fusil de chasse, je donnais la dot pour mon fils aîné. Le chef de poste était content, le commis aussi. 

Depuis 1930, le commis fait la grimace. Il dit : "Les Blancs de l'autre côté de l'eau (l'Atlantique) ne veulent plus notre cacao, leurs usines sont fermées, leurs banques sont mortes." En 1931, il m'a donné 95 centimes. En 1932, 60 centimes ! Pour 3000 kg, cela fait 1800 francs seulement. L'impôt seul mange 72 francs. Il reste 1728 francs pour nourrir 15 bouches, payer les porteurs, acheter la houe, le machette, le quinquina contre la fièvre. 

Le commis m'a dit : "Ngono, ne plante plus, ça ne vaut plus la peine." Mais la terre, c'est notre vie. Le chef de poste m'a menacé de prison si je ne payais pas l'impôt. J'ai dû vendre ma chèvre, emprunter au traitant grec à 50\% d'intérêt. Mes fils ont arrêté l'école des Pères pour aller couper le bois chez le Blanc. C'est la faute aux Blancs de New York qui ont mangé l'argent du monde. »
\end{definition}

\courssub{Consignes}
\textbf{Travail en groupe (4 élèves, 45 min) — Puis rédaction individuelle (15 min).}

\begin{enumerate}
    \item \textbf{Analyse documentaire (15 min).} Pour chaque document, notez : 
    \begin{itemize}
        \item Nature (carte, graphique, tableau, témoignage) et date.
        \item Information principale (une phrase).
        \item Un chiffre ou fait précis marquant.
    \end{itemize}
    \item \textbf{Établissement de la chaîne de causalité (15 min).} Reliez les documents entre eux pour construire l'explication complète. Complétez le schéma ci-dessous en y insérant des flèches et des mots-clés (crédit, demande, cours mondial, prix producteur, revenu, impôt, niveau de vie) :
    \begin{center}
    \begin{tikzpicture}[node distance=1.2cm, auto, >=Stealth]
        \node[draw, rounded corners, fill=popblueL] (A) {Krach Wall Street (Oct. 1929)};
        \node[draw, rounded corners, fill=poporangeL, right=of A] (B) { ? };
        \node[draw, rounded corners, fill=popgreenL, right=of B] (C) { ? };
        \node[draw, rounded corners, fill=poppurpleL, right=of C] (D) { ? };
        \node[draw, rounded corners, fill=poppinkL, right=of D] (E) { Ruine planteur camerounais (1932) };
        \draw[->, thick] (A) -- (B);
        \draw[->, thick] (B) -- (C);
        \draw[->, thick] (C) -- (D);
        \draw[->, thick] (D) -- (E);
    \end{tikzpicture}
    \end{center}
    \item \textbf{Calcul quantitatif (5 min).} Calculez le \textbf{taux de variation en pourcentage} du prix payé au planteur entre 1928 et 1932. Formule : $\frac{V_{final} - V_{initial}}{V_{initial}} \times 100$. Interprétez le résultat en une phrase.
    \item \textbf{Rédaction individuelle (15 min).} Rédigez une réponse argumentée de \textbf{10 à 15 lignes} à la question : \emph{« Pourquoi le cacao de grand-père ne se vendait plus en 1932 ? »}. Votre texte doit :
    \begin{itemize}
        \item Présenter une thèse claire en introduction.
        \item Développer la chaîne de causalité en 3 étapes minimum (Finance $\rightarrow$ Commerce mondial $\rightarrow$ Économie coloniale).
        \item Citer explicitement \textbf{au moins 3 faits précis} (chiffres, dates, citations) tirés des \textbf{documents 1, 2, 3, 4}.
        \item Conclure sur la responsabilité du système capitaliste mondial et de la politique coloniale (impôts maintenus).
    \end{itemize}
\end{enumerate}

\courssub{Ressources à mobiliser}
\begin{methode}[Calcul du taux de variation (évolution en \%)]
    \begin{enumerate}
        \item Identifier la valeur initiale ($V_i$) et la valeur finale ($V_f$).
        \item Appliquer la formule : $T = \frac{V_f - V_i}{V_i} \times 100$.
        \item Si $T < 0$ : baisse de $|T|\%$. Si $T > 0$ : hausse de $T\%$.
        \item Toujours écrire une phrase d'interprétation : « Le prix a baissé de X\% entre [date] et [date], ce qui signifie... »
    \end{enumerate}
\end{methode}

\begin{methode}[Commentaire de document historique (méthode « 3C » : Contexte, Contenu, Conclusion)]
    \begin{enumerate}
        \item \textbf{Contexte :} Auteur, date, nature, destination, événement lié.
        \item \textbf{Contenu :} Idées principales, chiffres clés, arguments, vocabulaire spécifique.
        \item \textbf{Conclusion/Intérêt :} Ce que le document apporte à la démonstration (preuve, illustration, limite).
    \end{enumerate}
\end{methode}

\begin{aretenir}[Repères chronologiques clés de la leçon]
    \begin{itemize}
        \item \textbf{1914-1918 :} Première Guerre mondiale (guerre totale, blocus, mobilisation coloniale).
        \item \textbf{1919 :} Traité de Versailles, mandats SDN (Cameroun partagé France/GB).
        \item \textbf{1924-1929 :} « Années folles », prospérité apparente, spéculation boursière, surproduction agricole.
        \item \textbf{24-29 oct. 1929 :} Krach de New York (Jeudi noir, Mardi noir).
        \item \textbf{1930-1933 :} Grande Dépression : protectionnisme, effondrement du commerce mondial (-60\%), chômage de masse, crise bancaire.
        \item \textbf{1933 :} New Deal (Roosevelt), Hitler au pouvoir, Conférence de Londres (échec monétaire).
    \end{itemize}
\end{aretenir}

\courssub{Démarche et corrigé commenté}
\begin{exoresolu}[Enquête : Pourquoi le cacao de grand-père ne se vendait plus en 1932 ?]
\textbf{Étape 1 : Analyse documentaire (résumé attendu).}
\begin{itemize}
    \item \textbf{Doc 1 (Carte) :} Montre que le Cameroun est un territoire colonial partagé en deux mandats (Fr/GB) mais unifié par l'exportation cacaoyère vers Douala. La région du Centre (zone verte) est un bassin principal. \emph{Fait clé :} Tout le cacao passe par Douala, port unique, contrôlé par des firmes européennes (SCOA, etc.).
    \item \textbf{Doc 2 (Courbes) :} Corrélation visuelle forte. Le Dow Jones (bleu) s'effondre de 350 (1929) à 45 (1932) soit \textbf{-87\%}. L'indice du cacao (orange) chute de 108 (1929) à 35 (1932) soit \textbf{-68\%}. La chute du cacao commence légèrement avant le krach (surproduction 1928) mais s'accélère brutalement après 1929. \emph{Fait clé :} Synchronisation de l'effondrement boursier et de l'effondrement des matières premières.
    \item \textbf{Doc 3 (Tableau) :} Transmission au producteur local. Prix Douala : 2,50 F (1928) $\rightarrow$ 0,60 F (1932). Chute de \textbf{76\%} en 4 ans. Le tableau mentionne les causes macro : protectionnisme Hawley-Smoot (1930), crise bancaire européenne (1931), abandon étalon-or.
    \item \textbf{Doc 4 (Témoignage) :} Réalité vécue. Revenu passé de 7500 F à 1800 F (-76\%). Impôt inchangé (72 F). Endettement usuraire (50\%). Déscolarisation. Menace administrative. \emph{Fait clé :} « Les Blancs de l'autre côté de l'eau... leurs banques sont mortes. »
\end{itemize}

\textbf{Étape 2 : Chaîne de causalité complétée (Schéma).}
\begin{center}
\begin{tikzpicture}[node distance=1.1cm, auto, >=Stealth, every node/.style={font=\small, align=center}]
    \node[draw, rounded corners, fill=popblueL, text width=2.2cm] (A) {Krach Wall Street (Oct. 1929)\\(Doc 2)};
    \node[draw, rounded corners, fill=poporangeL, text width=2.5cm, right=of A] (B) {Effondrement du crédit mondial \& de la demande industrielle (USA/UE)\\(Doc 2, 3)};
    \node[draw, rounded corners, fill=popgreenL, text width=2.5cm, right=of B] (C) {Chute cours mondiaux cacao (-68\%)\\(Doc 2, 3)};
    \node[draw, rounded corners, fill=poppurpleL, text width=2.5cm, right=of C] (D) {Effondrement prix producteur Douala (-76\%) \& maintien impôts\\(Doc 3, 4)};
    \node[draw, rounded corners, fill=poppinkL, text width=2.5cm, right=of D] (E) {Ruine planteur : revenu divisé par 4, dettes, déscolarisation\\(Doc 4)};
    \draw[->, thick, popdark] (A) -- node[above] {Crash boursier $\rightarrow$ panique bancaire} (B);
    \draw[->, thick, popdark] (B) -- node[above] {Moins de demande $\rightarrow$ surstocks} (C);
    \draw[->, thick, popdark] (C) -- node[above] {Firmes baissent prix d'achat} (D);
    \draw[->, thick, popdark] (D) -- node[above] {Revenu < charges fixes (impôts)} (E);
\end{tikzpicture}
\end{center}

\textbf{Étape 3 : Calcul du taux de variation (1928 $\rightarrow$ 1932).}
\begin{align*}
    V_i &= 2,50 \text{ F/kg} \quad (\text{Doc 3, année 1928}) \\
    V_f &= 0,60 \text{ F/kg} \quad (\text{Doc 3, année 1932}) \\
    T &= \frac{0,60 - 2,50}{2,50} \times 100 = \frac{-1,90}{2,50} \times 100 = -0,76 \times 100 = \mathbf{-76\%}
\end{align*}
\textbf{Interprétation :} Le prix payé au planteur a baissé de \textbf{76\%} en quatre ans, soit une division par plus de 4 du revenu brut pour une production équivalente, anéantissant la capacité d'épargne et de paiement des charges fixes (impôts, outils).

\textbf{Étape 4 : Proposition de rédaction individuelle (Modèle corrigé — environ 13 lignes).}
\begin{quote}\itshape
Le cacao de grand-père ne se vendait plus à un prix rémunérateur en 1932 parce que la crise de 1929 s'est transmise violemment de la finance new-yorkaise à l'économie de plantation camerounaise par trois mécanismes enchaînés. 

Premièrement, le krach de Wall Street (octobre 1929) a provoqué l'effondrement du crédit international et de la demande industrielle américaine et européenne, comme le montre la chute synchronisée de l'indice Dow Jones (de 350 à 45 points entre 1929 et 1932, Doc 2) et de l'indice du cours mondial du cacao (de 108 à 35, base 100 en 1925, Doc 2). 

Deuxièmement, les maisons de négoce (SCOA, etc.) ont répercuté cette baisse sur le prix payé au producteur à Douala, qui s'effondre de 2,50 F/kg en 1928 à 0,60 F/kg en 1932, soit une baisse de 76\% (calcul effectué d'après Doc 3). 

Troisièmement, l'administration coloniale a maintenu la pression fiscale (impôt de capitation inchangé à 72 F pour la famille Ngono, Doc 4) alors que le revenu brut passait de 7500 F à 1800 F (Doc 4). Ngono témoigne : « Le commis m'a dit : les Blancs de l'autre côté de l'eau ne veulent plus notre cacao, leurs banques sont mortes » (Doc 4). 

En conclusion, la ruine du planteur n'est pas une fatalité climatique mais la conséquence directe de l'interdépendance capitaliste : la spéculation boursière américaine a détruit le pouvoir d'achat mondial, et le système colonial a refusé d'amortir le choc pour le producteur africain.
\end{quote}
\end{exoresolu}

\courssub{Bilan de l'activité}
\begin{aretenir}[Ce que cette activité a permis de mettre en évidence]
\begin{enumerate}
    \item \textbf{La mondialisation de la crise :} Le krach boursier n'est pas un événement abstrait ; il détruit le revenu réel d'un paysan camerounais en 3 ans via la chute des cours des matières premières (cacao, café, coton, caoutchouc).
    \item \textbf{L'asymétrie coloniale :} Les firmes européennes (Doc 1, 4) et l'État colonial (impôts Doc 4) transfèrent le risque de la crise sur le producteur le plus faible (le planteur), qui subit la baisse des prix d'achat sans voir baisser ses charges (impôts, outils, dettes).
    \item \textbf{La méthode historique :} Croiser des sources quantitatives (courbes, tableaux) et qualitatives (témoignage, carte) permet de construire une explication causale robuste, chiffrée et humanisée.
    \item \textbf{Compétences Bac :} Calcul de taux de variation, lecture de graphiques à double échelle, rédaction argumentée structurée (Thèse / Arguments chiffrés / Conclusion), citation précise des documents (Doc 1, Doc 2, Doc 3, Doc 4).
\end{enumerate}
\end{aretenir}

\begin{attention}[Pièges à éviter lors de la rédaction]
\begin{enumerate}
    \item Ne pas confondre \emph{cours mondial} (bourse de Londres/NY, Doc 2) et \emph{prix au planteur} (comptoir Douala, Doc 3) : l'écart est la marge des négociants.
    \item Ne pas écrire « la crise a touché le Cameroun » sans prouver \textbf{comment} (mécanisme : cours $\rightarrow$ prix producteur $\rightarrow$ revenu $\rightarrow$ impôt).
    \item Oublier de citer les documents explicitement : « d'après le document 3 », « le témoignage du document 4 montre que... ».
    \item Ne pas calculer le pourcentage de baisse (compétence transversale exigible).
\end{enumerate}
\end{attention}

\begin{savaistu}[Le Cameroun et la crise de 1930 : au-delà du cacao]
La crise a touché toutes les filières d'exportation : le \textbf{caoutchouc} (prix divisé par 10), le \textbf{coton}, le \textbf{bois}. Les recettes douanières du budget colonial s'effondrent, entraînant une baisse des investissements publics (routes, santé, écoles). Paradoxalement, l'administration augmente la pression fiscale (impôt de capitation, prestations) pour équilibrer son budget, aggravant la paupérisation rurale. Cela alimente les premières contestations nationalistes (ex: mouvement des \emph{Vétérans} 1930, grèves portuaires Douala 1931) qui préparent l'après-guerre.
\end{savaistu}

\newpage

\coursec{Méthodes et exercices résolus}

\courssub{Méthodes et exercices résolus}

\begin{methode}[Analyser un document historique (texte, image, tableau, carte) — Démarche en 4 temps]
\begin{enumerate}
\item \textbf{Identifier le document} : nature (texte officiel, discours, article de presse, affiche, photo, statistique, carte), auteur, date, source, contexte historique immédiat.
\item \textbf{Dégager l'idée principale} : quel est le message central ? Quelle thèse l'auteur défend-il ? Quelle information brute livre le document (chiffres, faits visibles) ?
\item \textbf{Expliquer et contextualiser} : mobiliser ses connaissances pour expliquer les références implicites, les mots-clés, les acteurs. Situer le document dans la chronologie du chapitre (ex : 1914, 1917, 1919, 1929, 1933).
\item \textbf{Conclure sur la portée / l'intérêt} : en quoi ce document éclaire-t-il le thème étudié (guerre totale, traité de Versailles, krach, New Deal) ? Quelles sont ses limites (subjectivité, partialité, champ restreint) ?
\end{enumerate}
\end{methode}

\begin{exoresolu}[Commentaire d'une affiche de propagande de la Première Guerre mondiale]
\textbf{Document :} Affiche française de 1915 « Pour la France, versez votre or. L'or combat pour la victoire ». On y voit une pièce d'or (un Napoléon) se transformant en obus qui détruit un aigle allemand.
\tcblower
\textbf{Solution détaillée :}

\noindent \textbf{1. Identification}
\begin{itemize}
\item \textbf{Nature :} Affiche de propagande (image + slogan), document iconographique officiel.
\item \textbf{Auteur / Émetteur :} Gouvernement français (souscription nationale), probablement via la Banque de France ou le Trésor public.
\item \textbf{Date :} 1915 (première année de guerre, après l'échec de la guerre de mouvement, début de la guerre d'usure).
\item \textbf{Contexte :} La guerre s'enlise (batailles d'Artois, de Champagne). Les besoins financiers sont colossaux (coût journalier de la guerre, achats d'armements aux USA). L'État français lance des emprunts nationaux et appelle à l'épargne populaire (« or »).
\end{itemize}

\noindent \textbf{2. Description et idée principale}
\begin{itemize}
\item \textbf{Visuel :} Métamorphose : une pièce d'or (symbole d'épargne, de richesse nationale) $\rightarrow$ obus (arme offensive) $\rightarrow$ destruction de l'aigle allemand (symbole du Reich, de l'ennemi).
\item \textbf{Texte :} « Pour la France, versez votre or. L'or combat pour la victoire ».
\item \textbf{Idée principale :} L'épargne privée (l'or thésaurisé par les Français) est une arme directe. Le citoyen participe à l'effort de guerre en apportant son or à l'État. L'argent = munitions = victoire.
\end{itemize}

\noindent \textbf{3. Explication et contextualisation (Mobilisation des connaissances)}
\begin{itemize}
\item \textbf{Guerre totale :} Ce document illustre la \cle{mobilisation économique et financière} de l'arrière. La distinction front/arrière s'efface.
\item \textbf{Financement :} La France ne peut payer ses importations (charbon, blé, acier, armes) qu'en or ou en devises. La « bataille de l'or » est vitale pour la balance des paiements.
\item \textbf{Propagande :} Utilisation de symboles forts (or = force, aigle = ennemi vaincu) pour susciter le patriotisme et vaincre la réticence à se séparer de son épargne. Censure et contrôle de l'opinion (Union sacrée).
\item \textbf{Chiffres clés :} 4 grands emprunts nationaux (1915, 1916, 1917, 1918) + emprunts à court terme. L'or rentre massivement (ex : environ 3,5 milliards de francs en or pour le 1er emprunt).
\end{itemize}

\noindent \textbf{4. Portée et limites}
\begin{itemize}
\item \textbf{Portée :} Témoigne de l'effort financier consenti par la population civile. Montre la puissance de l'État pour capter l'épargne. Illustre le lien \textbf{argent-guerre}.
\item \textbf{Limites :} Document officiel, unilatéral. Ne montre pas l'inflation galopante, la dette publique qui explose, ni l'endettement vis-à-vis des USA (la France devient débitrice). Ne montre pas les inégalités de sacrifice (impôt sur le revenu créé en 1914 mais faible, inflation qui appauvrit les rentiers et ouvriers).
\end{itemize}

\noindent \textbf{Conclusion :} Cette affiche est une source de premier ordre pour comprendre la \cle{guerre totale} : elle révèle comment l'État transforme l'épargne populaire en puissance de feu, condition \textit{sine qua non} de la tenue du front jusqu'en 1918.
\end{exoresolu}

\begin{methode}[Rédiger un développement structuré (type dissertation / question de synthèse) — Plan en 3 parties]
\begin{enumerate}
\item \textbf{Analyser le sujet} : définir les termes clés, délimiter l'espace (Europe, Monde) et le temps (dates bornes), identifier la problématique (cause/conséquence, rupture/continuité, ampleur).
\item \textbf{Construire le plan détaillé} (3 parties équilibrées, 2 à 3 sous-parties chacune) :
\begin{itemize}
\item \textbf{Partie I :} Les causes / Les origines / L'enclenchement (facteurs structurels + facteur déclencheur).
\item \textbf{Partie II :} Le déroulement / Les mécanismes / Les formes (phases chronologiques ou aspects thématiques : militaire, économique, social).
\item \textbf{Partie III :} Les conséquences / Les héritages / Le nouvel ordre (immédiates et à long terme, bilans).
\end{itemize}
\item \textbf{Rédiger l'introduction} : Accroche $\rightarrow$ Définitions $\rightarrow$ Débordement (contexte) $\rightarrow$ Problématique $\rightarrow$ Annonce du plan.
\item \textbf{Rédiger le développement} : Une idée directrice par paragraphe = 1 phrase d'intro + arguments précis (dates, noms, chiffres, concepts) + mini-conclusion/transition.
\item \textbf{Rédiger la conclusion} : Bilan synthétique (réponse à la problématique) $\rightarrow$ Ouverture (lien avec le chapitre suivant : traités de paix $\rightarrow$ fragilités $\rightarrow$ 1939 / crise 1929 $\rightarrow$ totalitarismes $\rightarrow$ 1939).
\end{enumerate}
\end{methode}

\begin{exoresolu}[Sujet type Probatoire : « Les conséquences de la Première Guerre mondiale ont-elles réglé les problèmes qui en étaient à l'origine ? »]
\tcblower
\textbf{Solution détaillée :}

\noindent \textbf{1. Analyse du sujet \& Problématique}
\begin{itemize}
\item \textbf{Termes clés :} « Conséquences » (traité de Versailles, nouvelles frontières, SDN, réparations, démographie, économie), « Problèmes à l'origine » (nationalismes, impérialismes, alliances rigides, course aux armements, question des Balkans, rivalités franco-allemande et anglo-allemande).
\item \textbf{Bornes :} 1919 (traité de Versailles) -- 1939 (déclenchement 2e GM) ou 1929 (crise) comme tournant.
\item \textbf{Problématique :} Le traité de Versailles et l'ordre de 1919 ont-ils apporté une paix durable en résolvant les contentieux (territoriaux, nationalistes, économiques) ou ont-ils porté en germe les germes d'un nouveau conflit par leur caractère punitif et incomplet ?
\end{itemize}

\noindent \textbf{2. Plan détaillé}

\noindent \textbf{Introduction} : « La der des ders » (Clemenceau) ? 10 millions de morts. Traité de Versailles (28 juin 1919). \textbf{Problématique} ci-dessus. \textbf{Plan :} I. Une tentative de règlement territorial et institutionnel. II. Des règlements économiques et sécuritaires sources de tensions. III. L'échec structurel : les germes de la revanche et de l'instabilité (1919-1939).

\noindent \textbf{3. Développement}

\noindent \textbf{I. Une tentative de règlement territorial et institutionnel (L'apaisement apparent)}
\begin{itemize}
\item \textbf{A. Redessin des frontières selon le « droit des peuples » (Wilson / 14 points) :} Disparition des empires multinationaux (Autriche-Hongrie, Ottoman, Russe). Création d'États-nations : Pologne, Tchécoslovaquie, Yougoslavie, Pays Baltes. Restitution Alsace-Lorraine à la France.
\item \textbf{B. Création de la SDN (Société des Nations) :} Outil de sécurité collective, arbitrage, désarmement. Siège à Genève. Succès initiaux (conflit Åland, Haute-Silésie, réfugiés Fridtjof Nansen).
\item \textbf{C. Démantèlement militaire allemand :} Armée limitée à 100 000 hommes, pas d'aviation, pas de chars, Rhénanie démilitarisée. Désarmement imposé aux vaincus seulement.
\end{itemize}

\noindent \textbf{II. Des règlements économiques et sécuritaires sources de tensions (L'injustice vécue)}
\begin{itemize}
\item \textbf{A. Le « Diktat » de Versailles :} Traité imposé sans négociation (mai-juin 1919). Article 231 (responsabilité de guerre) $\rightarrow$ humiliation nationale allemande. Perte de 13\% territoire, 10\% population, 75\% fer, 26\% charbon.
\item \textbf{B. Réparations intenables :} 132 milliards marks-or (1921). Crise de la Ruhr (1923), hyperinflation. Plan Dawes (1924) / Young (1929) : dépendance au dollar US. Sentiment d'esclavage financier.
\item \textbf{C. Insécurité française vs révisionnisme allemand :} France veut sécurité (alliances Petite Entente, traité Rhineland). Allemagne (Streseemann) joue l'apaisement apparent mais vise la révision (Locarno 1925 : frontières OUEST garanties, EST négociables).
\item \textbf{D. Exclusion de l'URSS et des USA :} USA refusent SDN (isolationnisme). URSS paria (révolution bolchévique). SDN privée de 2 puissances majeures.
\end{itemize}

\noindent \textbf{III. L'échec structurel : les germes de la revanche et de l'instabilité (1919-1939)}
\begin{itemize}
\item \textbf{A. Nationalismes insatisfaits :} Minorités allemandes en Tchécoslovaquie (Sudètes), Pologne (Dantzig, Corridor), Autriche (Anschluss interdit). Minorités hongroises (Traité de Trianon). Irrédentisme italien (Trente, Trieste, Dalmatie promises non tenues $\rightarrow$ « victoire mutilée »).
\item \textbf{B. Fragilité économique mondiale :} Europe endettée vis-à-vis USA. Allemagne paye réparations $\rightarrow$ Alliés payent dettes US $\rightarrow$ cercle vicieux dépendant des capitaux US. Crise 1929 brise l'équilibre précaire (fin plans Dawes/Young).
\item \textbf{C. Montée des révisionnismes totalitaires :} Hitler (1933) : programme explicite (détruire Versailles, Lebensraum, Anschluss). Italie (Mussolini) : empire colonial. Japon : Mandchourie (1931). SDN impuissante (retraits : Japon 33, Allemagne 33, Italie 37).
\item \textbf{D. Rupture définitive :} Remilitarisation Rhénanie (1936), Anschluss (1938), Munich (1938), Prague (1939). L'ordre de 1919 s'effondre.
\end{itemize}

\noindent \textbf{4. Conclusion}
\begin{itemize}
\item \textbf{Bilan :} L'ordre de 1919 a \textbf{échoué} à régler les problèmes de fond. Il a gelé les nationalismes au lieu de les résoudre, humilié l'Allemagne sans la détruire, ruiné l'Europe sans la reconstruire, et créé une SDN sans dents ni membres majeurs.
\item \textbf{Ouverture :} Ce « traité de paix » de 20 ans (Foch) a posé les bases de la \cle{Seconde Guerre mondiale}. Seule la défaite totale de 1945 et la Guerre froide permettront (via CECA/CEE/UE et ONU) une construction européenne pacifique et une sécurité collective intégrant l'Allemagne et l'URSS/Russie (provisoirement).
\end{itemize}
\end{exoresolu}

\begin{methode}[Construire et exploiter une frise chronologique thématique — Repères temporels et causalités]
\begin{enumerate}
\item \textbf{Sélectionner les dates charnières} (10 à 15 max) : bornes du thème, événements déclencheurs, tournants, aboutissements. Regrouper par \textbf{thèmes} (Politique/Diplomatie, Militaire, Économique/Social, Idéologique).
\item \textbf{Tracer l'axe} : Graduation régulière (ex : 1 cm = 1 an ou 6 mois). Indiquer les bornes (1914-1919 / 1929-1939).
\item \textbf{Positionner les événements} : Au-dessus de l'axe : Politique/Diplomatie. Au-dessous : Militaire/Éco/Social. Utiliser des codes couleur (Rouge = Guerre/Conflit, Bleu = Diplomatie/Traités, Vert = Éco/Crise, Orange = Idéologie/Lois).
\item \textbf{Relier par des flèches causales} : Flèche continue = cause directe. Flèche pointillée = influence indirecte / condition. Annoter la flèche (ex : « Traité de Versailles $\xrightarrow{\text{humiliation/réparations}}$ Montée nazie »).
\item \textbf{Légender et titrer} : Titre précis, légende des codes, source si doc fourni.
\end{enumerate}
\end{methode}

\begin{exoresolu}[Frise chronologique commentée : « De la Grande Guerre à la Grande Crise (1914-1933) »]
\tcblower
\textbf{Solution détaillée :}

\noindent \textbf{1. Sélection des dates charnières (13 repères)}

\begin{center}
\begin{tabularx}{\linewidth}{|c|X|X|}\hline
\textbf{Date} & \textbf{Événement (Thème)} & \textbf{Signification causale} \\ \hline
1914 (28 juin) & Attentat de Sarajevo (Pol/Dip) & Déclencheur immédiat (étincelle). \\ \hline
1914 (août) & Déclenchement guerre / Plan Schlieffen échoue (Mil) & Guerre de mouvement $\rightarrow$ guerre de position (Marne). \\ \hline
1917 (avril) & Entrée en guerre des USA (Pol/Mil) & Bascule décisive : ressources fraîches, idéalisme wilsonien. \\ \hline
1917 (nov.) & Révolution bolchévique Russie (Pol/Idéo) & Retrait Russie (Brest-Litovsk 1918), peur rouge en Europe. \\ \hline
1918 (11 nov.) & Armistice de Rethondes (Mil) & Fin des combats, mais pas paix. Blocus continue. \\ \hline
1919 (28 juin) & Traité de Versailles (Dip) & Ordre de paix punitif, SDN, réparations, frontières nouvelles. \\ \hline
1923 & Occupation de la Ruhr / Hyperinflation ALL (Éco) & Échec réparations, trauma social allemand, Dawes 1924. \\ \hline
1924-1929 & « Années folles » / Stabilisation relative (Éco/Dip) & Prêts US $\rightarrow$ prospérité artificielle, Locarno, Briand-Kellogg. \\ \hline
1929 (24 oct.) & Jeudi noir / Krach Wall Street (Éco) & Effondrement boursier $\rightarrow$ crise bancaire $\rightarrow$ crise réelle. \\ \hline
1930 (juin) & Loi Hawley-Smoot USA (Éco) & Protectionnisme $\rightarrow$ effondrement commerce mondial (-60\%). \\ \hline
1931 (mai) & Krach Créditanstalt Vienne / Crise sterling (Éco) & Contagion européenne, fin étalon-or (UK sept 31). \\ \hline
1933 (30 janv.) & Hitler Chancelier (Pol/Idéo) & Prise de pouvoir légal $\rightarrow$ fin démocratie Weimar, réarmement. \\ \hline
1933 (mars) & New Deal Roosevelt (Éco/Pol) & Interventionnisme État US, relance par dépense publique. \\ \hline
\end{tabularx}
\end{center}

\noindent \textbf{2. Représentation graphique (Code TikZ corrigé : alignement précis sur l'axe temporel)}

\begin{popfigure}
\begin{tikzpicture}[scale=0.85, every node/.style={font=\footnotesize}]
% Axe principal : 1914 à 1933 (19 ans). 1 an = 0.75cm. Longueur ~14.25cm.
\draw[thick, ->] (0,0) -- (15,0) node[right] {Temps};
% Graduations annuelles
\foreach \y [count=\i from 0] in {1914,1915,1916,1917,1918,1919,1920,1921,1922,1923,1924,1925,1926,1927,1928,1929,1930,1931,1932,1933} {
    \pgfmathsetmacro{\x}{\i*0.75}
    \draw (\x, -0.1) -- (\x, 0.1) node[below] {\y};
}
% Événements AU-DESSUS (Politique/Diplo/Militaire)
% Coordonnées x = (Année - 1914) * 0.75
\node[above, align=center, popred] at (0, 1.8) {1914\\Sarajevo\\Guerre};
\node[above, align=center, popred] at (1.5, 2.1) {1916\\Verdun/Somme};
\node[above, align=center, popblue] at (2.25, 2.4) {1917\\USA / URSS};
\node[above, align=center, popred] at (3.0, 1.8) {1918\\Armistice};
\node[above, align=center, popblue] at (3.75, 2.2) {1919\\Versailles / SDN};
\node[above, align=center, popblue] at (8.25, 2.0) {1925\\Locarno};
\node[above, align=center, poppurple] at (11.25, 2.2) {1929\\Krach};
\node[above, align=center, poppurple] at (14.25, 2.0) {1933\\Hitler / New Deal};
% Flèches causales (au-dessus)
\draw[->, thick, popred, dashed] (0.5, 1.9) -- (1.2, 2.0) node[midway, above, sloped, xshift=-0.3cm] {Enlisement};
\draw[->, thick, popblue] (2.5, 2.3) -- (2.8, 1.9) node[midway, above, sloped] {Bascule};
\draw[->, thick, popblue] (3.2, 1.9) -- (3.5, 2.1) node[midway, above, sloped] {Paix imposée};
\draw[->, thick, poporange, dashed] (4.2, 2.1) -- (7.8, 2.0) node[midway, above, sloped, xshift=1cm] {Apaisement fragile};
\draw[->, thick, poppurple] (11.5, 2.1) -- (13.8, 2.0) node[midway, above, sloped] {Crise $\rightarrow$ Totalitarisme};
% Événements EN-DESSOUS (Éco/Social)
\node[below, align=center, popgreen] at (6.75, -1.8) {1923\\Ruhr / Hyperinflation};
\node[below, align=center, popgreen] at (9.4, -2.1) {1924-29\\Dawes / Prospérité $*$};
\node[below, align=center, poppurple] at (11.25, -1.8) {1929\\Krach};
\node[below, align=center, poppurple] at (12.75, -2.1) {1931\\Crise bancaire EU};
% Flèches causales (en-dessous)
\draw[->, thick, popgreen] (7.0, -1.9) -- (9.0, -2.0) node[midway, below, sloped] {Stabilisation $ (US) $};
\draw[->, thick, poppurple] (11.5, -1.9) -- (12.5, -2.0) node[midway, below, sloped] {Contagion};
\end{tikzpicture}
\legende{Frise thématique 1914-1933 : Au-dessus (Politique/Militaire), En-dessous (Économique). Flèches pleines = causalité forte ; pointillées = conditionnel/fragile. $*$ Prospérité dépendante des capitaux US.}
\end{popfigure}

\noindent \textbf{3. Exploitation / Analyse causale (Réponse type Probatoire)}
\begin{itemize}
\item \textbf{Rupture 1914-1919 :} La guerre (rouge) crée les conditions de la paix (bleu) mais la paix (Versailles) porte les germes de l'instabilité (flèche pointillée vers 1923 Ruhr).
\item \textbf{Stabilisation 1924-1929 (Zone verte) :} Artificielle car \textbf{dépendante des prêts US} (Dawes/Young). Flèche verte de 1923 vers 1924 annotée « Capitaux US ».
\item \textbf{Bascule 1929 (Violet) :} Le krach (éco) brise le cycle de la dette. Flèche violette vers 1931 (crise bancaire EU) puis vers 1933 (Hitler). \textbf{Lien direct :} Crise $\rightarrow$ Chômage (6 millions ALL) $\rightarrow$ Vote nazi $\rightarrow$ Hitler.
\item \textbf{Conclusion :} La frise montre l'\textbf{enchaînement déterministe} : Guerre $\rightarrow$ Paix injuste $\rightarrow$ Dette/Inflation $\rightarrow$ Dépendance US $\rightarrow$ Krach $\rightarrow$ Effondrement démocratique $\rightarrow$ Guerre. La SDN (1919) n'a pas cassé la chaîne.
\end{itemize}
\end{exoresolu}

\begin{methode}[Réaliser et commenter un croquis cartographique (Schémas de flux, fronts, aires d'influence) — Code graphique]
\begin{enumerate}
\item \textbf{Analyser le sujet} : Thème (Fronts 1914-1918, Propagande, Crise 1929, Commerce mondial, Alliances). Échelle (Monde, Europe, Bassin méditerranéen).
\item \textbf{Choisir le fond de carte simplifié} : Contours continents/pays au trait fin gris. Villes repères (capitales, ports, centres industriels). Fleuves/montagnes si pertinents (Rhénanie, Ruhr, Détroits).
\item \textbf{Appliquer le code graphique (Sémiologie graphique)} :
\begin{itemize}
\item \textbf{Surfaces (aires) :} Zones d'influence, empires, pays victorieux/vaincus, zones occupées, aires de crise. Hachures ou couleurs claires (ex : hachures rouges = fronts, jaune = zones industrielles, bleu = zones d'influence US).
\item \textbf{Lignes (flux, frontières, axes) :} Traits pleins épais = frontières nouvelles (1919) ou flux majeurs (commerce, prêts US). Traits pointillés = frontières anciennes (1914) ou flux secondaires. Flèches proportionnelles = flux financiers (réparations, dettes, capitaux US), migrations, offensives.
\item \textbf{Points (symboles) :} Carrés = capitales, cercles = ports/centres industriels (Ruhr, Pittsburgh), triangles = ressources (charbon, fer), étoiles = lieux de traités (Versailles, Locarno, Washington).
\item \textbf{Figurés de surface/linéaires/ponctuels} : Respecter l'ordre visuel (surfaces en fond, lignes au milieu, points au-dessus).
\end{itemize}
\item \textbf{Construire la légende organisée} : Trois colonnes (Figuré / Signification / Information). Grouper par thème (Politique, Militaire, Économique). Titre de la légende = Thème du croquis.
\item \textbf{Compléter les éléments obligatoires} : Titre précis (Thème + Espace + Temps), Orientation (flèche Nord), Échelle indicative (ex : 1 cm $\approx$ 500 km), Source (si doc fourni), Auteur/Date.
\item \textbf{Rédiger le commentaire structuré} : 1) Description globale (qué voit-on ?). 2) Analyse par thème (frontières, flux, pôles de puissance). 3) Interprétation / Synthesis (quels rapports de force, quelles dynamiques ?). 4) Limites (échelle, choix arbitraires, données manquantes).
\end{enumerate}
\end{methode}

\begin{exemplebox}[Exemple de croquis : « L'Europe des traités de 1919-1920 »]
\textbf{Consigne :} Réaliser un croquis montrant les nouveaux États, les pertes territoriales allemandes/russes/autrichiennes, les zones de plébiscite et la Rhénanie démilitarisée.
\textbf{Figurés suggérés :}
\begin{itemize}
\item \textbf{Fond :} Europe 1914 (gris clair) vs Europe 1920 (contours noirs épais).
\item \textbf{Aires :} Nouveaux États (Pologne, Tchécoslovaquie, Yougoslavie, Hongrie, Autriche) en couleurs pastel distinctes. Allemagne 1919 (hachures rouges fines). Zones de plébiscite (hachures croisées). Rhénanie (bande bleue hachurée).
\item \textbf{Lignes :} Frontières 1914 (trait fin pointillé gris). Frontières 1919 (trait plein noir épais). Ligne de démarcation polonaise (Curzon) trait mixte.
\item \textbf{Points :} Capitales (carrés noirs). Dantzig (cercle rouge « Ville libre »). Saar (triangle charbon). Memel (losange).
\item \textbf{Légende :} Organisée en « Frontières », « États », « Statuts spéciaux », « Ressources/Enjeux ».
\end{itemize}
\textbf{Commentaire type :} « Ce croquis met en évidence l'application du droit des peuples (création États-nations) mais aussi ses limites : minorités allemandes enfermées (Sudètes, Corridor), Allemagne amputée mais intacte démographiquement, Rhénanie démilitarisée source de tension future. L'Europe centrale devient un « cordon sanitaire » fragile face à l'URSS. »
\end{exemplebox}

\newpage

\coursec{Exercices}

\courssub{Je vérifie mes connaissances}

\begin{exercice}[QCM : les repères essentiels]
Pour chaque question, une seule réponse est exacte. Entoure la bonne lettre.
\begin{enumerate}
\item L'attentat de Sarajevo, événement déclencheur de la Première Guerre mondiale, a lieu en :
\begin{enumerate}
\item[a)] 1905 \quad b)] 1914 \quad c)] 1918
\end{enumerate}
\item La première bataille de la Marne, qui arrête l'avance allemande vers Paris, se déroule en :
\begin{enumerate}
\item[a)] septembre 1914 \quad b)] novembre 1918 \quad c)] octobre 1929
\end{enumerate}
\item Le traité qui met officiellement fin à la guerre avec l'Allemagne est :
\begin{enumerate}
\item[a)] le traité de Brest-Litovsk \quad b)] le traité de Versailles \quad c)] le traité de Berlin
\end{enumerate}
\item Le krach boursier de Wall Street (le « jeudi noir ») a lieu le :
\begin{enumerate}
\item[a)] 11 novembre 1918 \quad b)] 4 février 1922 \quad c)] 24 octobre 1929
\end{enumerate}
\item Le programme de relance économique lancé par le président Roosevelt à partir de 1933 s'appelle :
\begin{enumerate}
\item[a)] le New Deal \quad b)] le plan Marshall \quad c)] la SDN
\end{enumerate}
\end{enumerate}
\end{exercice}

\begin{exercice}[Vrai ou faux ? Justifie]
Réponds par VRAI ou FAUX, puis justifie ta réponse en une ou deux phrases.
\begin{enumerate}
\item La Première Guerre mondiale oppose la Triple-Entente à la Triple-Alliance.
\item La guerre de 1914-1918 est restée une guerre de mouvement du début à la fin.
\item Le Cameroun allemand a été épargné par les combats de la Première Guerre mondiale.
\item Après la guerre, le Cameroun est placé sous mandat de la Société des Nations et partagé entre la France et la Grande-Bretagne.
\item La crise de 1929 n'a touché que les États-Unis.
\item La crise de 1929 a provoqué une baisse des cours des matières premières comme le cacao.
\end{enumerate}
\end{exercice}

\begin{exercice}[Définis les mots-clés]
Donne une définition courte et précise (deux ou trois lignes maximum) de chacun des termes suivants :
\begin{enumerate}
\item \cle{Guerre totale}
\item \cle{Mandat} (de la Société des Nations)
\item \cle{Krach boursier}
\item \cle{Surproduction}
\item \cle{New Deal}
\end{enumerate}
\end{exercice}

\begin{exercice}[Calculs directs sur des données de la crise]
Aux États-Unis, le nombre de chômeurs passe d'environ 1,5 million en 1929 à environ 12,8 millions en 1933.
\begin{enumerate}
\item Calcule l'augmentation du nombre de chômeurs entre 1929 et 1933.
\item Calcule le taux de variation du nombre de chômeurs entre 1929 et 1933 (en \%, arrondi à l'unité).
\item En 1933, la population active américaine est d'environ 51 millions de personnes. Calcule le taux de chômage en 1933 (en \%, arrondi à l'unité).
\end{enumerate}
\end{exercice}

\courssub{Je m'entraîne}

\begin{exercice}[Frise chronologique]
Construis une frise chronologique graduée de 1914 à 1933 et places-y les événements suivants : début de la Première Guerre mondiale ; première bataille de la Marne ; bataille de Verdun ; entrée en guerre des États-Unis ; armistice du 11 novembre 1918 ; traité de Versailles ; confirmation par la SDN des mandats sur le Cameroun (1922) ; krach de Wall Street ; arrivée de Roosevelt au pouvoir. Pour chaque événement, indique la date exacte (année, et le mois lorsqu'il est connu).
\end{exercice}

\begin{exercice}[Croquis de la campagne du Cameroun]
À partir d'un fond de carte muet du Cameroun fourni par ton professeur, réalise un croquis légendé de la campagne du Cameroun (1914-1916) :
\begin{enumerate}
\item Reporte et nomme : Douala, Yaoundé, Garoua, Mora, le fleuve Sanaga.
\item Trace par des flèches les axes de progression des troupes franco-britanniques (depuis le Nigeria à l'ouest, le Tchad au nord et le Congo français au sud).
\item Indique par un symbole la retraite des troupes allemandes vers Yaoundé, puis vers la Guinée espagnole (1916).
\item Rédige un titre, une légende numérotée et une orientation (flèche du nord).
\end{enumerate}
\end{exercice}

\begin{methode}[Réussir un croquis d'Histoire au Probatoire]
\begin{enumerate}
\item Dessine d'abord le contour du territoire au crayon, simplement, sans chercher la précision.
\item Place les lieux (villes, fleuves) par des points ou des symboles nommés.
\item Trace les flèches des offensives en respectant le sens réel des déplacements.
\item Rédige un titre complet (quoi, où, quand) : « La campagne du Cameroun, 1914-1916 ».
\item Numérote la légende et ajoute la flèche du nord : un croquis sans titre ni légende perd des points.
\end{enumerate}
\end{methode}

\begin{exercice}[Diagramme en barres du chômage]
Voici le nombre de chômeurs (en millions) aux États-Unis et en Allemagne :
\begin{center}
\begin{tabularx}{\linewidth}{|l|X|X|X|X|X|}
\hline
\rowcolor{popblueL} Année & 1929 & 1930 & 1931 & 1932 & 1933 \\
\hline
États-Unis & 1,5 & 4,3 & 8,0 & 12,1 & 12,8 \\
\hline
Allemagne & 1,3 & 3,1 & 4,5 & 5,6 & 4,8 \\
\hline
\end{tabularx}
\end{center}
\begin{enumerate}
\item Construis sur ton cahier un diagramme en barres groupées (une barre par pays et par année), avec titre, axes légendés et légende.
\item Pour chaque pays, indique l'année où le chômage est maximal.
\item Compare l'évolution du chômage dans les deux pays entre 1929 et 1932 en deux phrases.
\end{enumerate}
\end{exercice}

\begin{exercice}[Du krach de New York au planteur de cacao]
Remets dans l'ordre logique les étapes suivantes qui expliquent comment la crise de 1929 atteint le Cameroun, puis rédige un paragraphe de 6 à 8 lignes qui les enchaîne avec des connecteurs (d'abord, ensuite, par conséquent, enfin) :
\begin{itemize}
\item baisse des revenus des planteurs camerounais et misère rurale ;
\item krach boursier de Wall Street (octobre 1929) ;
\item effondrement des cours mondiaux du cacao et du café ;
\item faillites des banques et des entreprises aux États-Unis puis en Europe ;
\item baisse de la demande européenne en matières premières coloniales ;
\item chômage massif et baisse du pouvoir d'achat des consommateurs.
\end{itemize}
\end{exercice}

\courssub{Je me prépare au Bac}

\begin{exercice}[Dissertation guidée — type Probatoire]
\textbf{Sujet :} Montrez que la Première Guerre mondiale fut une guerre totale et exposez ses conséquences pour le Cameroun.
\begin{enumerate}
\item Analyse le sujet : dégage les deux parties imposées et définis l'expression « guerre totale » dans une introduction de 5 à 8 lignes (accroche, définition, problématique, annonce du plan).
\item Première partie : à l'aide de trois arguments précis (mobilisation des hommes, mobilisation de l'économie et des femmes, mobilisation des esprits et des colonies), montre que le conflit de 1914-1918 a engagé toutes les ressources des États belligérants. Rédige cette partie en 12 à 15 lignes.
\item Deuxième partie : expose les conséquences de la guerre pour le Cameroun (campagne de 1914-1916, fin de la colonisation allemande, partage franco-britannique, statut de mandat de la SDN confirmé en 1922). Rédige cette partie en 10 à 12 lignes.
\item Rédige une conclusion de 4 à 6 lignes qui répond à la problématique et ouvre sur la crise de 1929.
\end{enumerate}
\end{exercice}

\begin{methode}[Construire une dissertation d'Histoire en deux parties]
\begin{enumerate}
\item Souligne les mots-clés du sujet : ici « guerre totale » (partie 1) et « conséquences pour le Cameroun » (partie 2).
\item Dans l'introduction : accroche (une date, un fait), définition des termes, problématique (une question), annonce du plan en deux parties.
\item Dans chaque partie : une idée par paragraphe, appuyée par un fait précis daté et localisé.
\item Dans la conclusion : réponds clairement à la problématique, puis ouvre sur un sujet voisin (ici, l'après-guerre et la crise de 1929).
\end{enumerate}
\end{methode}

\begin{exercice}[Étude de documents — type Baccalauréat technique]
\textbf{Document 1 :} photographie décrite : soldats dans une tranchée boueuse de la Somme (1916), casques, sacs de terre, barbelés, terrain dévasté par les obus.

\textbf{Document 2 :} pertes militaires (morts, en millions, valeurs arrondies) :
\begin{center}
\begin{tabularx}{\linewidth}{|l|X|X|X|X|}
\hline
\rowcolor{popblueL} Pays & Allemagne & France & Royaume-Uni & Russie \\
\hline
Morts militaires (millions) & 2,0 & 1,4 & 0,9 & 1,8 \\
\hline
\end{tabularx}
\end{center}
\begin{enumerate}
\item Présente le document 1 (nature, date, lieu, auteur probable) en 3 à 4 lignes.
\item À partir du document 1, dégage deux caractéristiques de la guerre de tranchées (conditions de vie des soldats, violence des combats, immobilité du front).
\item À partir du document 2, calcule le total des morts militaires de ces quatre pays et classe-les du plus touché au moins touché.
\item En t'appuyant sur les deux documents et sur tes connaissances, rédige un commentaire de 12 à 15 lignes montrant que la guerre de 1914-1918 fut une guerre de masse au bilan humain catastrophique.
\end{enumerate}
\end{exercice}

\begin{exercice}[Question de synthèse — la crise de 1929, une crise mondiale]
\textbf{Données :} entre 1929 et 1932, la production industrielle mondiale baisse d'environ 38 \% ; le commerce mondial chute d'environ 65 \% ; le cours du cacao passe d'environ 22 centimes le kilogramme (équivalent monétaire de l'époque) à moins de 10 centimes ; le chômage touche environ 30 millions de personnes dans le monde en 1932.
\begin{enumerate}
\item Rappelle les causes de la prospérité américaine des années 1920 et explique en quoi elle reposait sur des fragilités (surproduction, spéculation, crédit facile). (6 à 8 lignes)
\item Explique le mécanisme du krach du 24 octobre 1929 et les raisons de sa propagation à l'Europe puis aux colonies africaines. (8 à 10 lignes)
\item À l'aide des données chiffrées, montre que la crise de 1929 est une crise mondiale qui frappe aussi le Cameroun. (6 à 8 lignes)
\item Présente en 5 à 6 lignes une réponse apportée à la crise : le New Deal de Roosevelt (à partir de 1933).
\end{enumerate}
\end{exercice}



\newpage

\coursec{Corrigés des exercices}

\coursec{Corrigés des exercices — La Première Guerre mondiale et la crise de 1929}

\begin{corrige}[Exercice 1 — QCM : les repères essentiels]
\begin{enumerate}
\item \textbf{b)} L'attentat de Sarajevo a lieu le 28 juin 1914 : l'héritier du trône d'Autriche-Hongrie, l'archiduc François-Ferdinand, est assassiné par le nationaliste serbe Gavrilo Princip. Cet événement déclenche la crise de juillet puis la guerre.
\item \textbf{a)} La première bataille de la Marne se déroule en septembre 1914 ; elle stoppe l'avance allemande vers Paris et marque le début de la guerre de position.
\item \textbf{b)} Le traité de Versailles est signé le 28 juin 1919. Le traité de Brest-Litovsk (mars 1918) ne concerne que la Russie et l'Allemagne ; le traité de Berlin date de 1878.
\item \textbf{c)} Le krach de Wall Street a lieu le jeudi 24 octobre 1929, le « jeudi noir ».
\item \textbf{a)} Le New Deal (« nouvelle donne ») est le programme de relance de Franklin D. Roosevelt, élu en 1932 et entré en fonctions en mars 1933. Le plan Marshall date de 1947 et la SDN est créée en 1919.
\end{enumerate}
\textbf{Point de vigilance :} ne pas confondre l'armistice (11 novembre 1918, arrêt des combats) et le traité de Versailles (28 juin 1919, paix officielle).
\end{corrige}

\begin{corrige}[Exercice 2 — Vrai ou faux ? Justifie]
\begin{enumerate}
\item \textbf{VRAI.} La Triple-Entente (France, Royaume-Uni, Russie) affronte la Triple-Alliance (Allemagne, Autriche-Hongrie, Italie — qui change de camp en 1915).
\item \textbf{FAUX.} Après l'échec du plan Schlieffen et la bataille de la Marne (septembre 1914), le front de l'Ouest se fige : c'est la guerre de position et des tranchées jusqu'en 1918. La guerre de mouvement ne dure que quelques mois en 1914.
\item \textbf{FAUX.} Le Cameroun allemand (Kamerun) est le théâtre de la campagne du Cameroun (1914-1916) : les troupes franco-britanniques en chassent les Allemands, qui se replient vers la Guinée espagnole en 1916.
\item \textbf{VRAI.} Le Cameroun est partagé entre la France (la plus grande partie, à l'est) et la Grande-Bretagne (deux bandes à l'ouest, accolées au Nigeria) ; la Société des Nations confirme ces mandats en 1922.
\item \textbf{FAUX.} Partie des États-Unis, la crise gagne l'Europe (retrait des capitaux américains, faillites bancaires dès 1931) puis les colonies africaines par la baisse de la demande et des cours des matières premières : c'est une crise mondiale.
\item \textbf{VRAI.} L'effondrement de la demande européenne fait chuter les cours du cacao et du café, ce qui ruine de nombreux planteurs camerounais.
\end{enumerate}
\textbf{Point de vigilance :} une réponse « vrai » ou « faux » sans justification ne rapporte aucun point : la justification datée est exigée.
\end{corrige}

\begin{corrige}[Exercice 3 — Définis les mots-clés]
\begin{enumerate}
\item \cle{Guerre totale} : guerre qui mobilise toutes les ressources d'un État — les hommes (soldats et civils), l'économie (usines, femmes au travail), les finances, la propagande et les empires coloniaux — pour écraser l'adversaire.
\item \cle{Mandat} : mission confiée par la Société des Nations à une puissance victorieuse (France ou Grande-Bretagne pour le Cameroun) pour administrer un territoire pris aux vaincus, en principe pour le préparer à se gouverner lui-même.
\item \cle{Krach boursier} : effondrement brutal et rapide des cours des actions à la Bourse, provoqué par des ventes massives et paniques, comme à Wall Street le 24 octobre 1929.
\item \cle{Surproduction} : production de biens (industriels ou agricoles) supérieure à ce que les consommateurs peuvent acheter, ce qui entraîne la baisse des prix, les stocks invendus puis les faillites.
\item \cle{New Deal} : « nouvelle donne » ; politique de relance menée par le président américain Roosevelt à partir de 1933 : grands travaux publics, aide aux chômeurs et aux paysans, réforme des banques, intervention de l'État dans l'économie.
\end{enumerate}
\textbf{Point de vigilance :} une bonne définition comporte la nature du terme, un élément de contexte (date, lieu) et, si possible, un exemple.
\end{corrige}

\begin{corrige}[Exercice 4 — Calculs directs sur des données de la crise]
\begin{enumerate}
\item Augmentation : $12{,}8 - 1{,}5 = 11{,}3$ millions de chômeurs supplémentaires entre 1929 et 1933.
\item Taux de variation :
\[ \frac{12{,}8 - 1{,}5}{1{,}5} \times 100 = \frac{11{,}3}{1{,}5} \times 100 \approx 753\,\%. \]
Le nombre de chômeurs a été multiplié par environ $12{,}8 / 1{,}5 \approx 8{,}5$, soit une hausse d'environ 753 \%.
\item Taux de chômage en 1933 :
\[ \frac{12{,}8}{51} \times 100 \approx 25\,\%. \]
Environ un actif américain sur quatre est sans emploi en 1933.
\end{enumerate}
\textbf{Point de vigilance :} le taux de variation se calcule toujours par rapport à la valeur de départ (1929), jamais par rapport à la valeur d'arrivée.
\end{corrige}

\begin{corrige}[Exercice 5 — Frise chronologique]
Dates attendues sur la frise (graduée de 1914 à 1933) :
\begin{itemize}
\item Août 1914 : début de la Première Guerre mondiale (entrée en guerre des grandes puissances européennes).
\item Septembre 1914 : première bataille de la Marne.
\item Février--décembre 1916 : bataille de Verdun.
\item Avril 1917 : entrée en guerre des États-Unis aux côtés de l'Entente.
\item 11 novembre 1918 : armistice de Rethondes, fin des combats.
\item 28 juin 1919 : traité de Versailles.
\item 1922 : la Société des Nations confirme les mandats français et britannique sur le Cameroun.
\item 24 octobre 1929 : krach de Wall Street (« jeudi noir »).
\item Mars 1933 : Roosevelt entre en fonctions et lance le New Deal.
\end{itemize}

\begin{popfigure}
\begin{tikzpicture}[scale=0.9, every node/.style={font=\small}]
% Axe principal
\draw[thick, ->] (0,0) -- (19,0) node[right] {Années};
% Graduations principales (tous les 2 ans)
\foreach \y/\x in {1914/0, 1916/2, 1918/4, 1920/6, 1922/8, 1924/10, 1926/12, 1928/14, 1930/16, 1932/18} {
    \draw (\x, -0.1) -- (\x, 0.1) node[below=2pt] {\x};
}
% Événements (position relative approximative)
% 1914: Août (0.15), Sept (0.25)
\draw[popblue, thick] (0.15, 0.5) -- (0.15, 0.2);
\node[popblue, above, align=center] at (0.15, 0.55) {Août 1914\\Début GM};
\draw[popblue, thick] (0.25, -0.7) -- (0.25, -0.4);
\node[popblue, below, align=center] at (0.25, -0.75) {Sept 1914\\Marne};

% 1916: Verdun (milieu 1916 ~ 2.5)
\draw[poporange, thick] (2.5, 0.5) -- (2.5, 0.2);
\node[poporange, above, align=center] at (2.5, 0.55) {Fév--Déc 1916\\Verdun};

% 1917: US entry (Avril ~ 3.25)
\draw[popgreen, thick] (3.25, -0.7) -- (3.25, -0.4);
\node[popgreen, below, align=center] at (3.25, -0.75) {Avril 1917\\Entrée USA};

% 1918: Armistice (Nov ~ 4.8)
\draw[poppurple, thick] (4.8, 0.5) -- (4.8, 0.2);
\node[poppurple, above, align=center] at (4.8, 0.55) {11 Nov 1918\\Armistice};

% 1919: Versailles (Juin ~ 5.4)
\draw[poppink, thick] (5.4, -0.7) -- (5.4, -0.4);
\node[poppink, below, align=center] at (5.4, -0.75) {28 Juin 1919\\Versailles};

% 1922: Mandats Cameroun (~ 8)
\draw[popgold, thick] (8, 0.5) -- (8, 0.2);
\node[popgold, above, align=center] at (8, 0.55) {1922\\Mandats Cameroun};

% 1929: Krach (Oct ~ 15.75)
\draw[popdark, thick] (15.75, -0.7) -- (15.75, -0.4);
\node[popdark, below, align=center] at (15.75, -0.75) {24 Oct 1929\\Krach};

% 1933: New Deal (Mars ~ 19.1 -> hors cadre, on met à 18.9)
\draw[popblue, thick] (18.9, 0.5) -- (18.9, 0.2);
\node[popblue, above, align=center] at (18.9, 0.55) {Mars 1933\\New Deal};
\end{tikzpicture}
\legende{Frise chronologique 1914--1933 : repères majeurs de la Grande Guerre et de la crise de 1929.}
\end{popfigure}

\textbf{Point de vigilance :} les événements doivent être placés dans l'ordre croissant, à peu près à la bonne distance les uns des autres, et chaque repère doit porter sa date.
\end{corrige}

\begin{corrige}[Exercice 6 — Croquis de la campagne du Cameroun]
Éléments attendus sur le croquis :
\begin{enumerate}
\item \textbf{Lieux reportés et nommés :} Douala (sur la côte), Yaoundé (au centre-sud), Garoua (au nord, sur la Bénoué), Mora (à l'extrême nord, où les Allemands résistent jusqu'en 1916), le fleuve Sanaga (au centre).
\item \textbf{Flèches des offensives franco-britanniques :} depuis le Nigeria vers l'ouest du Cameroun ; depuis le Tchad vers le nord (Garoua, Mora) ; depuis le Congo français vers le sud (Yaoundé) ; débarquement à Douala (1914).
\item \textbf{Retraite allemande :} flèche de repli vers Yaoundé (1915), puis vers la Guinée espagnole (1916).
\item \textbf{Titre complet :} « La campagne du Cameroun, 1914-1916 » ; légende numérotée ; flèche du nord.
\end{enumerate}

\begin{popfigure}
\begin{tikzpicture}[scale=1.1, every node/.style={font=\small}]
% Contour très simplifié du Cameroun (polygone)
\draw[thick, popdark] (0,0) -- (4,0) -- (4.5,1) -- (4,3) -- (1,3.5) -- (0,2) -- cycle;
% Fleuve Sanaga
\draw[popblue, thick] (1.5,3.3) .. controls (2,2.5) and (2.5,1.5) .. (3,0.5) node[right, pos=0.5, popblue] {Sanaga};
% Fleuve Bénoué (Nord)
\draw[popblue, thick] (3.5,3) -- (4.5,2.5) node[right, popblue] {Bénoué};

% Villes
\node[circle, fill=popdark, inner sep=1.5pt, label=below:Douala] (Douala) at (0.5,0.5) {};
\node[circle, fill=popdark, inner sep=1.5pt, label=right:Yaoundé] (Yaounde) at (1.8,2.2) {};
\node[circle, fill=popdark, inner sep=1.5pt, label=above:Garoua] (Garoua) at (3.8,2.5) {};
\node[circle, fill=popdark, inner sep=1.5pt, label=above:Mora] (Mora) at (4.2,3.2) {};

% Flèches offensives (Franco-Britanniques)
% Depuis Nigeria (Ouest) -> Douala / Yaoundé
\draw[->, thick, popgreen] (-0.5, 0.5) -- (Douala);
\node[popgreen, left] at (-0.5, 0.5) {Nigeria};
\draw[->, thick, popgreen] (-0.2, 1.5) .. controls (0.5, 1.8) .. (Yaounde);

% Depuis Tchad (Nord-Est) -> Garoua / Mora
\draw[->, thick, poporange] (5, 3) -- (Mora);
\node[poporange, right] at (5, 3) {Tchad};
\draw[->, thick, poporange] (4.8, 2.2) -- (Garoua);

% Depuis Congo (Sud) -> Yaoundé
\draw[->, thick, poppurple] (1.5, -0.5) -- (Yaounde);
\node[poppurple, below] at (1.5, -0.5) {Congo français};

% Débarquement Douala 1914
\draw[->, thick, popblue] (0.5, -0.5) -- (Douala);
\node[popblue, below] at (0.5, -0.5) {Débarquement 1914};

% Retraite allemande
\draw[->, thick, dashed, poppink] (Mora) .. controls (3.5, 2.8) .. (Garoua);
\draw[->, thick, dashed, poppink] (Garoua) .. controls (3, 2) .. (Yaounde);
\node[poppink, align=center] at (3.5, 1.5) {Retraite\\allemande\\1915-1916};
\draw[->, thick, dashed, poppink] (Yaounde) .. controls (0.5, 1.5) .. (-0.5, 1);
\node[poppink, left, align=center] at (-0.5, 1){Guinée\\(espagnole)};

% Flèche Nord
\node[above] at (4.5, 3.8) {\textbf{N} $\uparrow$};

% Légende
\begin{scope}[shift={(5.5, 2.5)}]
\node[anchor=west] at (0, 1.5) {\textbf{Légende}};
\draw[->, thick, popgreen] (0, 1) -- (0.5, 1) node[right] {Offensives franco-britanniques};
\draw[->, thick, dashed, poppink] (0, 0.5) -- (0.5, 0.5) node[right] {Retraite allemande};
\draw[->, thick, popblue] (0, 0) -- (0.5, 0) node[right] {Débarquement / Flux};
\fill[popdark] (0.1, -0.5) circle (1.5pt) node[right] {Villes principales};
\draw[popblue, thick] (0, -1) -- (0.5, -1) node[right] {Fleuves};
\end{scope}
\end{tikzpicture}
\legende{Croquis schématique : La campagne du Cameroun (1914-1916).}
\end{popfigure}

\textbf{Barème indicatif (sur 10) :} contour et lieux exacts (3 pts) ; flèches des offensives correctement orientées (3 pts) ; retraite allemande (2 pts) ; titre, légende, nord (2 pts).

\textbf{Point de vigilance :} un croquis sans titre ou sans légende numérotée est pénalisé, même si le fond est juste.
\end{corrige}

\begin{corrige}[Exercice 7 — Diagramme en barres du chômage]
\begin{enumerate}
\item Le diagramme comporte en abscisse les années (1929 à 1933), en ordonnée le nombre de chômeurs en millions (gradué de 0 à 14), deux barres juxtaposées par année (une couleur par pays), un titre (« Le chômage aux États-Unis et en Allemagne, 1929-1933 ») et une légende.
\item Maximum : 1933 pour les États-Unis (12,8 millions) ; 1932 pour l'Allemagne (5,6 millions).
\item Entre 1929 et 1932, le chômage augmente fortement et régulièrement dans les deux pays : il est multiplié par environ 8 aux États-Unis (de 1,5 à 12,1 millions) et par plus de 4 en Allemagne (de 1,3 à 5,6 millions). La hausse est donc plus violente aux États-Unis, berceau de la crise.
\end{enumerate}

\begin{popfigure}
\begin{tikzpicture}
\begin{axis}[
    ybar,
    enlargelimits=0.15,
    width=\linewidth,
    height=8cm,
    ylabel={Nombre de chômeurs (millions)},
    xlabel={Années},
    symbolic x coords={1929,1930,1931,1932,1933},
    xtick=data,
    legend style={at={(0.5,-0.15)}, anchor=north, legend columns=-1},
    ymin=0, ymax=14,
    bar width=6pt,
    tick label style={font=\small},
    label style={font=\small},
    legend style={font=\small},
]
\addplot[fill=popblue!70, draw=popblue] coordinates {
    (1929,1.5) (1930,4.3) (1931,8.0) (1932,12.1) (1933,12.8)
};
\addplot[fill=poporange!70, draw=poporange] coordinates {
    (1929,1.3) (1930,3.0) (1931,4.5) (1932,5.6) (1933,6.0)
};
\legend{États-Unis, Allemagne}
\end{axis}
\end{tikzpicture}
\legende{Évolution du chômage aux États-Unis et en Allemagne (1929-1933). Sources : données historiques approximatives.}
\end{popfigure}

\textbf{Point de vigilance :} chaque barre doit être proportionnelle à sa valeur ; les axes doivent porter la grandeur et l'unité (millions de chômeurs).
\end{corrige}

\begin{corrige}[Exercice 8 — Du krach de New York au planteur de cacao]
Ordre logique : krach de Wall Street $\rightarrow$ faillites des banques et des entreprises $\rightarrow$ chômage massif et baisse du pouvoir d'achat $\rightarrow$ baisse de la demande européenne en matières premières $\rightarrow$ effondrement des cours du cacao et du café $\rightarrow$ baisse des revenus des planteurs et misère rurale.

\emph{Paragraphe attendu (modèle) :} D'abord, le krach boursier de Wall Street, en octobre 1929, ruine les spéculateurs et provoque la faillite des banques et des entreprises, aux États-Unis puis en Europe. Ensuite, les licenciements se multiplient : le chômage massif fait chuter le pouvoir d'achat des consommateurs. Par conséquent, la demande européenne en matières premières coloniales diminue fortement. Les cours mondiaux du cacao et du café s'effondrent alors. Enfin, les planteurs camerounais, qui vendaient leur récolte à bas prix, voient leurs revenus s'effondrer : la misère gagne les campagnes. La crise de 1929 est donc bien une crise mondiale, qui frappe jusqu'aux colonies africaines.

\textbf{Point de vigilance :} la chaîne causale doit être respectée : c'est la baisse du pouvoir d'achat qui provoque la baisse de la demande, et non l'inverse.
\end{corrige}

\begin{corrige}[Exercice 9 — Dissertation guidée — type Probatoire]
\emph{Corrigé indicatif (éléments attendus) :}

\textbf{1. Introduction.} Accroche : le 11 novembre 1918, l'armistice met fin à plus de quatre ans de combats. Définition : la guerre totale est une guerre qui mobilise toutes les ressources d'un État (hommes, économie, esprits, empires). Problématique : en quoi la guerre de 1914-1918 fut-elle une guerre totale et quelles en furent les conséquences pour le Cameroun ? Annonce du plan en deux parties.

\textbf{2. Première partie (une guerre totale).} Mobilisation des hommes : des millions de soldats appelés (environ 8 millions de Français), civils à l'arrière. Mobilisation de l'économie : usines tournées vers l'armement, femmes au travail (munitionnettes), emprunts d'État. Mobilisation des esprits et des colonies : propagande, censure, « bourrage de crâne » ; soldats et porteurs venus des empires (tirailleurs sénégalais, porteurs camerounais).

\textbf{3. Deuxième partie (conséquences pour le Cameroun).} Campagne du Cameroun (1914-1916) : combats, réquisitions, déplacements de populations ; fin de la colonisation allemande (retraite vers la Guinée espagnole, 1916) ; partage franco-britannique (France à l'est, Grande-Bretagne à l'ouest) ; statut de territoire sous mandat de la SDN, confirmé en 1922.

\textbf{4. Conclusion.} La guerre de 1914-1918 est bien la première guerre totale ; elle redessine la carte du Cameroun. Ouverture : la paix est fragile et le monde entre bientôt dans la crise économique de 1929.

\textbf{Barème indicatif (sur 20) :} introduction (4 pts) ; première partie, trois arguments datés (6 pts) ; deuxième partie, faits précis sur le Cameroun (6 pts) ; conclusion et ouverture (2 pts) ; présentation et orthographe (2 pts).

\textbf{Points de vigilance :} respecter le plan en deux parties imposé par le sujet ; chaque argument doit être appuyé par un fait daté et localisé ; ne pas traiter la crise de 1929 dans le développement (elle n'appartient qu'à l'ouverture).
\end{corrige}

\begin{corrige}[Exercice 10 — Étude de documents — type Baccalauréat technique]
\begin{enumerate}
\item Le document 1 est une photographie prise en 1916 sur le front de la Somme (nord de la France), pendant la Première Guerre mondiale. L'auteur est probablement un photographe de guerre, soldat ou reporter, témoin direct des combats. C'est une source primaire, contemporaine des événements.
\item Deux caractéristiques de la guerre de tranchées : des conditions de vie épouvantables (boue, froid, promiscuité dans la tranchée) ; un front immobile et des combats d'une violence extrême (terrain dévasté par les obus, barbelés, sacs de terre protecteurs).
\item Total : $2{,}0 + 1{,}4 + 0{,}9 + 1{,}8 = 6{,}1$ millions de morts militaires. Classement du plus touché au moins touché : Allemagne (2,0), Russie (1,8), France (1,4), Royaume-Uni (0,9).

\begin{center}
\begin{tabularx}{\linewidth}{|l|X|}\hline
\textbf{Pays} & \textbf{Pertes militaires (millions)} \\ \hline
Allemagne & 2,0 \\ \hline
Russie & 1,8 \\ \hline
France & 1,4 \\ \hline
Royaume-Uni & 0,9 \\ \hline
\textbf{Total} & \textbf{6,1} \\ \hline
\end{tabularx}
\end{center}

\item \emph{Commentaire attendu (plan) :} introduction présentant les deux documents (nature, date, sujet) ; développement montrant d'abord la guerre de masse (des millions d'hommes mobilisés, enfermés dans les tranchées comme le montre le document 1, armes industrielles : artillerie, mitrailleuses, gaz), ensuite le bilan humain catastrophique (6,1 millions de morts pour les seuls quatre pays du document 2, soit près de 10 millions de morts militaires au total dans le monde, sans compter les blessés, les mutilés et les victimes civiles) ; conclusion : la guerre de 1914-1918, première guerre totale et industrielle, laisse une Europe saignée et traumatisée.
\end{enumerate}
\textbf{Barème indicatif (sur 20) :} présentation du document 1 (3 pts) ; analyse du document 1 (3 pts) ; calcul et classement exacts (3 pts) ; commentaire construit croisant documents et connaissances (9 pts) ; présentation (2 pts).

\textbf{Points de vigilance :} présenter un document, c'est toujours donner sa nature, sa date, son lieu et son auteur probable ; dans le commentaire, citer explicitement les documents (« le document 1 montre que... ») et ne pas se contenter de réciter le cours.
\end{corrige}

\begin{corrige}[Exercice 11 — Question de synthèse — la crise de 1929, une crise mondiale]
\begin{enumerate}
\item Dans les années 1920, les États-Unis, enrichis par la guerre et premiers créanciers du monde, connaissent une forte croissance : production de masse (automobile, électroménager), travail à la chaîne, essor du crédit. Mais cette prospérité est fragile : les usines produisent plus que les salaires ne permettent d'acheter (surproduction), les paysans s'endettent, et la Bourse monte artificiellement grâce à des achats d'actions à crédit (spéculation).
\item Le 24 octobre 1929, la méfiance gagne les investisseurs : tous veulent vendre leurs actions en même temps et les cours s'effondrent à Wall Street. Les banques, qui avaient prêté l'argent des spéculateurs, font faillite ; les entreprises ferment et licencient. La crise gagne l'Europe car les Américains rapatrient leurs capitaux et cessent leurs prêts (faillites bancaires en Allemagne et en Autriche dès 1931). Elle atteint ensuite les colonies africaines : la baisse de la demande européenne fait s'effondrer les cours des matières premières comme le cacao et le café.
\item Les chiffres montrent une crise d'ampleur mondiale : la production industrielle mondiale recule d'environ 38 \% et le commerce mondial d'environ 65 \% entre 1929 et 1932 ; environ 30 millions de chômeurs sont comptés dans le monde en 1932. Le Cameroun est directement touché : le cours du cacao, sa principale richesse d'exportation, perd plus de la moitié de sa valeur (de 22 à moins de 10 centimes le kilogramme, soit une baisse d'environ 55 \%), ce qui ruine les planteurs et réduit les revenus de la colonie.
\item À partir de 1933, le président Roosevelt lance le New Deal (« nouvelle donne ») : l'État intervient dans l'économie en lançant de grands travaux publics pour créer des emplois, en aidant les chômeurs et les paysans, en contrôlant les banques et en encadrant la production pour éviter la surproduction. Cette politique redonne confiance et relance progressivement l'économie américaine.
\end{enumerate}
\textbf{Barème indicatif (sur 20) :} question 1 (4 pts) ; question 2 (6 pts) ; question 3, avec exploitation des chiffres (6 pts) ; question 4 (4 pts).

\textbf{Points de vigilance :} la question 3 exige d'utiliser les données chiffrées fournies (38 \%, 65 \%, 30 millions, cours du cacao) et non des généralités ; pour le cacao, justifier la chute par un calcul de baisse (de 22 à moins de 10 centimes, soit plus de la moitié de la valeur perdue).
\end{corrige}

\newpage

\coursec{Fiche bilan}

\begin{popfigure}
\begin{tikzpicture}[
  mindmap,
  concept color=popblue,
  level 1 concept/.append style={level distance=4.5cm, sibling angle=360/7},
  level 2 concept/.append style={level distance=3cm},
  every node/.style={font=\small, text=white},
  grow cyclic,
  text width=2.8cm,
  align=center
]
\node[concept] {Guerre 14--18 \& Crise 1929}
  child[concept color=poporange] {node[concept, align=center]{Causes\\Alliances, nationalismes,\\crise juil. 1914}
    child[concept color=poporange!70] {node[concept, align=center]{Sarajevo\\28 juin 1914}}}
  child[concept color=popgreen] {node[concept, align=center]{Guerre totale\\Tranchées, fronts,\\entrée US 1917}
    child[concept color=popgreen!70] {node[concept, align=center]{Verdun, Somme\\1916}}}
  child[concept color=poppurple] {node[concept, align=center]{Conséquences\\Traité Versailles,\\SDN, nouvelles cartes}
    child[concept color=poppurple!70] {node[concept, align=center]{14 pts Wilson\\Réparations}}}
  child[concept color=poppink] {node[concept, align=center]{Années 1920\\Prospérité US,\\crédit, spéculation}
    child[concept color=poppink!70] {node[concept, align=center]{Fordisme\\Bourse en hausse}}}
  child[concept color=popgold] {node[concept, align=center]{Krach 1929\\Jeudi noir 24 oct.,\\mardi noir 29 oct.}
    child[concept color=popgold!70] {node[concept, align=center]{Effondrement\\cours bourse}}}
  child[concept color=popdark] {node[concept, align=center]{Grande Dépression\\Chômage massif,\\protectionnisme}
    child[concept color=popdark!70] {node[concept, align=center]{Crise bancaire\\commerce mondial}}}
  child[concept color=popblue!80!black] {node[concept, align=center]{Issues politiques\\New Deal, totalitarismes,\\voie vers 1939}
    child[concept color=popblue!60!black] {node[concept, align=center]{Hitler 1933\\Roosevelt 1933}}};
\end{tikzpicture}
\legende{Carte mentale de synthèse : causes, déroulement et suites de la Grande Guerre ; prospérité, krach et dépression des années 1920--1930.}
\end{popfigure}

\begin{bilan}
\courssub{Définitions clés}
\begin{itemize}
  \item \cle{Guerre totale} : mobilisation de toutes les ressources (humaines, économiques, morales) d'un pays ; front et arrière confondus.
  \item \cle{Système d'alliances} : Triple-Alliance (Allemagne, Autriche-Hongrie, Italie) vs Triple-Entente (France, Royaume-Uni, Russie).
  \item \cle{Traité de Versailles} (28 juin 1919) : paix imposée à l'Allemagne (pertes territoriales, désarmement, réparations, clause de responsabilité).
  \item \cle{SDN} (Société des Nations) : organisation internationale née en 1920 pour garantir la paix collective ; échec face aux révisionnismes.
  \item \cle{Spéculation boursière} : achat d'actions à crédit (marge) en pariant sur la hausse continue des cours.
  \item \cle{Krach de 1929} : effondrement brutal des cours à Wall Street (octobre 1929), point de départ de la Grande Dépression.
  \item \cle{Grande Dépression} : crise économique mondiale (1929--années 1930) : surproduction, déflation, chômage de masse, effondrement du commerce international.
  \item \cle{New Deal} : politique de relance keynésienne de Roosevelt (1933) : intervention de l'État, régulation bancaire, grands travaux.
\end{itemize}

\courssub{Repères chronologiques indispensables}
\begin{center}
\begin{tabularx}{\linewidth}{|c|X|}\hline
\textbf{Date} & \textbf{Événement majeur} \\ \hline
1914 (28 juin -- 4 août) & Attentat de Sarajevo ; engrenage des alliances ; entrée en guerre \\ \hline
1917 (avril) & Entrée en guerre des États-Unis \\ \hline
1918 (11 nov.) & Armistice de Rethondes ; fin des combats \\ \hline
1919 (28 juin) & Traité de Versailles ; création de la SDN \\ \hline
1924--1929 & « Années de prospérité » aux USA ; plan Dawes, spéculation \\ \hline
1929 (24 \& 29 oct.) & Krach de Wall Street (Jeudi noir, Mardi noir) \\ \hline
1930--1933 & Extension mondiale : faillites bancaires, protectionnisme (Hawley-Smoot) \\ \hline
1933 (mars) & Roosevelt président ; lancement du New Deal \\ \hline
1933 (janv.) & Hitler chancelier en Allemagne \\ \hline
\end{tabularx}
\end{center}

\courssub{Méthodes express (Bac)}
\begin{enumerate}
  \item \textbf{Commentaire de document :} 1) Identifier (nature, auteur, date, source). 2) Contextualiser (rappeler l'événement lié). 3) Analyser (idées forces, vocabulaire, chiffres). 4) Conclure (portée, limites, lien avec le cours).
  \item \textbf{Dissertation / Rédaction d'invention :} 1) Analyser le sujet (mots-clés, bornes spatio-temporelles). 2) Problématiser. 3) Plan dialectique (Thèse / Antithèse / Synthèse) ou thématique (Causes / Déroulement / Conséquences). 4) Rédiger en paragraphes argumentés (1 idée = 1 paragraphe + exemple précis). 5) Conclusion ouverte.
  \item \textbf{Croquis / Carte :} Titre précis, légende organisée (hiérarchie), figurés adaptés (points, surfaces, flèches), orientation, échelle indicative, propreté.
\end{enumerate}
\end{bilan}

\begin{aretenir}[Checklist avant le Bac]
\begin{itemize}
  \item $\square$ Je sais expliquer l'engrenage des alliances et le rôle de l'attentat de Sarajevo (juillet 1914).
  \item $\square$ Je caractérise la guerre totale (tranchéisation, mobilisation industrielle, propagande, entrée US 1917).
  \item $\square$ Je cite les clauses principales du Traité de Versailles et les faiblesses de la SDN.
  \item $\square$ Je décris la prospérité américaine des années 1920 (fordisme, crédit, spéculation boursière).
  \item $\square$ Je distingue les causes structurelles (surproduction, inégalités, endettement) du déclencheur (krach d'octobre 1929).
  \item $\square$ Je trace la propagation de la crise : crise boursière $\to$ crise bancaire $\to$ crise économique $\to$ crise sociale/mondiale.
  \item $\square$ Je compare les réponses : New Deal (État-providence, relance) vs solutions autoritaires (Allemagne, Italie, Japon).
  \item $\square$ Je sais situer sur une frise : 1914, 1917, 1918, 1919, 1929, 1933.
  \item $\square$ Je suis capable de commenter un document (discours de Wilson, affiche de propagande, graphique du chômage, extrait du New Deal).
  \item $\square$ Je sais réaliser un croquis : « Les fronts de la Grande Guerre » ou « La propagation de la crise de 1929 ».
  \item $\square$ J'utilise le vocabulaire précis : armistice, réparations, clause de culpabilité, protectionnisme, déflation, keynésianisme.
  \item $\square$ Je relie les conséquences de 1918-1919 aux causes de la Seconde Guerre mondiale (révisionnisme, humiliation allemande, faillite de la SDN).
\end{itemize}
\end{aretenir}

\findecours

\end{document}
